\documentclass[11pt,lang=cn,math=cm,thmcnt=section,usesamecnt]{elegantbook}
\usepackage{amssymb}
\usepackage{mathtools}
\IfFontExistsTF{KaiTi}{\setCJKfamilyfont{zhkai}[AutoFakeBold=true]{KaiTi}}{}
\IfFontExistsTF{FangSong}{\setCJKfamilyfont{zhfs}[AutoFakeBold=true]{FangSong}}{}
\usepackage[noautomatic]{imakeidx}
\makeindex[title=术语索引,intoc,columns=2]
\indexsetup{othercode=\small\setstretch{1.03}\renewcommand{\indexspace}{\par\vspace{3pt}}}
\usepackage{tikz-cd}
\usepackage{needspace}
\usepackage{float}
\usepackage{longtable}
\usepackage{booktabs}
\usepackage{array}
\usepackage{xurl}
\hypersetup{unicode=true,colorlinks=true,linkcolor=structurecolor,urlcolor=structurecolor,citecolor=structurecolor,bookmarksnumbered=true,pdftitle={朗兰兹纲领入门课程讲义（中文译注本）},pdfauthor={Johannes Anschütz; Ben Heuer}}
\setcounter{tocdepth}{1}
\setcounter{secnumdepth}{0}
\setlength{\emergencystretch}{3em}
\setlength{\intextsep}{8pt plus 2pt minus 2pt}
\tcbset{common/.append style={before skip=4pt,after skip=4pt}}
\renewcommand{\xchaptertitle}{第\thechapter{}讲}
\title{朗兰兹纲领入门课程讲义}
\subtitle{含 Ben Heuer 的三讲}
\author{Johannes Anschütz}
\date{2020 年 7 月 22 日修订稿}
\makeatletter
\let\nativeNormalsize\normalsize
\renewcommand{\normalsize}{\nativeNormalsize\fontsize{11pt}{13.8pt}\selectfont\setlength{\abovedisplayskip}{8pt plus 2pt minus 2pt}\setlength{\belowdisplayskip}{8pt plus 2pt minus 2pt}\setlength{\abovedisplayshortskip}{4pt plus 2pt minus 1pt}\setlength{\belowdisplayshortskip}{6pt plus 2pt minus 2pt}}
\makeatother
\DeclareMathOperator{\Gal}{Gal}
\DeclareMathOperator{\GL}{GL}
\DeclareMathOperator{\SL}{SL}
\DeclareMathOperator{\SO}{SO}
\DeclareMathOperator{\Hom}{Hom}
\DeclareMathOperator{\End}{End}
\DeclareMathOperator{\Spec}{Spec}
\DeclareMathOperator{\Tr}{Tr}
\DeclareMathOperator{\Res}{Res}
\DeclareMathOperator{\Ind}{Ind}
\DeclareMathOperator{\Pic}{Pic}
\DeclareMathOperator{\Sym}{Sym}
\DeclareMathOperator{\Lie}{Lie}
\DeclareMathOperator{\diag}{diag}
\DeclareMathOperator{\Frob}{Frob}
\DeclareMathOperator{\ord}{ord}
\DeclareMathOperator{\vol}{vol}
\DeclareMathOperator{\Div}{Div}
\DeclareMathOperator{\Aut}{Aut}
\DeclareMathOperator{\Sat}{Sat}
\newcommand{\Q}{\mathbb{Q}}
\newcommand{\N}{\mathbb{N}}
\newcommand{\Z}{\mathbb{Z}}
\newcommand{\R}{\mathbb{R}}
\newcommand{\C}{\mathbb{C}}
\newcommand{\F}{\mathbb{F}}
\newcommand{\A}{\mathbb{A}}
\newcommand{\Af}{\mathbb{A}_f}
\newcommand{\OO}{\mathcal{O}}
\newcommand{\Gm}{\mathbb{G}_m}
\newcommand{\HH}{\mathbb{H}}
\newcommand{\Pone}{\mathbb{P}^1}
\newcommand{\g}{\mathfrak{g}}
\newcommand{\p}{\mathfrak{p}}
\newcommand{\m}{\mathfrak{m}}
\newcommand{\n}{\mathfrak{n}}



\newcommand{\Fq}{\overline{\mathbb{F}}_p}
\newcommand{\Ql}{\overline{\mathbb{Q}}_\ell}
\newcommand{\cH}{\mathcal{H}}
\newcommand{\cZ}{\mathcal{Z}}
\newcommand{\cA}{\mathcal{A}}
\newcommand{\cE}{\mathcal{E}}
\newcommand{\cF}{\mathcal{F}}
\newcommand{\cL}{\mathcal{L}}
\newcommand{\cK}{\mathcal{K}}
\newcommand{\cP}{\mathcal{P}}
\newcommand{\cM}{\mathcal{M}}
\newcommand{\cS}{\mathcal{S}}
\newcommand{\cD}{\mathcal{D}}
\newcommand{\znote}[1]{\footnote{\textbf{【译注】}#1}}
\newcommand{\znoteplain}[1]{\footnote{\textbf{【译者按】}#1}}
\newcommand{\sourcepage}[1]{\phantomsection\label{src:#1}}
\newcommand{\idx}[1]{\index{#1|hyperpage}}
\setlength{\skip\footins}{10pt plus 4pt minus 2pt}
\newenvironment{references}{\begin{description}[leftmargin=3.2em,labelwidth=2.6em,itemsep=2pt]}{\end{description}}
\newcommand{\refitem}[2]{\item[\hypertarget{#1}{[#2]}]}
\begin{document}
\begin{titlepage}
\centering
\vspace*{2.2cm}
{\color{structurecolor}\rule{\textwidth}{1.2pt}}\par
\vspace{1.4cm}
{\Huge\bfseries 朗兰兹纲领入门课程讲义\par}
\vspace{0.9cm}
{\Large 含 Ben Heuer 的三讲\par}
\vspace{0.6cm}
{\large\normalfont Lecture notes for an introductory course\\ on the Langlands program\par}
\vspace{1.2cm}
{\color{structurecolor}\rule{\textwidth}{0.6pt}}\par
\vfill
{\large Johannes Anschütz\par}
\vspace{0.4cm}
{\large 2020 年 7 月 22 日修订稿\par}
\vspace{1.6cm}
{\large 中文译注本\quad 数译 · SY02-38\par}
\vspace{0.8cm}
\end{titlepage}
\hypersetup{pageanchor=true}
\frontmatter
\tableofcontents
\chapter*{译者说明}
\addcontentsline{toc}{chapter}{译者说明}

本书是 Johannes Anschütz《Lecture notes for an introductory course on the Langlands program》（2020 年 7 月 22 日修订稿，103 页）的中文译注本，其中第 7、8、10 讲由 Ben Heuer 撰写，标题中已分别标明。

\section*{底本与译名}

底本为作者在波恩大学开设的朗兰兹纲领入门课程的讲义，共 12 讲：前六讲从模形式出发，经 adèle 化、Hecke 代数与强重数一，把新形式组织成自守表示，并陈述 $\mathrm{GL}_2$ 上的 Langlands 互反；第 7、8 讲（Heuer）讲模曲线作为椭圆曲线模空间、Eichler--Shimura 关系；第 9 讲用上同调实现附着于新形式的 Galois 表示；第 10 讲（Heuer）用模 $p$ 同余与 Deligne--Serre 提升处理权 $1$ 形式；第 11、12 讲把整套陈述推广到一般约化群（$L$-群、Satake 参数、$L$-参数、Buzzard--Gee 猜想），并以 Shimura 簇收尾。

译文尽量保留讲义的条目层级与标签（「上次」「要点」「事实」「注」「略述」），文献引用（如 \texttt{[GH19, Lemma 3.7.1]}）与定义、定理、公式的编号一律原样保留，不重编。

\section*{体例}

\begin{itemize}
\item 页边不显示原文页码；每页正文前在源码中以注释形式保留页标记，供校对核对，不进入成品。
\item 译者添加的解释以脚注形式出现，标有「【译注】」者均有可追溯来源（见脚注内的\emph{来源}），标有「【译者按】」者为本译本的补充说明、无外部依据。
\item 原文中的笔误（例如 \S 5 中 $\Gamma_0(m)$ 的 $\Gamma_0$ 条件、\S 4 中 $Y$ 的矩阵分解、\S 11 中 Casimir 元的一段中间计算等）均在译注中说明原文写法与校正依据，正文按校正后的含义翻译。
\item 术语与同主题前作《Langlands 纲领导论》（SY02-09）对齐；少数与《朗兰兹纲领导论（全书中文译注）》（SY02-24）不同的译法，在首次出现处以括号附英文原词。
\end{itemize}

\section*{致谢与声明}

感谢原作者 Johannes Anschütz 与 Ben Heuer 写出并公开这份讲义。本译注本为非商业学术译本，底本版权归原作者所有；若原作者不愿本译本流传，请告知，本库将立即撤下。译文中一切错误由本译本负责。

\mainmatter
\chapter{关于本讲稿的说明}

% ===== PDF p.002 =====
这是朗兰兹纲领入门课程的讲义（略作修订）。\znoteplain{原文脚注：由于本课程原计划与同名、但最终被取消的三学期项目（trimester program）平行开设，最初的题目是「朗兰兹纲领的算术」。然而这个题目并未准确描述课上所讲的内容，因此我在修订稿中改换了题目。}

课程最终的目标是：对 $\Q$ 上的一般约化群\idx{约化群}，陈述（大多不给出任何证明的迹象）朗兰兹纲领中出现的各种陈述、对象、猜想……，并解释它们如何特殊化为 $\mathrm{GL}_2$ 的情形。另一个目标是给出附着于模形式\idx{模形式}的 Galois 表示的构造（重点在权 $1$ 与权 $2$ 的情形）。

原本还计划纳入更多关于局部 Langlands 纲领\idx{Langlands 纲领}的材料。但此事被推迟了：2020/2021 年冬季学期，Peter Scholze 宣布将在波恩大学开设这方面的课程。

为使讲义更易读，我把一些原本只在课堂上口头说过的话也写了进来。

\textbf{免责声明}：我远不是对朗兰兹纲领有严肃知识或理解的人，因此讲义中我可能把问题过度简化或过度复杂化，表述不准确甚至出错，并漏掉细微之处。因此请自担风险地使用本讲义，确切的、正确的陈述请查阅所引文献。欢迎对本讲义提出任何意见或提示！

我要感谢课程参与者在课堂上提出的有趣问题。此外，我要衷心感谢 Ben Heuer 采纳我的建议讲了三次课（哪三讲在相应的标题中标明）。

\chapter{朗兰兹纲领概观}

% ===== PDF p.003 =====
本讲我们想对朗兰兹纲领作一个粗略的介绍，并给出它算术意义的第一个例子。

回忆：域 $k$ 上连通、光滑、仿射的代数群 $G$ 称为\textbf{约化的}（reductive），若其基变换 $G_{\bar k}$ 到 $k$ 的代数闭包上不包含同构于加法群 $\mathbb{G}_a$ 的闭正规子群；等价地，$G_{\bar k}$ 的幂幺根（unipotent radical）平凡。

对每个素数 $p$，令
\[
\Z_p := \varprojlim_n \Z/p^n \Z,\qquad \Q_p := \Z_p\otimes_\Z \Q
\]
分别为 $p$ 进整数环与 $p$ 进数域。令
\[
\Af := \Bigl\{(x_p)_p\in\prod_p \Q_p \Bigm| x_p\in\Z_p \text{ 除有限多个素数 } p \text{ 外}\Bigr\}
\]
为 $\Q$ 的\textbf{有限阿代尔环\idx{阿代尔环}}（ring of finite adèles），并令
\[
\A := \Af\times\R
\]
为 $\Q$ 的\textbf{完全阿代尔\idx{阿代尔}环}。

回忆 $\A$ 是局部紧拓扑环。即：$\Af$ 上存在唯一的环拓扑使子环 $\prod_p\Z_p\subseteq\Af$ 为开；$\A$ 上则取它与 $\R$ 上通常拓扑的积拓扑。一个重要性质是：子环 $\Q\subseteq\A$（对角嵌入）上的子空间拓扑是离散的。

朗兰兹纲领的起点是 $\Q$ 上的一个约化群 $G$，最重要的例子是 $G=\GL_n$（某个 $n\ge1$）。

上述事实蕴涵：$G(\A)\cong G(\Af)\times G(\R)$ 自然地是局部紧拓扑群，且 $G(\Q)\subseteq G(\A)$ 是离散子群。具体地，取嵌入 $G\subseteq\Spec(\Q[X_1,\dots,X_m])$（某个 $m\ge1$），并取诱导嵌入 $G(\A)\subseteq\A^m$ 的子空间拓扑；这个拓扑与嵌入的选取无关。细节见 \texttt{[GH19, Theorem 2.2.1.]}。

我们将使用如下事实：

\begin{itemize}
\item 任意局部紧拓扑群 $H$ 上都存在右不变 Haar 测度，见 \texttt{[GH19, Section 3.2.]}，且它在 $\R_{>0}$ 中的正标量倍意义下唯一。
\item $G(\A)$ 上的右不变 Haar 测度同时也是左不变测度，并下降为 $G(\Q)\backslash G(\A)$ 上的 $G(\A)$-不变测度，见 \texttt{[GH19, Lemma 3.5.4.]}、\texttt{[GH19, Lemma 3.5.3]}。
\end{itemize}

现在引入朗兰兹纲领中关于 $G$ 的最终感兴趣的对象。

\noindent\textbf{定义 2.1}\phantomsection\label{def:2.1}\quad 令
\[
L^2(G(\Q)\backslash G(\A))
\]
为可测函数 $f:G(\Q)\backslash G(\A)\to\C$ 且
\[
\int_{G(\Q)\backslash G(\A)}|f|^2<\infty
\]
构成的空间，

% ===== PDF p.004 =====
其中积分关于 $G(\Q)\backslash G(\A)$ 上的 $G(\A)$-不变测度。

与通常一样，若两个可测函数在一个零测集之外相等，则在 $L^2$ 空间中视为同一个元素。

空间 $L^2(G(\Q)\backslash G(\A))$ 通过内积
\[
(f_1,f_2)\mapsto\int_{G(\Q)\backslash G(\A)}f_1(g)\overline{f_2(g)}\,\mathrm{d}g
\]
自然成为 Hilbert 空间。$G(\A)$ 经右平移作用在 $G(\Q)\backslash G(\A)$ 上，由此诱导出 $G(\A)$ 在 $L^2(G(\Q)\backslash G(\A))$ 上的（左）作用；由于所选取的测度是右不变的，这个作用是酉作用，即保持内积。

现在我们可以（以极为粗略的形式）陈述朗兰兹纲领的一个主要目标：

\begin{quote}
把 Hilbert 空间 $L^2(G(\Q)\backslash G(\A))$ 作为 $G(\A)$-表示，按算术数据加以分解。
\end{quote}

我们会略微说明「分解」是什么意思。记
\[
\widehat{G(\A)}
\]
为（Hilbert 空间上的）不可约酉 $G(\A)$-表示的同构类集。由于分析上的问题，期望存在形如
\[
L^2(G(\Q)\backslash G(\A))\ \text{“}\cong\text{”}\ \bigoplus_{[\pi]\in\widehat{G(\A)}}\pi^{\oplus m_\pi}
\]
的直和分解（其中 $m_\pi\in\N\cup\{\infty\}$ 是某个不可约酉 $G(\A)$-表示 $\pi$ 的重数）是过于天真的。

让我们考察 $G=\Gm=\GL_1$ 的例子。此时
\[
G(\Q)\backslash G(\A)\cong\Q^\times\backslash\A^\times
\]
是局部紧 abelian 群，因此可以用（抽象的）Fourier 理论来分析 $\GL_1(\A)=\A^\times$-表示
\[
L^2(\Q^\times\backslash\A^\times).
\]
下面回忆抽象 Fourier 理论的一些结果，依 \texttt{[DE14, Chapter 3]}。

\begin{itemize}
\item 设 $A$ 为任意局部紧 abelian 群（如 $\R$、$S^1\cong\R/\Z$、$\A^\times$、$\Z/n$、$\Q_p$ 等）。
\item 设 $\widehat A$ 为 $A$ 的不可约酉表示的同构类集。
\item 则 $\widehat A$ 与酉连续特征 $A\to S^1$ 的集合 $\mathrm{Hom}_{\mathrm{cts}}(A,S^1)$ 一一对应；也就是说，每个 $\pi\in\widehat A$ 都是一维的，并且（通过逐点乘法与紧开拓扑）本身又成为局部紧 abelian 群。
\item 此外 $A\cong\widehat{\widehat A}$ 作为拓扑群（同构由 $a\mapsto(\chi\mapsto\chi(a))$ 给出），并且 $A$ 离散当且仅当 $\widehat A$ 紧。\znoteplain{原文脚注：“$\Rightarrow$” 取 $\coprod_I\Z\twoheadrightarrow A$，则 $\widehat A$ 嵌入 $\prod_I S^1$。“$\Leftarrow$” 用 $A\cong\widehat{\widehat A}$ 与紧开拓扑的定义。}
\end{itemize}

% ===== PDF p.005 =====
Plancherel 定理是抽象 Fourier 理论的主要结果。对我们而言，它的意义在于给出了空间 $L^2(A)$ 作为 $A$-表示的另一种描述。

\noindent\textbf{定理 2.2}（Plancherel 定理，见 \texttt{[DE14, Chapter 3.4]}）\phantomsection\label{thm:2.2}\quad Fourier 变换
\[
\mathcal F:L^1(A)\cap L^2(A)\longrightarrow L^2(\widehat A),\qquad
f\mapsto\Bigl(\chi\mapsto\int_A f(x)\chi(x)\,\mathrm{d}\mu(x)\Bigr)
\]
延拓为 Hilbert 空间的酉同构
\[
\mathcal F:L^2(A)\xrightarrow{\ \sim\ }L^2(\widehat A).
\]
当我们让 $A$ 通过
\[
a\cdot g(\chi):=\chi(a)g(\chi),\qquad a\in A,\ g\in L^2(\widehat A),\ \chi\in\widehat A
\]
作用在 $L^2(\widehat A)$ 上时，Fourier 变换 $\mathcal F$ 是 $A$-等变同构。因此作为 $A$-表示有
\[
L^2(A)\cong L^2(\widehat A).
\]

让我们讨论两个例子：$A=S^1$ 与 $A=\R$。

\begin{itemize}
\item 若 $A=S^1$，则
\[
\Z\cong\widehat{A},\qquad n\mapsto\chi_n,\quad \chi_n(z)=z^n,\ z\in S^1.
\]
于是
\[
L^2(S^1)\cong\widehat{\bigoplus_{n\in\Z}}\C\chi_n
\]
是 $S^1$-等变子空间 $\C\chi_n\subseteq L^2(S^1)$ 的 Hilbert 空间直和。具体地说：每个 $L^2$ 函数 $f:S^1\to\C$ 都可以唯一地写成
\[
f=\sum_{n\in\Z}a_n\chi_n,\qquad a_n\in\C,\quad \sum_{n\in\Z}|a_n|^2<\infty.
\]
这样我们就找到了酉 $S^1$-表示 $L^2(S^1)$ 分解为不可约子空间的一个漂亮分解，分析上的问题并不大。不过，我们不得不考虑 Hilbert 空间直和 $\widehat{\bigoplus}_{n\in\Z}\C\chi_n$，而不是（代数的）表示直和。
\item 若 $A=\R$，则 $\widehat{\R}\cong\R$，同构由
\[
\R\to\widehat{\R},\qquad x\mapsto\chi_x,
\]
给出，其中 $\chi_x:\R\to S^1$，$y\mapsto e^{2\pi ixy}$。同构
\[
\mathcal F:L^2(\R)\xrightarrow{\ \sim\ }L^2(\R)
\]
因而蕴涵：每个 $f\in L^2(\R)$ 都可以写成
\[
f(x)=\int_\R g(y)\chi_y(x)\,\mathrm{d}y
\]

% ===== PDF p.006 =====
的形式，其中 $g\in L^2(\R)$ 唯一。因此
\[
L^2(\R)\cong\int_\R\chi_y\,\mathrm{d}y
\]
是表示的\textbf{Hilbert 积分}（见 \texttt{[GH19, Section 3.7]}），这个概念不像 Hilbert 空间直和那样直接。注意我们无法做得更好：对任意 $x\in\R$ 都有 $\chi_x\notin L^2(\R)$，而且 $L^2(\R)$ 不含 $\R$ 的任何不可约子表示，见 \texttt{[GH19, Section 3.7]}。\znoteplain{原文脚注：也没有任何不可约商；\texttt{[GH19, Lemma 3.7.1]} 中关于商的陈述，我认为是错的。}
\end{itemize}

让我们回到 $G=\Gm$，即 $A=\A^\times$ 的情形。

\begin{itemize}
\item 把每个 $n\in\Q^\times$ 分解为 $n=\pm p_1\cdots p_k$（$p_i$ 为素数），可推出
\[
\Q^\times\backslash\A^\times\cong\prod_p\Z_p^\times\times\R_{>0}.
\]
\item 我们暂时忽略因子 $\R_{>0}$，考察
\[
L^2\Bigl(\prod_p\Z_p^\times\Bigr)\cong\widehat{\bigoplus_{\chi}}\C\chi,
\]
其中 $\chi:\prod_p\Z_p^\times\to S^1$ 遍历 $\prod_p\Z_p^\times$ 的全部连续特征。注意每个特征 $\chi:\prod_p\Z_p^\times\cong\varprojlim_m(\Z/m)^\times\to S^1$ 都通过某个商 $(\Z/m)^\times$ 分解。
\item 这些特征 $\chi$ 是「算术」数据，即
\[
\Gal(\Q(\mu_\infty)/\Q)\cong\prod_p\Z_p^\times,
\]
其中 $\Q(\mu_\infty)=\bigcup_{n\in\N}\Q(e^{2\pi i/n})$ 是 $\Q$ 的分圆扩张。定义
\[
\alpha:\Gal(\overline\Q/\Q)\twoheadrightarrow\Gal(\Q(\mu_\infty)/\Q)\cong\prod_{p\ \text{素数}}\Z_p^\times
\]
为自然投影。
\item 因此对每个 $\chi:\prod_p\Z_p^\times\to S^1$，复合 $\chi\circ\alpha$ 都是 $\Gal(\overline\Q/\Q)$ 的一个（连续）一维 Galois 表示。
\item 由 Kronecker--Weber 定理，每个一维 Galois 表示
\[
\sigma:\Gal(\overline\Q/\Q)\to\C^\times
\]
都具有 $\sigma=\chi\circ\alpha$ 的形式，其中 $\chi$ 如上。
\end{itemize}

\noindent\textbf{定理 2.3}（Kronecker--Weber）\phantomsection\label{thm:2.3}\quad 域 $\Q(\mu_\infty)$ 是 $\Q$ 的极大 abelian 扩张；也就是说，每个 Galois 群 $\Gal(F/\Q)$ 为 abelian 的有限 Galois 扩张 $F/\Q$ 都包含在 $\Q(\mu_\infty)$ 中。

于是对 $G=\Gm$ 我们得到如下分解：
\[
L^2(\R_{>0}\Q^\times\backslash\A^\times)\cong\widehat{\bigoplus_{\sigma:\Gal(\overline\Q/\Q)\to\C^\times}}\C\chi_\sigma,
\]
% ===== PDF p.007 =====
其中 $\chi_\sigma$ 由 $\chi_\sigma\circ\alpha=\sigma$ 刻画。这个（简单的）例子是朗兰兹纲领所追求目标的原型。注意模掉 $\R_{>0}$ 是方便而且无害的。全空间
\[
L^2(\Q^\times\backslash\A^\times)
\]
是表示 $L^2(\R_{>0}\Q^\times\backslash\A^\times)$ 在 $\R_{>0}$ 的酉特征空间 $\widehat{\R}\cong\R$ 上的 Hilbert 积分。

现在转向一般的 $G$，讨论空间
\[
G(\Q)\backslash G(\A)
\]
的几何。这将比 $G=\GL_1$ 的情形复杂得多。

存在中心子群 $A_G\subseteq G(\R)$，$A_G\cong\R_{>0}^r$，使得
\[
[G]:=A_GG(\Q)\backslash G(\A)
\]
具有有限体积，见 \texttt{[GH19, Theorem 2.6.2.]}。具体地，$A_G$ 可以取作 $G$ 的中心 $Z(G)$ 的极大分裂子环面的连通分支；例如 $G=\GL_n$ 时 $A_{\GL_n}\subseteq\GL_n(\R)$ 是元素在 $\R_{>0}$ 中的标量矩阵构成的群，而 $A_{\SL_n}$ 平凡。

同样，全空间 $L^2(G(\Q)\backslash G(\A))$ 是表示 $L^2([G])$ 在 $\widehat{A_G}$ 上的 Hilbert 积分，因此用 $L^2([G])$ 代替 $L^2(G(\Q)\backslash G(\A))$ 可以省去一些分析上的麻烦。\znoteplain{原文脚注：精确地说，$G(\A)$ 上每个右不变 Haar 测度都下降为 $[G]$ 上的测度，见 \texttt{[GH19, Section 2.6]}。}

虽然我们只关心 $G(\A)$ 上任意的右不变测度\znoteplain{原文脚注：因为所得的 $L^2([G])$ 空间彼此同构。}，但请允许我们提到：$G(\A)$ 上存在一个典范测度，即\textbf{玉河测度}（Tamagawa measure），对这个测度而言，数
\[
\tau(G):=\mathrm{vol}([G])
\]
称为 $G$ 的\textbf{玉河数\idx{玉河数}}（Tamagawa number）。玉河数的知识在算术上很有意义，参见 \texttt{[Col11, Appendix B]}、\texttt{[Lur, Lecture 1]}。例如 $\tau(\SL_n)=1$，这蕴涵
\[
\mathrm{vol}(\SL_n(\R)/\SL_n(\Z))=\zeta(2)\zeta(3)\cdots\zeta(n),
\]
其中 $\zeta(s)$ 是 Riemann $\zeta$ 函数。

拓扑空间 $G(\Q)\backslash G(\A)$ 与 $[G]$ 都是某个实流形的 profinite 覆盖。例如 $G=\GL_n$ 时

% ===== PDF p.008 =====
$\GL_n(\Z)\backslash\bigl(\prod_p\GL_n(\Z_p)\times\GL_n(\R)\bigr)$（$\GL_n(\Z)$ 对角地作用）覆盖实流形 $\GL_n(\Z)\backslash\GL_n(\R)$；类似地 $[G]$ 覆盖 $\GL_n(\Z)\R_{>0}\backslash\GL_n(\R)$。见 \texttt{[GH19, Section 2.6]}、\texttt{[Del73, (0.1.4.1)]}。

更精确地说，对每个紧开子群 $K\subseteq G(\Af)$（也称为一个「level 子群」），商
\[
[G]/K
\]
都是实流形。

类似地，设 $K_\infty\subseteq G(\R)$ 为 $G(\R)$ 的极大紧子群（连通与否不论），例如 $G=\GL_n$ 时可取 $K_\infty\cong\SO_n(\R)$。由于 $K_\infty$ 与 $A_G\subseteq G(\R)$ 平凡相交，空间 $[G]$ 是 $[G]/K_\infty$ 上的主 $K_\infty$-丛（或称 $K_\infty$-挠子）。

下一讲我们将更详细地讨论 $G=\GL_2$ 的情形。现在只提一点：若 $K_\infty\cong\SO_2(\R)$，则由 Möbius 变换给出同构
\[
A_G\backslash\GL_2(\R)/K_\infty\xrightarrow{\ \sim\ }\mathbb{H}^\pm,\qquad
g=\begin{pmatrix}a&b\\c&d\end{pmatrix}\mapsto g(i):=\frac{a\cdot i+b}{c\cdot i+d},
\]
其中
\[
\mathbb{H}^\pm:=\{z\in\C\mid \operatorname{Im}(z)\neq0\}
\]
是上半平面与下半平面之并。特别地，
\[
[\GL_2]/K_\infty\cong\GL_2(\Z)\backslash\Bigl(\GL_2\Bigl(\prod_p\Z_p\Bigr)\times\mathbb{H}^\pm\Bigr).
\]

设 $K\subseteq G(\Af)$ 为紧开子群，则集合
\[
G(\Q)\backslash G(\Af)/K=\bigsqcup_{i=1}^m G(\Q)g_iK
\]
有限（见 \texttt{[GH19, Theorem 2.6.1]}），并且
\[
[G]/K_\infty K\cong\bigsqcup_{i=1}^m\Gamma_i\backslash X,
\]
其中
\[
X:=A_G\backslash G(\R)/K_\infty
\]
是实流形，而
\[
\Gamma_i:=G(\Q)\cap G(\R)g_iKg_i^{-1}\subseteq G(\R)
\]
是离散子群（同余子群的一个例子，见 \texttt{[GH19, Section 2.6]}）。当 $K$ 充分小时，$\Gamma_i$ 在 $X$ 上自由且真不连续地作用，见 \texttt{[GH19, Definition 15.2]}。此时态射 $X\to\Gamma_i\backslash X$ 是以 $\Gamma_i$ 为覆盖群的（拓扑意义下的）覆盖。我们最关心的情形是 $G=\GL_2$，此时一切都变得更具体：\[\Gamma_i\backslash X\cong\Gamma\backslash\mathbb{H}^\pm,\]

其中 $\Gamma\subseteq\GL_2(\Z)$ 是包含
\[
\Gamma(m):=\ker(\GL_2(\Z)\to\GL_2(\Z/m))
\]
（某个 $m\ge0$）的离散子群。当 $m\ge3$ 时，$\Gamma(m)$ 在 $\mathbb{H}^\pm$ 上的作用自由。这个例子将在下一讲提供与模形式的联系。\znote{原文此处作 “For K is sufficiently small”，语法有误，译文按 “When $K$ is sufficiently small” 理解。}

注意对一般的 $G$，通常 $\dim(\Gamma_i\backslash X)>0$。因此 $[G]$ 的几何、以及空间 $L^2([G])$，都比 $G=\Gm$ 的情形复杂得多。

实流形 $\Gamma_i\backslash X$ 本身是很有趣的实流形，因为它们具有算术意义。我们将会看到，当 $G=\GL_2$ 时这些 Riemann 曲面是拟射影的，并典范地定义在数域上。

若空间 $[G]$ 紧（这当且仅当任何 $\Gm\subseteq G$ 都落在中心内，见 \texttt{[GH19, Theorem 2.6.2]}），则 $L^2([G])$ 的分析更简单。一个具体的例子是 $G$ 取 $\Q$ 上非分裂四元数代数的单位群，即具有表现

% ===== PDF p.009 =====
\[
\langle x,y\mid x^2=a,\ y^2=b,\ xy=-yx\rangle
\]
（适当选取 $a,b\in\Q^\times$）的 $\Q$-代数的单位群。注意当 $ab<0$ 时 $G(\R)\cong\GL_2(\R)$，因此 $\Gamma_i\backslash X$ 仍然是上半平面的商。四元数代数的更多细节见 \texttt{[GS17, Chapter 1]}。

$\Q$ 上的约化群大量存在。产生大量例子的一种特别容易的方式是通过 Weil 限制。为此设 $F/\Q$ 为有限扩张，$H$ 在 $F$ 上约化。则函子 $R\mapsto H(R\otimes_\Q F)$（定义在 $\Q$-代数范畴上）可由 $\Q$ 上的约化群
\[
G:=\Res_{F/\Q}(H)
\]
表出，见 \texttt{[CGP15, Appendix A]} 或 \texttt{[GH19, Section 1.4]}。对这个约化群 $G$ 有
\[
G(\Q)\backslash G(\A)=H(F)\backslash H(\A_F),
\]
其中 $\A_F$ 是 $F$ 的阿代尔环。这蕴涵：就朗兰兹纲领而言，考虑 $\Q$ 上所有约化群等价于考虑所有数域上所有约化群。

当 $H=\mathbb{G}_{m,F}$ 时，$L^2([G])$ 的期望分解（大致上）等价于 $F$ 的类域论。顺带一提：$[G]$ 体积的有限性包含了 $F$ 的类数有限性与 $F$ 上的 Dirichlet 单位定理。

另一个例子来源：设 $F$ 为虚二次域，$H=\GL_{2,F}$。此时 $G(\R)\cong\GL_2(\C)$，可取 $K_\infty\cong U(2)$（酉群），而
\[
X=A_G\backslash G(\R)/K_\infty
\]
是三维双曲空间 $\mathbb{H}^3$。此时的商 $\Gamma_i\backslash X$ 称为\textbf{算术双曲 3-流形}，见 \texttt{[Thu82, Section 4 \& 5]}。更精确地说，双曲 3-空间 $\mathbb{H}^3$ 可以模型在四元数
\[
\{q=x+y\cdot i+z\cdot j\mid x,y\in\R,\ z>0\}
\]
上，而 $\begin{pmatrix}a&b\\c&d\end{pmatrix}\in\GL_2(\C)$ 通过 $q\mapsto(aq+b)(cq+d)^{-1}$ 作用在 $\mathbb{H}^3$ 上，其中 $\C=\{x+y\cdot i\mid x,y\in\R\}$。

讨论完一些例子之后，我们现在可以给出朗兰兹纲领中居于中心地位的自守表示\idx{自守表示}（在 $L^2$ 意义下）的定义。

\noindent\textbf{定义 2.4}（\texttt{[GH19, Definition 3.3]}）\phantomsection\label{def:2.4}\quad 设 $G$ 为 $\Q$ 上的约化群。一个\textbf{自守表示}（在 $L^2$ 意义下）是不可约酉 $G(\A)$-表示 $\pi$，它同构于 $L^2([G])$ 的某个子商（subquotient）。

以后我们还会讨论另一种更代数的自守表示概念。在上述定义中，$\pi$ 必然在 $A_G$ 上平凡。这个假设通常是无害的：由不可约酉表示的 Schur 引理，对任意不可约酉 $G(\A)$-表示 $\pi$ 都能找到特征 $\chi:G(\R)\to\C^\times$，使得 $\chi\otimes\pi$ 在 $A_G$ 上平凡，见 \texttt{[GH19, Section 6.5]}。例如 $G=\GL_n$ 时总可以取 $\chi$ 为形如 $g\mapsto|\det(g)|^s$（某个 $s\in i\R$）的特征。

空间 $L^2([G])$ 分解为「离散」部分与「连续」部分，见 \texttt{[GH19, Section 10.4]}。从分析的角度看，$[G]$ 紧的情形更容易，因为此时连续部分消失。

现在我们更精确地解释「离散」部分与「连续」部分的含义。回忆 $\widehat{G(\A)}$ 记不可约

% ===== PDF p.010 =====
酉 $G(\A)$-表示的同构类集。类比于局部紧 abelian 群 $A$ 上 $\widehat A$ 的紧开拓扑，集合 $\widehat G$ 有一个自然的拓扑，即 \textbf{Fell 拓扑}，见 \texttt{[GH19, Section 3.8]}。粗略地说，Fell 拓扑可以这样描述：取 $\pi\in\widehat{G(\A)}$、$x\in\pi$ 以及连续 $\C$-线性同态 $\phi:\pi\to\C$，则
\[
f_{x,\phi}:G\to\C,\qquad g\mapsto\phi(gx)
\]
称为 $\pi$ 的一个\textbf{矩阵系数}。粗略地说，$\pi,\pi'\in\widehat{G(\A)}$ 在 Fell 拓扑中彼此接近，若它们的矩阵系数在紧开拓扑下彼此接近。

下面的定理对 $G(\A)$ 的酉表示更一般地成立。我们不需要定理 2.5 中积分的精确含义，只需指出它的出现是因为 $G(\A)$ 的某些表示「贡献」于 $L^2([G])$ 却不作为子表示出现（类似于 $\R$ 的 Fourier 理论中的情形）。

\noindent\textbf{定理 2.5}（\texttt{[GH19, Theorem 3.9.4]}）\phantomsection\label{thm:2.5}\quad 存在可测的重数函数 $m:\widehat G\to\{1,2,\dots,\infty\}$ 与 $\widehat G$ 上的测度 $\mu$，使得
\[
L^2([G])\cong\int_{\widehat G}\Bigl(\widehat{\bigoplus^{m(\pi)}}\pi\Bigr)\,\mathrm{d}\mu(\pi).
\]
$m$ 与 $\mu$ 在「零测集上作改变」的意义下唯一。

借助这个定理，我们现在可以粗略解释 $L^2([G])$ 的「离散」部分与「连续」部分：满足 $\mu(\pi)>0$ 的点 $\pi\in\widehat{G(\A)}$ 对 $\widehat{G(\A)}$ 上的测度 $\mu$ 而言是「离散」的，它们以重数 $m(\pi)$ 作为子表示出现在 $L^2([G])$ 中。这引出分解
\[
L^2([G])\cong L^2_{\mathrm{disc}}([G])\oplus L^2_{\mathrm{cont}}([G]),
\]
其中 $L^2_{\mathrm{disc}}([G])$ 是 $G(\A)$ 的全部不可约酉子表示的 Hilbert 和，$L^2_{\mathrm{cont}}([G])$ 是它的正交补，见 \texttt{[GH19, Section 9.1]}。

Langlands 在一项开创性的工作中用 $G$ 的（真）抛物子群\idx{抛物子群}描述了连续部分。回忆：按定义，闭连通子群 $P\subseteq G$（在 $\Q$ 上）称为\textbf{抛物子群}，若 $G/P$ 是 $\Spec(\Q)$ 上的 proper 概形；而 $P$ 对其单根取商所得为其 \textbf{Levi 商} $M$，见 \texttt{[GH19, Section 1.9]}。例如对 $\GL_n$，取分解 $n=n_1+n_2+\dots+n_k$，令 $P$ 为 $\GL_n$ 中块大小为 $n_1,\dots,n_k$ 的块上三角矩阵构成的子群；其 Levi 商为块对角矩阵构成的群。$\GL_n$ 的所有抛物子群在共轭下都呈这种形状。

Langlands 随后证明：连续部分 $L^2_{\mathrm{cont}}([G])$ 可以通过对 $P\subseteq G$（$P\neq G$）的所有 Levi 商的离散部分 $L^2_{\mathrm{disc}}([M])$ 作「诱导」来描述，见 \texttt{[GH19, Section 10.4]}。这就是他的「Eisenstein 级数\idx{Eisenstein 级数}」理论，也是朗兰兹纲领的一个起点。

对本课程而言，$L^2_{\mathrm{disc}}([G])$（实际上是其中恰当的「尖点」子空间）将比整个空间 $L^2([G])$ 更重要。顺带一提：$[G]$ 紧当且仅当 $G$ 没有（定义在 $\Q$ 上的）真抛物子群。这与前面所说的「此时 $L^2([G])=L^2_{\mathrm{disc}}([G])$」是一致的。

\chapter{从模形式到自守表示 I}

% ===== PDF p.011 =====
从本讲开始，我们想为 $\GL_2$ 构造与模形式相伴的自守表示。

先回忆上次讨论过的内容：
\begin{itemize}
\item 设 $G$ 是 $\Q$ 上的约化群，$\A=\Af\times\R$ 是阿代尔环。
\item 我们引入空间
\[
[G]=A_GG(\Q)\backslash G(\A),
\]
称之为\textbf{阿代尔商}（adélic quotient），它具有有限的 $G(\A)$-不变体积。
\item 我们陈述过：朗兰兹纲领的目标是按算术数据分解酉 $G(\A)$-表示
\[
L^2([G]).
\]
\item 在例子 $G=\Gm$ 中我们证明了
\[
L^2(\R_{>0}\Q^\times\backslash\A^\times)\cong\widehat{\bigoplus_{\sigma:\Gal(\overline\Q/\Q)\to\C^\times}}\C\chi_\sigma,
\]
作为「所期望的算术数据来自 $\Gal(\overline\Q/\Q)$ 的 Galois 表示」的一个例证。
\item 随后我们讨论了 $[G]$ 的几何：
\begin{itemize}
\item 取 $K_\infty\subseteq G(\R)$ 为某个极大紧连通子群，
\item 取 $K\subseteq G(\Af)$ 为某个紧开子群。
\item 则 $[G]\to[G]/K$ 是 $K$-挠子（特别地，是一个 profinite 覆盖），$[G]/K\to[G]/KK_\infty$ 是 $K_\infty$-挠子，而当 $K$ 充分小时 $[G]/KK_\infty$ 是实流形；更精确地，
\[
[G]/KK_\infty=\bigsqcup_{i=1}^m\Gamma_i\backslash X
\]
是若干算术流形的无交并，其中 $\Gamma_i\subseteq G(\R)$ 是由 $K$ 决定的（离散）同余子群，而
\[
X:=A_G\backslash G(\R)/K_\infty.
\]
\end{itemize}
\item 此外我们定义：$G(\A)$ 的不可约酉表示 $\pi$ 称为\textbf{自守的}（在 $L^2$ 意义下），若它同构于 $L^2([G])$ 的某个子商。
\item 最初的例子（当时只是隐含提及）如下：
\begin{itemize}
\item 若 $G=\Gm$，则 $G(\A)$ 的自守表示（在 $L^2$ 意义下）恰是 $\A^\times$ 的酉特征，且这些特征在 $\Q^\times\cdot\R_{>0}$ 上平凡。
\item 若 $G$ 是一般的群，则平凡表示（即 $[G]$ 上的常函数）是自守的，因为 $[G]$ 有有限体积；并且它落在 $L^2([G])$ 的离散部分 $L^2_{\mathrm{disc}}([G])$ 中。
\end{itemize}
\item 对一般的 $G$，存在正交分解
\[
L^2([G])\cong L^2_{\mathrm{disc}}([G])\oplus L^2_{\mathrm{cont}}([G]),
\]
其中 $L^2_{\mathrm{disc}}([G])$ 是 $G(\A)$ 的所有不可约子表示之和的闭包。
% ===== PDF p.012 =====
\item 最后我们提到，Langlands 用 $L^2_{\mathrm{disc}}([M])$ 描述了 $L^2_{\mathrm{cont}}([G])$，其中 $M$ 遍历抛物子群 $P\subseteq G$（$P\neq G$，且定义在 $\Q$ 上）的 Levi 商。
\end{itemize}

今天我们将更详细地讨论 $G=\GL_2$ 的情形。首先，
\[
A_{\GL_2}\cong\R_{>0}
\]
（作为标量矩阵嵌入 $\GL_2(\R)$）。上次我们断言
\[
\GL_2(\Q)\backslash\GL_2(\A)\cong\GL_2(\Z)\backslash(\GL_2(\widehat\Z)\times\GL_2(\R)),
\]
其中
\[
\widehat\Z:=\varprojlim_m\Z/m\cong\prod_p\Z_p.
\]
下面解释它为什么成立，参见 \texttt{[Tho, Theorem 1]}。

群 $\GL_2(\Af)$ 是诸群 $\GL_2(\Q_p)$ 关于紧开子群 $\GL_2(\Z_p)$ 的限制乘积，即
\[
\GL_2(\Af)=\Bigl\{(A_p)_p\in\prod_p\GL_2(\Q_p)\Bigm| A_p\in\GL_2(\Z_p)\ \text{除有限多个 }p\text{ 外}\Bigr\}
\]
（参见 \texttt{[GH19, Proposition 2.3.1]}）。中国剩余定理蕴涵：$\Q\subseteq\Af$（对角嵌入）是稠密的。此外，由素因子分解，
\[
\Af^\times=\Q^\times\cdot\widehat\Z^\times.
\]

群 $\GL_2(\Z_p)$、$\GL_2(\Q_p)$ 由初等矩阵与对角矩阵生成，这蕴涵
\[
\GL_2(\Af)=\GL_2(\Q)\GL_2(\widehat\Z),
\]
从而
\[
\GL_2(\Q)\backslash\GL_2(\Af)\cong\GL_2(\Q)\cap\GL_2(\widehat\Z)\backslash\GL_2(\widehat\Z)\cong\GL_2(\Z)\backslash\GL_2(\widehat\Z).
\]
由此可推出所想要的\znote{此处关键的一步是 $\SL_2$ 的强逼近定理：$\SL_2(\Q)\SL_2(\widehat\Z)=\SL_2(\Af)$，即 $\SL_2(\Q)$ 在 $\SL_2(\Af)$ 中是稠密的；再配合 $\Af^\times=\Q^\times\widehat\Z^\times$ 处理行列式，就得到上面的双陪集分解。原讲义把这一步归为「初等矩阵与对角矩阵生成」，细节可参见 \texttt{[Tho, Theorem 1]}。[来源：原文第 12 页与 \texttt{[Tho, Theorem 1]}]}
\[
\GL_2(\Q)\backslash\GL_2(\A)\cong\GL_2(\Z)\backslash(\GL_2(\widehat\Z)\times\GL_2(\R)).
\]
上述证明其实对一切 $n\ge1$ 都成立，并表明
\[
\GL_n(\Q)\backslash\GL_n(\A)\cong\GL_n(\Z)\backslash(\GL_n(\widehat\Z)\times\GL_n(\R)).
\]

对 $m\ge1$ 定义
\[
K(m)=\ker(\GL_2(\widehat\Z)\to\GL_2(\Z/m)).
\]
这些群在所有紧开子群 $\GL_2(\Af)$ 中共尾。我们有
\[
\GL_2(\Q)\backslash\GL_2(\Af)/K(m)\cong\GL_2(\Z)\backslash\GL_2(\widehat\Z)/K(m),
\]
且 $\GL_2(\widehat\Z)/K(m)\cong\GL_2(\Z/m)$。群 $\SL_2(\Z/m)$ 甚至由初等矩阵生成\znoteplain{原文脚注：对某个环 $R$ 中的任意可逆元 $a$，利用欧几里得算法与如下「神奇恒等式」：
\[
\begin{pmatrix}a&0\\0&a^{-1}\end{pmatrix}
=\begin{pmatrix}1&a\\0&1\end{pmatrix}
\begin{pmatrix}1&0\\-a^{-1}&1\end{pmatrix}
\begin{pmatrix}1&a\\0&1\end{pmatrix}
\begin{pmatrix}0&-1\\1&0\end{pmatrix}.
\]}。这蕴涵态射 $\SL_2(\Z)\to\SL_2(\Z/m)$ 是满射。于是我们得到
\[
\GL_2(\Z)\backslash\GL_2(\widehat\Z)/K(m)\cong\{\pm1\}\backslash(\Z/m)^\times,
\]
% ===== PDF p.013 =====
同构由行列式给出。现在定义
\[
K_\infty:=\SO_2(\R)\subseteq\GL_2(\R)
\]
以及
\[
X:=A_{\GL_2}\backslash\GL_2(\R)/K_\infty.
\]
则 $X\cong\HH^\pm$（上半平面与下半平面之并），同构由态射
\[
A_{\GL_2}\backslash\GL_2(\R)/K_\infty\xrightarrow{\ \sim\ }\HH^\pm,\qquad
g=\begin{pmatrix}a&b\\c&d\end{pmatrix}\mapsto g(i):=\frac{a\cdot i+b}{c\cdot i+d}
\]
给出。于是最终得到
\[
\GL_2(\Q)\backslash(\GL_2(\Af)/K(m)\times X)\cong\bigsqcup_{\{\pm1\}\backslash(\Z/m)^\times}\Gamma(m)\backslash\HH^\pm,
\]
其中
\[
\Gamma(m):=K(m)\cap\GL_2(\Z)=\ker(\GL_2(\Z)\to\GL_2(\Z/m)).
\]

$\GL_2(\Q)\backslash(\GL_2(\Af)/K(m)\times X)$ 不连通这件事也许显得烦人，但后面会表明这其实是优点。注意：若把 $K\subseteq\GL_2(\Af)$ 换成某个使
\[
\det:K\to\widehat\Z^\times
\]
为满射的紧开子群 $K$，例如取
\[
K_1(m):=\Bigl\{\begin{pmatrix}a&b\\c&d\end{pmatrix}\in\GL_2(\Af)\Bigm| c\equiv0,\ d\equiv1\bmod m\Bigr\},
\]
则 $\GL_2(\Q)\backslash(\GL_2(\Af)/K\times X)$ 连通。

$G=\GL_2$ 时 $[G]$ 体积的有限性，现在可以用直接论证来检验。$K_\infty K(m)$ 的紧性蕴涵：只需证明空间
\[
[\GL_2]/K_\infty K(m)
\]
有有限体积。把这个空间写成无交并，甚至可以归约到「$\Gamma(m)\backslash\HH^\pm$ 对 $m\ge1$ 有有限体积」这一陈述。至多差 $\R_{>0}$ 中的一个标量，$\HH^\pm$ 上 $\GL_2(\R)$-不变测度由体积形式
\[
\frac{1}{y^2}\,\mathrm{d}x\wedge\mathrm{d}y
\]
给出。确实，我们须证明对 $g\in\GL_2(\R)$ 有
\[
g^*\Bigl(\frac{1}{y^2}\,\mathrm{d}x\wedge\mathrm{d}y\Bigr)=\frac{1}{y^2}\,\mathrm{d}x\wedge\mathrm{d}y.
\]
现在把 $\mathrm{d}x\wedge\mathrm{d}y$ 写成 $\frac{i}{2}\,\mathrm{d}z\wedge\mathrm{d}\bar z$，其中 $z=x+iy$。容易验证上述陈述对
\[
g=\begin{pmatrix}1&a\\0&1\end{pmatrix},\qquad a\in\R,
\]
对 $A_{\GL_2}$，对
\[
g=\begin{pmatrix}a&0\\0&1\end{pmatrix},\qquad a\in\R^\times,
\]
最后对
\[
g=\begin{pmatrix}0&1\\1&0\end{pmatrix},
\]
即对 $z\mapsto\frac{1}{z}=\frac{\bar z}{|z|^2}$ 成立。由前面提到的神奇恒等式，可以推出对下三角矩阵的不变性，从而对 $\GL_2(\R)$ 全体成立。

% ===== PDF p.014 =====
现在证明 $\Gamma(m)\backslash\HH^\pm$ 有有限体积。$\SL_2(\Z)$ 在 $\HH$ 上的一个已知基本域由集合
\[
D:=\{z\in\HH\mid |z|>1,\ |\operatorname{Re}(z)|<1/2\}
\]
给出，而 $\GL_2(\Z)\backslash\HH^\pm$ 体积的有限性由
\[
\int_D\frac{1}{y^2}\,\mathrm{d}x\wedge\mathrm{d}y<\infty
\]
推出。最后，由于 $\Gamma(m)\subseteq\GL_2(\Z)$ 是有限指数的，$\Gamma(m)\backslash\HH^\pm$ 有有限体积，如所愿。

朗兰兹纲领尤其包含模形式与 $\Gal(\overline\Q/\Q)$ 的二维（$\ell$-进）Galois 表示之间的紧密联系；讨论这一联系将是本课程的一大主题。首先回忆模形式的经典定义，参见 \texttt{[DS05]}。为此，设 $\Gamma\subseteq\SL_2(\Z)$ 为同余子群（后面假设它充分小），并固定 $k\in\Z$。

\noindent\textbf{定义 3.1}\phantomsection\label{def:3.1}\quad 函数 $f:\HH^\pm\to\C$ 称为 $\Gamma$ 的\textbf{权 $k$ 模形式}，若
\begin{itemize}
\item 对 $\gamma=\begin{pmatrix}a&b\\c&d\end{pmatrix}\in\Gamma$ 有 $f(\gamma\cdot z)=(cz+d)^kf(z)$；
\item $f$ 全纯；
\item $f$ 在 $\Gamma$ 的尖点处全纯（稍后解释）。
\end{itemize}
令 $M_k(\Gamma)$ 为 $\Gamma$ 的权 $k$ 模形式构成的 $\C$-向量空间。只满足前两个条件、而不必在尖点处全纯的函数，我们也称为\textbf{弱模形式}。

回忆：上半平面 $\HH$ 自然 $\SL_2(\Q)$-等变地嵌入集合
\[
\HH^*:=\HH\cup\Pone(\Q),
\]
后者带 Satake 拓扑。同余子群 $\Gamma\subseteq\SL_2(\Z)$ 在 $\Pone(\Q)$ 上的轨道称为 $\Gamma$ 的\textbf{尖点}，商
\[
\Gamma\backslash\HH^*
\]
是 $\Gamma\backslash\HH$ 的自然紧化。把 $\Pone(\Q)$ 的另一份加到下半平面 $\HH^-$ 上，对 $\Gamma\subseteq\GL_2(\Z)$ 就得到类似的对象与概念。

「在尖点处全纯」这一条件现在意味着：对每个尖点 $\sigma$，$\Gamma$ 的弱模形式 $f$ 都有 Fourier 展开，即把该尖点平移到 $\infty\in\Pone(\Q)$ 之后，存在等式
\[
f(z)=\sum_{n\in\Z}a_nq^{n/h},
\]
其中 $q=e^{2\pi iz}$，$h$ 是所谓 $f$ 的\textbf{最小周期}，$a_n\in\C$，$n\in\Z$。若 $n<0$ 时 $a_n=0$，则称弱模形式 $f$ 在该尖点处全纯；若它在每个尖点处都全纯，则称它在尖点处全纯。若进一步对每个尖点都有 $a_0=0$，则称 $f$ 为\textbf{尖点形式}。记
\[
S_k(\Gamma)\subseteq M_k(\Gamma)
\]
为 $\Gamma$ 的权 $k$ 尖点形式子空间。

下一个目标是把 $\Gamma\subseteq\GL_2(\Z)$ 的权 $k\in\Z$ 的（弱）模形式典范地描述为某个全纯线丛
% ===== PDF p.015 =====
$\omega^{\otimes k}$ 在 $\Gamma\backslash\HH^\pm$ 上（至少当 $\Gamma$ 充分小时）的截面。

做这件事之前，我们先讨论一些有用的「纤维丛的下降」构造，它们会被频繁使用。

考虑一个集合 $Z$，以及一个自由右作用在 $Z$ 上的群 $H$。则投影
\[
Z\to Z/H
\]
配得上称为 $H$-\textbf{挠子}。设 $F$ 是带 $H$（左）作用的集合。$Z$ 关于 $H$ 与 $F$ 的\textbf{收缩积}（contracted product）
\[
Z\times^HF
\]
定义为 $Z\times F$ 对 $H$ 的作用 $h\cdot(z,f)\mapsto(zh^{-1},hf)$ 的商。经由自然投影 $(z,f)\mapsto z$，我们得到 $Z/H$ 上的「纤维丛」
\[
Z\times^HF\to Z/H,
\]
其纤维（非典范地）同构于 $F$。确实，对类 $c\in Z/H$ 与满足 $z\mapsto c$ 的 $z\in Z$ 的选取，得到同构
\[
F\cong(Z\times^HF)\times_{Z/H}c,\qquad f\mapsto(z,f).
\]
注意 $-\times^H-$ 表示收缩积，而 $-\times_{Z/H}-$ 表示纤维积（这是两个互不相干的构造）。特别地，存在典范同构
\[
Z\times F\cong(Z\times^HF)\times_{Z/H}Z,\qquad(z,f)\mapsto((z,f),z),
\]
即纤维丛 $Z\times^HF$ 拉回到 $Z$ 之后典范地平凡化。\znoteplain{原文此处说拉回到 $Z/H$；但按上式应为拉回到 $Z$（$Z/H$ 是底空间）。}

我们可以方便地描述纤维丛 $Z\times^HF\to Z/H$ 的截面。

\noindent\textbf{引理 3.2}\phantomsection\label{lem:3.2}\quad 截面 $s:Z/H\to Z\times^HF$（即 $Z/H$ 上的态射）的集合，典范地等同于满足
\[
\phi(zh)=h^{-1}\phi(z)
\]
（对一切 $z\in Z$，$h\in H$）的函数 $\phi:Z\to F$ 的集合。

\noindent\textbf{证明}．给定 $s:Z/H\to Z\times^HF$ 与 $z\in Z$，由 $H$ 在 $Z$ 上作用的自由性，存在唯一元素 $f\in F$ 使
\[
s(z)=(z,f).
\]
令 $\phi(z):=f$。则
\[
(z,\phi(z))=s(z)=s(zh)=(zh,\phi(zh))=(z,h\phi(zh)),
\]
这蕴涵
\[
\phi(z)=h\phi(zh)
\]
如所愿。反之，给定满足 $\phi(zh)=h^{-1}\phi(z)$（对一切 $z\in Z$，$h\in H$）的函数 $\phi:Z\to F$，函数
\[
Z\to Z\times F,\qquad z\mapsto(z,\phi(z))
\]
按定义对 $H$ 的作用不变，从而下降为截面 $s:Z/H\to Z\times^HF$。两个构造互为逆。\hfill$\square$

% ===== PDF p.016 =====
反之，假设给定态射 $Y\to Z/H$，以及 $Z$ 上的同构
\[
g:Y\times_{Z/H}Z\xrightarrow{\ \sim\ }Z\times F
\]
（对某个集合 $F$）。取 $z\in Z$，其像为 $c\in Z/H$。则对任意 $h\in H$ 得到同构
\[
\sigma_h:F\xrightarrow{\ g^{-1}|_{z\times F}\ }Y_c\times_cz\xrightarrow{\ \cdot h\ }Y_c\times_czh\xrightarrow{\ g|_{g^{-1}(zh\times F)}\ }F,
\]
它们定义了 $H$ 在 $F$ 上的作用。此外可以验证，通过这个作用
\[
Y\cong Z\times^HF.
\]
用更花哨的话说：带 $Z$ 上「平凡化」的「纤维丛」$Y\to Z/H$，等价于带 $H$ 作用的集合。

我们再提一个变体：当 $Z=G$ 是群、$H\subseteq G$ 是子群时。$G/H$ 上的 $G$-\textbf{等变集合} $Y$ 是 $G/H$ 上的空间 $\pi:Y\to G/H$，连同 $G$ 的左作用，使 $\pi$ 为 $G$-等变。此时有典范平凡化
\[
G\times^{1\cdot H}Y\cong Y\times_{G/H}G,\qquad(g,y)\mapsto(gy,g),
\]
由此可推出：$G/H$ 上 $G$-等变集合的范畴恰好等价于带 $H$（左）作用的集合的范畴，因为对 $H$-集合 $F$，收缩积
\[
G\times^HF
\]
自然地是 $G$-等变的（让 $G$ 通过左乘作用在左因子上）。注意我们要求 $H$ 右作用在 $Z$（相应地 $G$）上，但当然也可以类似地考虑 $H\backslash Z$ 上的「纤维丛」。

我们把构造放在集合范畴中，但显然若 $Z$、$H$、$F$ 是拓扑空间（实流形、复流形……）且 $H$ 的作用连续（光滑、全纯……），我们可以在拓扑空间（实流形、复流形……）等范畴中作同样的构造。

现在开始把模形式描述为模曲线\idx{模曲线}上的线丛的截面。

首先注意，线丛
\[
\OO_{\Pone_\C}(k)
\]
（$k\in\Z$）自然地是 $\GL_2(\C)$-等变的。事实上，取张量幂后只需考虑 $k=1$。但 $\OO_{\Pone_\C}(1)$ 相伴于自然的 $\Gm$-挠子
\[
\mathbb{A}^2_\C\setminus\{0\}\to\Pone_\C,\qquad(z_1,z_2)\mapsto[z_1:z_2],
\]
后者自然地 $\GL_2(\C)$-等变。其次注意，嵌入
\[
\HH^\pm\subseteq\Pone_\C,\qquad z\mapsto[z:1]
\]
是 $\GL_2(\R)$-等变的，于是对 $k\in\Z$ 得到 $\GL_2(\R)$-等变线丛
\[
\OO_{\HH^\pm}(k):=\OO_{\Pone_\C}(k)\times_{\Pone_\C}\HH^\pm,
\]
它位于 $\HH^\pm$ 上。

% ===== PDF p.017 =====
$i\in\HH^\pm$ 在 $\GL_2(\R)$ 中的稳定化子等同于
\[
\C^\times\cong\R_{>0}\SO_2(\R)\subseteq\GL_2(\R).
\]
利用
\[
\HH^\pm\cong\GL_2(\R)/\C^\times,
\]
在先前提到的纤维丛构造下，$\OO_{\HH^\pm}(k)$ 对应于（全纯）表示
\[
\chi_k:\C^\times\to\GL_1(\C)=\C^\times,\qquad z\mapsto z^k.
\]
注意引理 3.2 蕴涵：$\OO_{\HH^\pm}(k)$ 在 $\HH^\pm$ 上的光滑截面等同于满足
\[
\phi(gz)=z^{-k}\phi(g),\qquad g\in\GL_2(\R),\ z\in\C^\times
\]
的光滑函数 $\phi:\GL_2(\R)\to\C$。这一点后面会用到。从现在起我们假设 $\Gamma$ 充分小，即它自由作用在 $\HH^\pm$ 上。对 $\Gamma$ 在 $\OO_{\HH^\pm}(k)$ 上的作用取商，得到 $\Gamma\backslash\HH^\pm$ 上的全纯线丛
\[
\omega^{\otimes k}:=\Gamma\backslash\OO_{\HH^\pm}(k).
\]
为把 $\omega^{\otimes k}$ 的截面与 $\Gamma$ 的权 $k$（弱）模形式联系起来，我们须把 $\GL_2(\R)$ 在 $\OO_{\HH^\pm}(k)$ 上的作用写得更显式。对 $k\in\Z$ 定义 $\HH^\pm$ 上的全纯 $\GL_2(\R)$-等变线丛
\[
L_k:=\HH^\pm\times\C,
\]
$\GL_2(\R)$ 通过
\[
g\cdot(z,\lambda):=\Bigl(\frac{az+b}{cz+d},(cz+d)^k\lambda\Bigr)
\]
左作用，其中 $g=\begin{pmatrix}a&b\\c&d\end{pmatrix}\in\GL_2(\R)$。为记号方便，引入函数（一个「自守因子」）\znote{自守因子 $j(g,z)=cz+d$ 有简单的几何意义：$\operatorname{Im}(gz)=\det(g)|j(g,z)|^{-2}\operatorname{Im}(z)$，且它满足余圈关系 $j(gh,z)=j(g,hz)j(h,z)$。权 $k$ 的变换律正是用 $j$ 的 $k$ 次幂写出的。[来源：原文第 17、25 页]}
\[
j:\GL_2(\R)\times\HH^\pm\to\C^\times,\qquad
\Bigl(g=\begin{pmatrix}a&b\\c&d\end{pmatrix},z\Bigr)\mapsto j(g,z):=cz+d.
\]
留作练习：验证映射
\[
L_1^*\to\OO_{\HH^\pm}(1)^*=\HH^\pm\times_{\Pone_\C}\mathbb{A}^2_\C\setminus\{0\},\qquad
(z,\lambda)\mapsto(z,(\lambda z,\lambda))
\]
是 $\HH^\pm$ 上 $\GL_2(\R)$-等变线丛的全纯同构。这里（以及下文）$(-)^*$ 表示线丛中零截面的补集。取张量幂得到
\[
L_k\cong\OO_{\HH^\pm}(k)
\]
% ===== PDF p.018 =====
对一切 $k\in\Z$ 成立，作为 $\GL_2(\R)$-等变全纯线丛\znoteplain{原文脚注：$L_k$ 与 $\OO_{\HH^\pm}(k)$ 作为 $\GL_2(\R)$-等变线丛抽象同构，也可由观察到两者都对应于 $\C^\times$-表示 $z\mapsto z^k$ 推出。这一点已对 $\OO_{\HH^\pm}$ 指出，其成立是因为在嵌入 $\C^\times\hookrightarrow\GL_2(\R)$ 下，元素 $z\in\C^\times$ 映到矩阵 $g=\begin{pmatrix}\operatorname{Re}(z)&-\operatorname{Im}(z)\\\operatorname{Im}(z)&\operatorname{Re}(z)\end{pmatrix}$，且 $j(g,i)=z$。}。

与引理 3.2 类似\znoteplain{原文脚注：这里的情形与引理 3.2 略有不同，即这些 $\Gamma$-等变向量丛并不相伴于 $\Gamma$ 的某个表示。}，我们得到：全纯截面
\[
\HH^\pm\to L_k,\qquad z\mapsto(z,f(z))
\]
是 $\Gamma$-等变的，当且仅当 $f(z)$ 满足
\[
f(\gamma\cdot z)=j(\gamma,z)^kf(z)
\]
（$\gamma\in\Gamma$），即
\[
f\Bigl(\frac{az+b}{cz+d}\Bigr)=(cz+d)^kf(z),
\]
只要 $\gamma=\begin{pmatrix}a&b\\c&d\end{pmatrix}\in\Gamma$。换言之，我们证明了
\[
H^0(\Gamma\backslash\HH^\pm,\omega^{\otimes k})
\]
等同于 $\Gamma$ 的弱模形式空间。特别地，有自然嵌入
\[
M_k(\Gamma)\subseteq H^0(\Gamma\backslash\HH^\pm,\omega^{\otimes k}),
\]
其像由「在尖点处全纯」这一条件定义。下面分析在尖点处全纯对 $\omega^{\otimes k}$ 的截面意味着什么。

我们断言 $\omega^{\otimes k}$ 典范地延拓到 $\Gamma\backslash\HH^\pm$ 的（Satake）紧化 $\overline{\Gamma\backslash\HH^\pm}$。尖点 $\infty\in\Pone(\Q)$ 在 $\SL_2(\Z)$ 中的稳定化子 $\Gamma_\infty$ 是元素
\[
\begin{pmatrix}1&a\\0&1\end{pmatrix},\qquad a\in\Z
\]
构成的子群。因此只需看 $\omega$ 沿嵌入
\[
\Gamma_\infty\backslash(\HH\cup\{\infty\})
\]
的（典范）延拓。记
\[
L_k^+:=L_k|_{\HH}.
\]
则 $\Gamma_\infty$ 在 $L_k^+=\HH\times\C$ 上的作用只通过第一个因子（因为对一切 $z\in\HH$ 与任意 $\gamma\in\Gamma_\infty$ 有 $j(\gamma,z)=1$）。注意经由指数映射
\[
\Gamma_\infty\backslash\HH\cong\mathbb{D}^\times:=\{q\in\C^\times\mid 0<|q|<1\}
\]
（讨论 Fourier 展开时已隐含使用了这一点）。于是典范地
\[
\Gamma_\infty\backslash L_k^+\cong\mathbb{D}^\times\times\C,
\]
它显然延拓到 $\mathbb{D}:=\{q\in\C\mid|q|<1\}$，即延拓为 $\mathbb{D}\times\C$。这便证明了 $\omega^{\otimes k}$ 典范地延拓到 $\Gamma\backslash\HH^\pm$ 的上述（Satake）紧化 $\overline{\Gamma\backslash\HH^\pm}$。我们把这个延拓仍记作 $\omega^{\otimes k}$。全纯性条件
% ===== PDF p.019 =====
在尖点处全纯，现在翻译为：截面 $f\in H^0(\Gamma\backslash\HH^\pm,\omega^{\otimes k})$ 延拓为 $\overline{\Gamma\backslash\HH^\pm}$ 上 $\omega^{\otimes k}$ 的截面\znote{「在尖点处全纯」等价于 $q$-展开 $f=\sum_{n\in\Z}a_nq^{n/h}$ 中不出现负幂次；而「在尖点处全纯且常数项为零」则进一步要求 $a_0=0$。几何上，前者恰是截面沿紧化边界全纯延拓，后者是截面在边界除子 $D$ 上消没。[来源：原文第 14、19 页]}。

再多用一点理论：Riemann 曲面 $\Gamma\backslash\HH^\pm$ 是 $\C$ 上的代数曲线，其典范紧化\idx{典范紧化}的解析化是 $\overline{\Gamma\backslash\HH^\pm}$。由于 $M_k(\Gamma)$ 是 $H^0(\Gamma\backslash\HH^\pm,\omega^{\otimes k})$ 在 $H^0(\overline{\Gamma\backslash\HH^\pm},\omega^{\otimes k})$ 中的像，由 GAGA 可推出 $M_k(\Gamma)$ 在 $\C$ 上有限维。此外可以证明 $\omega$ 在 $\Gamma\backslash\HH^\pm$ 上丰沛，这蕴涵 $k<0$ 时 $M_k(\Gamma)=0$，而 $k\gg0$ 时 $M_k(\Gamma)$ 变得「大」。事实上，用 Riemann--Roch 定理可以计算 $M_k(\Gamma)$ 的维数。最后注意
\[
S_k(\Gamma)\cong H^0(\Gamma\backslash\HH^\pm,\omega^{\otimes k}(-D)),
\]
其中 $D=\overline{\Gamma\backslash\HH^\pm}\setminus\Gamma\backslash\HH^\pm$ 是无穷远处的（既约）除子。更多细节参见 \texttt{[DS05, Chapter 3, Chapter 7]}。

现在解释如何从 $\omega^{\otimes k}$ 的截面过渡到 $\GL_2(\A)$（或 $[\GL_2]$）上的函数。为此，我们把 $\Gamma\subseteq\GL_2(\Z)$ 遍历所有同余子群时的模形式，解释为 $K\subseteq\GL_2(\Af)$ 遍历所有紧开子群时的「模形式」。

设 $K\subseteq\GL_2(\Af)$ 为紧开子群，并写
\[
\GL_2(\Q)\backslash(\GL_2(\Af)/K\times\HH^\pm)\cong\bigsqcup_{i=1}^m\Gamma_i\backslash\HH^\pm,
\]
其中 $\Gamma_i\subseteq\GL_2(\Z)$ 是某些同余子群。则由 $\HH^\pm$ 拉回，$\OO_{\HH^\pm}(k)$ 定义了
\[
\GL_2(\Q)\backslash(\GL_2(\Af)/K\times\HH^\pm)
\]
上的（复）线丛 $\omega^{\otimes k}$（我们隐含假设 $K$ 充分小，并用了 $\OO_{\HH^\pm}(k)$ 的 $\GL_2(\Q)$-等变性）。

$\Gamma_i$（$i=1,\dots,m$）的模形式空间 $M_k(\Gamma_i)$（甚至弱模形式）嵌入到全纯截面空间
\[
H^0(\GL_2(\Q)\backslash(\GL_2(\Af)/K\times\HH^\pm),\omega^{\otimes k}).
\]
令
\[
M_k(K):=\bigoplus_{i=1}^mM_k(\Gamma_i).
\]
则 $M_k(K)\subseteq H^0(\GL_2(\Q)\backslash(\GL_2(\Af)/K\times\HH^\pm),\omega^{\otimes k})$ 同样由「在所有尖点处全纯」这一条件定义。我们甚至可以去掉 $K$，用 $\omega^{\otimes k}$ 的拉回定义
\[
\GL_2(\Q)\backslash(\GL_2(\Af)\times\HH^\pm)
\]
上的 $G(\Af)$-等变线丛，仍记作 $\omega^{\otimes k}$。注意这个空间不再是复流形，而只是复流形沿有限覆盖映射的逆向极限。但仍可通过
\[
H^0(\GL_2(\Q)\backslash(\GL_2(\Af)\times\HH^\pm),\omega^{\otimes k})
:=\varinjlim_KH^0(\GL_2(\Q)\backslash(\GL_2(\Af)/K\times\HH^\pm),\omega^{\otimes k})
\]
定义它上面 $\omega^{\otimes k}$ 的「全纯截面」的类似物。

% ===== PDF p.020 =====
此外，我们得到（大空间）
\[
M_k:=\varinjlim_KM_k(K),
\]
即权 $k$ 的模形式。\znoteplain{原文写作 “modular forms of level $k$”，其中 “level” 应为 “weight”（级与权是两个概念）。}后面我们会解释：对（充分小的）$K\subseteq\GL_2(\Af)$，$M_k(K)$ 是 $M_k$ 上某个典范 $K$-作用的 $K$-不变量。

由围绕引理 3.2 的纤维丛讨论我们知道：$\OO_{\HH^\pm}(k)$ 沿
\[
\GL_2(\R)\to\HH^\pm,\qquad g\mapsto g\cdot i
\]
的拉回典范平凡（作为 $\GL_2(\R)$-等变丛），且 $\OO_{\HH^\pm}(k)$ 在 $\HH^\pm$ 上的光滑截面典范等同于满足
\[
\phi(gz)=z^{-k}\phi(g),\qquad g\in\GL_2(\R),\ z\in\C^\times
\]
的光滑函数 $\phi:\GL_2(\R)\to\C$。具体地，把光滑截面
\[
\HH^\pm\to L_k=\HH^\pm\times\C,\qquad z\mapsto(z,f(z))
\]
对应到光滑函数
\[
\phi_f:\GL_2(\R)\to\C,\qquad g\mapsto j(g,i)^{-k}f(gi).
\]
注意 $L_k$ 在 $\HH^\pm$ 上的光滑截面就是 $\HH^\pm$ 上的光滑函数（但这个等同不是 $\GL_2(\R)$-等变的！）。

若把光滑性概念略加推广，我们可以把 $\HH^\pm$ 换成 $\GL_2(\Af)\times\HH^\pm$ 做类似的事。

\noindent\textbf{定义 3.3}\phantomsection\label{def:3.3}\quad 设 $Y_i$（$i\in I$）是实流形的共滤逆向系统，使每个过渡映射 $Y_i\to Y_j$ 都是有限覆盖。令 $Y:=\varprojlim_{i\in I}Y_i$，$\pi_i:Y\to Y_i$ 为典范投影。则称函数
\[
\phi:Y\to\C
\]
为\textbf{光滑的}，若存在 $i\in I$ 与光滑函数 $\phi_i:Y_i\to\C$ 使 $\phi=\phi_i\circ\pi_i$。

注意按同样的模式，若每个 $Y_i$ 还是复流形，也可定义 $Y$ 上的全纯函数（或线丛截面）。

当然，我们想把这个术语用到 $\GL_2(\Af)\times\HH^\pm$（或 $\GL_2(\A)=\GL_2(\Af)\times\GL_2(\R)$）上，它是复流形 $\GL_2(\Af)/K\times\HH^\pm$（相应地实流形 $\GL_2(\Af)/K\times\GL_2(\R)$）的逆向极限，其中 $K$ 遍历 $\GL_2(\Af)$ 的紧开子群的共滤系统。我们得到：光滑截面
\[
\GL_2(\Af)\times\HH^\pm\to\GL_2(\Af)\times L_k,\qquad(g,z)\mapsto(g,z,f(g,z))
\]
等同于满足
\[
\phi(g,g_\infty z)=z^{-k}\phi(g,g_\infty)
\]
% ===== PDF p.021 =====
（对 $(g,g_\infty)\in\GL_2(\A)$ 与 $z\in\C^\times\subseteq\GL_2(\R)$）的光滑函数
\[
\phi:\GL_2(\Af)\times\GL_2(\R)\to\C.
\]
具体地，对 $f$ 令
\[
\phi_f(g,g_\infty):=j(g_\infty,i)^{-k}f(g,g_\infty i).
\]
注意截面
\[
\GL_2(\Af)\times\HH^\pm\to\GL_2(\Af)\times L_k,\qquad(g,z)\mapsto(g,z,f(g,z))
\]
在 $\GL_2(\Q)$ 的左作用下不变，当且仅当 $f$ 满足模性条件
\[
f(\gamma g,\gamma z)=j(\gamma,z)^kf(g,z).
\]
利用等式
\[
j(\gamma g_\infty,i)=j(\gamma,g_\infty i)j(g_\infty,i),
\]
我们得到：$f$ 定义 $L_k$ 的 $\GL_2(\Q)$-等变截面，当且仅当函数 $\phi_f$ 是 $\GL_2(\Q)$-不变的，即
\[
\phi_f(\gamma g,\gamma g_\infty)=\phi_f(g,g_\infty)
\]
对 $(g,g_\infty)\in\GL_2(\A)$ 成立。

限制到全纯截面，我们于是构造了自然态射
\[
H^0(\GL_2(\Q)\backslash(\GL_2(\Af)\times\HH^\pm),\omega^{\otimes k})\to C^\infty(\GL_2(\Q)\backslash\GL_2(\A))
\]
使 $f(g,z)\mapsto\phi_f(g,g_\infty):=j(g_\infty,i)^{-k}f(g,g_\infty i)$；这里我们把 $\GL_2(\Q)\backslash(\GL_2(\Af)\times\HH^\pm)$ 上 $\omega^{\otimes k}$ 的截面看作满足上述模性条件的函数
\[
\GL_2(\Af)\times\HH^\pm\to\C,\qquad(g,z)\mapsto f(g,z).
\]

回忆
\[
M_k=\varinjlim_KM_k(K),
\]
其中 $M_k(K)$ 是各种同余子群 $\Gamma_i\subseteq\GL_2(\Z)$ 的模形式空间之和
\[
\bigoplus_{i=1}^mM_k(\Gamma_i),
\]
而 $M_k$ 在 $H^0(\GL_2(\Q)\backslash(\GL_2(\Af)\times\HH^\pm),\omega^{\otimes k})$ 中的像由「在尖点处全纯」这一条件定义。

这样，我们就把模形式与 $\GL_2(\Q)\backslash\GL_2(\A)$ 上的光滑函数联系了起来。注意
\[
[\GL_2]=\GL_2(\Q)A_{\GL_2}\backslash\GL_2(\A),
\]
但我们的函数 $\phi_f$ 不在 $A_{\GL_2}$ 下不变，即
\[
\phi(g,g_\infty r)=r^{-k}\phi(g,g_\infty),\qquad r\in A_{\GL_2}\cong\R_{>0}.
\]
不过这不难修正。令
\[
|\cdot|_{\mathrm{ad\acute{e}lique}}:\A^\times\to\R_{>0},\qquad((x_p)_p,x_\infty)\mapsto\prod_p|x_p|_p|x_\infty|
\]
为\textbf{adélic 绝对值}（它在 $\Q^\times$ 上平凡）。则对 $f\in M_k$，函数
\[
\widetilde\phi_f(g,g_\infty):=|\det(g,g_\infty)|^{k/2}_{\mathrm{ad\acute{e}lique}}\phi_f(g,g_\infty)
\]
在 $\GL_2(\Q)A_{\GL_2}$ 下不变，即定义了光滑函数 $[\GL_2]\to\C$。

% ===== PDF p.022 =====
我们必须澄清模形式何时给出 $L^2([\GL_2])$ 中的函数。为此取截面 $f\in H^0(\GL_2(\Q)\backslash(\GL_2(\Af)\times\HH^\pm),\omega^{\otimes k})$。问题是：
\[
\widetilde\phi_f\in C^\infty([\GL_2])
\]
何时属于 $L^2([\GL_2])$？取紧开子群 $K\subseteq\GL_2(\Af)$ 使
\[
f\in H^0(\GL_2(\Q)\backslash(\GL_2(\Af)/K\times\HH^\pm),\omega^{\otimes k}).
\]
函数 $\widetilde\phi_f$ 是 $L^2$ 的，当且仅当光滑函数 $\widetilde\phi_f\in C^\infty([\GL_2]/K)\subseteq C^\infty([\GL_2])$ 是 $L^2$ 的（因为 $K$ 紧）。这把问题归约到：对同余子群 $\Gamma\subseteq\GL_2(\Z)$ 与 $f\in H^0(\Gamma\backslash\HH^\pm,\omega^{\otimes k})$，函数
\[
\Gamma A_{\GL_2}\backslash\GL_2(\R)\to\C,\qquad
g\mapsto\widetilde\phi_f(g)=|\det(g)|^{k/2}(ci+d)^{-k}f(g(i))
\]
何时平方可积，即属于 $L^2$？我们断言：
\[
|\widetilde\phi_f(g)|^2=|f(g(i))|^2|\operatorname{Im}(g(i))|^k
\]
对任意 $g\in\GL_2(\R)$ 成立，等价地，
\[
|\det(g)|^{1/2}|ci+d|^{-1}=|\operatorname{Im}(g(i))|^{1/2}
\]
对 $g=\begin{pmatrix}a&b\\c&d\end{pmatrix}\in\GL_2(\R)$ 成立。\znoteplain{原文此式写作 $|\det(g)|^{-1/2}|ci+d|=|\operatorname{Im}(g(i))|^{1/2}$，与上一式互为倒数（取 $g=\mathrm{diag}(a,1)$ 即可验）。由 $\operatorname{Im}(g(i))=\det(g)/|ci+d|^2$ 知正确形式为 $|\det(g)|^{1/2}|ci+d|^{-1}=|\operatorname{Im}(g(i))|^{1/2}$，译文按此。}但这确实成立，因为两边都在下列作用下不变：
\begin{itemize}
\item 从右的 $K_\infty=\SO_2(\R)$ 作用，
\item 从左的子群 $\begin{pmatrix}1&*\\0&1\end{pmatrix}\subseteq\GL_2(\R)$ 作用，
\item 从左的 $\GL_2(\R)$ 中对角矩阵的作用。
\end{itemize}
于是
\[
\int_{\Gamma A_{\GL_2}\backslash\GL_2(\R)}|\widetilde\phi_f(g)|^2\,\mathrm{d}g
=\int_{\Gamma\backslash\HH^\pm}|f(z)|^2|y|^k\frac{1}{y^2}\,\mathrm{d}x\wedge\mathrm{d}y
\]
（回忆我们是在 $\HH^\pm$ 的「那个」不变测度上积分）。由于 $\Gamma\backslash\HH^\pm$ 只有有限多个尖点，这引导到尖点处的局部陈述。通过共轭可以归约到考察 $\infty\in\Pone(\Q)$。把 $f$ 在 $\infty$ 处的 Fourier 展开写成
\[
f(z)=\sum_{n\in\Z}a_nq^{n/h},
\]
其中 $h\ge0$（依赖于 $\Gamma$）。则
\[
|y|^kf(z)=2\pi|\log(|q|)|^k\sum_{n\in\Z}a_nq^{n/h},
\]
因为 $y=\log(|q|)$。现在对 $0<r\le1$，
\[
\int_{0<|q|\le r}|\log(|q|)|^{k-2}\Bigl|\sum_{n\in\Z}a_nq^{n/h}\Bigr|^2\frac{i}{2}\,\mathrm{d}q\wedge\mathrm{d}\bar q<\infty,
\]
只要 $n\le0$ 时 $a_n=0$，即只要 $f$ 是尖点的\znote{被积函数在 $q\to0$ 时的行为由常数项 $a_0$ 决定：若 $a_0=0$，则 $\bigl|\sum_{n\ge1}a_nq^{n/h}\bigr|^2=O(|q|^{2/h})$，而 $\int_0^r|\log|q||^{k-2}|q|^{2/h}\,\mathrm{d}|q|$ 收敛；故这里真正用到的是常数项消失，仅「在尖点处全纯」不足以保证平方可积。[来源：原文第 22 页与本讲的整段推导]}。

因此，作为前面讨论的总结，映射
\[
M_k\subseteq H^0(\GL_2(\Q)\backslash(\GL_2(\Af)\times\HH^\pm),\omega^{\otimes k})\to C^\infty([\GL_2]),\qquad f\mapsto\widetilde\phi_f
\]
诱导典范包含
\[
\widetilde\Phi:S_k\hookrightarrow L^2([\GL_2]),
\]
% ===== PDF p.023 =====
其中 $S_k:=\varinjlim_KS_k(K)$。

我们可以作如下观察：
\begin{itemize}
\item $\GL_2(\Af)$ 作用在
\[
\begin{gathered}
H^0(\GL_2(\Q)\backslash(\GL_2(\Af)\times\HH^\pm),\omega^{\otimes k}),\quad
C^\infty(\GL_2(\Q)\backslash\GL_2(\A)),\\
C^\infty([\GL_2]),\quad L^2([\GL_2]),\quad M_k,\quad S_k
\end{gathered}
\]
上。
\item 态射
\[
\Phi:M_k\to C^\infty(\GL_2(\A)),\qquad f\mapsto\phi_f
\]
是 $\GL_2(\A)$-等变的，但
\[
M_k\to C^\infty([\GL_2]),\qquad f\mapsto\widetilde\phi_f=|\det|^{k/2}_{\mathrm{ad\acute{e}lique}}\phi_f
\]
不是（原因在于因子 $|\det|^{k/2}_{\mathrm{ad\acute{e}lique}}$）。
\item 更精确地，对 $g\in\GL_2(\Af)$ 与变动的 $K\subseteq\GL_2(\Af)$，同构
\[
M_k(K)\xrightarrow{\ g\ }M_k(g^{-1}Kg)\subseteq M_k
\]
通过 $M_k=\varinjlim_KM_k(K)$ 定义了作用 $M_k\xrightarrow{\ g\ }M_k$。
\item 对任意固定的 $K\subseteq\GL_2(\Af)$，$\GL_2(\Af)$ 不作用在 $M_k(K)$、$S_k(K)$ 等上\znote{这正是必须取归纳极限的原因：$\GL_2(\Af)$ 的作用只把 $K$ 送到共轭子群 $g^{-1}Kg$，因此单个 $M_k(K)$ 一般不稳定；只有 $\varinjlim_KM_k(K)$ 才承载 $G(\Af)$ 的作用。也正因如此，$M_k$ 不是某一个同余子群的模形式空间，而是「所有级」的模形式。[来源：原文第 22、23 页与本讲的构造动机]}。
\end{itemize}

$\GL_2(\Af)$-作用的存在，是我们从 $\Gamma$-模形式转向 $K$-模形式的主要动机。

于是我们看到接下来的任务：
\begin{itemize}
\item 把 $S_k$ 描述为 $\GL_2(\Af)$-表示。
\item 刻画 $S_k$ 在 $L^2([\GL_2])$ 中的像（或 $M_k$ 在 $C^\infty(\GL_2(\Q)\backslash\GL_2(\A))$ 中的像）。
\end{itemize}

希望模形式与自守表示（在 $L^2$ 意义下）的联系现在更清楚了。对 $f\in S_k$，可以取由 $f$ 的 $\GL_2(\A)$-平移生成的子空间的闭包，并希望它定义一个不可约酉表示。注意这样得到的表示自动落在 $L^2([\GL_2])$ 的离散部分 $L^2_{\mathrm{disc}}([\GL_2])$ 中。

\chapter{从模形式到自守表示 II}

% ===== PDF p.024 =====
\textbf{今天：}
\begin{itemize}
\item 继续构造与尖点形式相伴的自守表示。
\end{itemize}

\textbf{上次：}
\begin{itemize}
\item 我们构造了 $\GL_2(\Af)$-等变嵌入
\[
\Phi:M_k\to C^\infty(\GL_2(\Q)\backslash\GL_2(\A)),
\]
它诱导嵌入
\[
\widetilde\Phi:S_k\to L^2([\GL_2]),
\]
后者在相差因子 $|\det|^{k/2}_{\mathrm{ad\acute{e}lique}}$ 的意义下是 $\GL_2(\Af)$-等变的；也就是说，态射
\[
\widetilde\Phi:S_k\otimes|\det|^{k/2}_{\mathrm{ad\acute{e}lique}}\to L^2([\GL_2])
\]
是 $\GL_2(\Af)$-等变的。
\item 这里
\[
M_k=\varinjlim_{K\subseteq\GL_2(\Af)}M_k(K),
\]
而
\[
M_k(K)=\bigoplus_{i=1}^mM_k(\Gamma_i)
\]
是同余子群 $\Gamma_i\subseteq\GL_2(\Z)$ 的模形式空间之和。
\item 类似地：
\[
S_k=\varinjlim_KS_k(K),
\]
其中 $S_k(K)$ 是同余子群的尖点形式空间之和。
\item 今天：如何刻画像
\[
\Phi(M_k)\subseteq C^\infty(\GL_2(\Q)\backslash\GL_2(\A))
\]
以及
\[
\widetilde\Phi(S_k)\subseteq L^2([\GL_2])？
\]
\item 下次：$S_k$ 如何作为 $\GL_2(\Af)$-表示分解？
\end{itemize}

上一讲中 $\Phi$ 的构造：
\begin{itemize}
\item 我们考虑了 $\GL_2(\R)$-等变线丛
\[
L_k\to\HH^\pm\cong\GL_2(\R)/\C^\times,
\]
它们相伴于 $\C^\times=A_{\GL_2}\SO_2(\R)$ 的表示 $z\mapsto z^k$。
\item 通过拉回，得到 $\GL_2(\Q)\backslash(\GL_2(\Af)\times\HH^\pm)$ 上的线丛 $\omega^{\otimes k}$。
\item $\omega^{\otimes k}$ 的光滑、全纯……截面等同于满足
\[
f(\gamma g,\gamma z)=j(\gamma,z)^kf(g,z)
\]
的光滑、全纯……函数
\[
f:\GL_2(\Af)\times\HH^\pm\to\C,\qquad(g,z)\mapsto f(g,z).
\]
% ===== PDF p.025 =====
\item 这里
\[
\gamma=\begin{pmatrix}a&b\\c&d\end{pmatrix}\in\GL_2(\Q),\qquad j(\gamma,z):=cz+d.
\]
\item 函数
\[
j:\GL_2(\R)\times\HH^\pm\to\C^\times,\qquad
\Bigl(\begin{pmatrix}a&b\\c&d\end{pmatrix},z\Bigr)\mapsto cz+d
\]
满足「余圈方程」
\[
j(gh,z)=j(g,hz)j(h,z)
\]
对一切 $g,h\in\GL_2(\R)$、$z\in\HH^\pm$ 成立。
\item 对如上 $f:\GL_2(\Af)\times\HH^\pm\to\C$，函数
\[
\phi_f:\GL_2(\Af)\times\GL_2(\R)\to\C,\qquad
(g,g_\infty)\mapsto j(g_\infty,i)^{-k}f(g,g_\infty i)
\]
是 $\GL_2(\Q)$-等变的，即属于 $C^\infty(\GL_2(\Q)\backslash\GL_2(\A))$。
\item 由引理 3.2 得到：映射 $f\mapsto\phi_f$ 把光滑截面
\[
\GL_2(\Q)\backslash(\GL_2(\Af)\times\HH^\pm)\to\omega^{\otimes k}
\]
与满足
\[
\phi(g,g_\infty z)=z^{-k}\phi(g,g_\infty)
\]
（对一切 $(g,g_\infty)\in\GL_2(\Af)\times\GL_2(\R)$、$z\in\C^\times$）的光滑函数
\[
\phi:\GL_2(\Q)\backslash\GL_2(\A)\to\C
\]
等同起来。
\end{itemize}

\textbf{要点与下一步目标：}
\begin{itemize}
\item 用 $\phi\in C^\infty(\GL_2(\Q)\backslash\GL_2(\A))$ 上的条件，描述 $M_k$ 在 $f\mapsto\phi_f$ 下的像。
\item 再加上尖点性条件，就决定了 $S_k$ 在
\[
\widetilde\Phi:S_k\to L^2([\GL_2])
\]
下的像：它是那些光滑尖点函数 $\phi$，使得 $|\det|^{-k/2}_{\mathrm{ad\acute{e}lique}}\phi$ 落在 $\Phi$ 的像中。
\item 我们已经知道，$\Phi$ 的像中的元素 $\phi$ 必须满足
\[
\phi(g,g_\infty z)=z^{-k}\phi(g,g_\infty)
\]
（$z\in\C^\times$，$(g,g_\infty)\in\GL_2(\A)$），并且这样的 $\phi$ 下降为满足 $\GL_2(\Q)$-模性的光滑函数
\[
f_\phi:\GL_2(\Af)\times\HH^\pm\to\C,\qquad(g,z)\mapsto f(g,z).
\]
\item 需要检验：
\begin{itemize}
\item $f_\phi$ 何时全纯？
\item $f_\phi$ 何时在尖点处全纯？
\end{itemize}
\end{itemize}

% ===== PDF p.026 =====
\textbf{$f_\phi$ 何时全纯？}这归约到如下问题：给定光滑函数
\[
f:\HH^\pm\to\C,
\]
$\phi_f:\GL_2(\R)\to\C$，$g\mapsto j(g,i)^{-k}f(gi)$ 上的什么条件保证 $f$ 全纯？我们概述将遵循的策略：
\begin{itemize}
\item Lie 代数
\[
\g:=\mathfrak{gl}_2(\R)\cong\mathrm{Mat}_{2,2}(\R)
\]
经右平移对 $\GL_2(\R)$ 的作用求导，而作用在 $C^\infty(\GL_2(\R))$ 上\znote{直观上，$f$ 全纯等价于 $\partial f/\partial\bar z=0$；而 $\partial/\partial\bar z$ 是 $\GL_2(\R)$ 的 Lie 代数作用的一个组合（见下文对 $Y$ 的构造），于是全纯性可以写成 Lie 代数层面的一个方程。这正是把「模形式的全纯性」翻译为「自守表示的无穷小条件」的出发点。[来源：原文第 26--29 页]}。
\item 把这个作用在 $\phi_f$ 上显式化。
\item 然后构造元素 $Y\in\g_\C$，使得
\[
Y*\phi_f=0
\]
当且仅当 $f$ 全纯。这里 $\g_\C:=\g\otimes_\R\C$ 是 $\g$ 的复化 Lie 代数。
\end{itemize}

\textbf{无穷小作用}（参见 \texttt{[GH19, Section 4.2]}）：我们开始解释「对 $\GL_2(\R)$ 的作用求导」是什么意思。
\begin{itemize}
\item 设 $H$ 为任意 Lie 群，Lie 代数 $\mathfrak{h}:=T_1H$。这里 $T_1H$ 表示实流形 $H$ 在单位元处的切空间，即 $\mathfrak{h}=T_1H$ 的元素是道路
\[
\gamma:(-\epsilon,\epsilon)\to H
\]
（$\epsilon>0$，$\gamma(0)=1$）的等价类，两条道路等价若它们在 $1$ 处（在某坐标卡中）有相同的导数。
\item 设 $V$ 是 $H$ 的任意表示，其中 $V$ 是 Hausdorff 拓扑 $\C$-向量空间。
\item 对 $X\in\mathfrak{h}$（由某条道路 $\gamma:(-\epsilon,\epsilon)\to H$，$\gamma(0)=1$ 定义）与 $v\in V$，考虑函数
\[
(-\epsilon,\epsilon)\setminus\{0\}\to V,\qquad t\mapsto\frac{\gamma(t)v-v}{t}.
\]
若该函数在 $t\to0$ 时的极限存在，我们记
\[
X*v:=\lim_{t\to0}\frac{\gamma(t)v-v}{t}.
\]
\item 元素 $v\in V$ 称为\textbf{光滑的}，若
\[
X_1*(\cdots*(X_n*v)\cdots)
\]
对一切 $X_1,\dots,X_n\in\mathfrak{h}$ 都存在。
\item 记 $V^{\mathrm{sm}}\subseteq V$ 为 $V$ 中光滑向量的子空间。
\item 则 $V^{\mathrm{sm}}$ 在 $H$ 下稳定（但可能与 $V$ 不同）。
\item 此外，$V^{\mathrm{sm}}$ 是 $\mathfrak{h}$ 的表示，或等价地是 $\mathfrak{h}$ 的包络代数
\[
U(\mathfrak{h}):=\bigoplus_{n\ge0}\mathfrak{h}^{\otimes n}/\langle X\otimes Y-Y\otimes X-[X,Y]\mid X,Y\in\mathfrak{h}\rangle
\]
上的模。
\item 这里 $[-,-]$ 是 $\mathfrak{h}$ 的 Lie 括号。
\end{itemize}

% ===== PDF p.027 =====
\textbf{$\g=\mathfrak{gl}_2(\R)$ 在 $\phi_f$ 上的显式作用：}现在我们把 $\mathfrak{gl}_2(\R)$ 在函数
\[
\phi_f(g):=j(g,i)^{-k}f(gi)
\]
上的作用写得更显式，其中 $\phi_f$ 相伴于某个截面
\[
f:\HH^\pm\to\OO_{\HH^\pm}(k).
\]
为此回忆等式
\[
j(gh,z)=j(g,hz)j(h,z)
\]
（$g,h\in\GL_2(\R)$，$z\in\HH^\pm$）。
\begin{itemize}
\item 我们已经知道
\[
\phi_f(gz)=z^{-k}\phi_f(g)
\]
（$z\in\C^\times$，$g\in\GL_2(\R)$，由引理 3.2）。注意
\[
\C\cong\Lie(\C^\times)=\Bigl\langle\begin{pmatrix}1&0\\0&1\end{pmatrix},\begin{pmatrix}0&-1\\1&0\end{pmatrix}\Bigr\rangle_\R\subseteq\g.
\]
求导可推出 $\C$ 通过线性形式
\[
\C\to\C,\qquad z\mapsto-kz
\]
作用在 $\phi_f$ 上。因此更精确地，
\[
\begin{pmatrix}1&0\\0&1\end{pmatrix}*\phi_f=-k\phi_f,
\qquad
\begin{pmatrix}0&-1\\1&0\end{pmatrix}*\phi_f=-ki\phi_f,
\]
因为 $\begin{pmatrix}0&-1\\1&0\end{pmatrix}$ 对应 $i\in\C$。
\item 现在考虑子群
\[
U:=\Bigl\{\begin{pmatrix}*&*\\0&1\end{pmatrix}\Bigr\}\subseteq\GL_2(\R)
\]
及其 Lie 代数
\[
\Lie(U)=\Bigl\langle\begin{pmatrix}1&0\\0&0\end{pmatrix},\begin{pmatrix}0&1\\0&0\end{pmatrix}\Bigr\rangle_\R\subseteq\g.
\]
\item 定义辅助函数
\[
\phi_f^\sharp:\GL_2(\R)\times\HH^\pm\to\C,\qquad(g,z)\mapsto j(g,z)^{-k}f(gz)
\]
（不要把它与上一讲引入的 $\widetilde\phi_f$ 混淆）\znote{原文把这里的辅助函数 $\phi_f^\sharp(g,z):=j(g,z)^{-k}f(gz)$ 与上一讲的 $\widetilde\phi_f=|\det|^{k/2}_{\mathrm{ad\acute{e}lique}}\phi_f(x)$ 用了同一个记号 $\widetilde\phi_f$（原文自己也提醒「不要混淆」）。为免混淆，译文把本讲的这个辅助函数记作 $\phi_f^\sharp$，而上一讲的仍记作 $\widetilde\phi_f$。[来源：原文第 27 页；另见第 3 讲第 21 页]}。
\item 对固定的 $g\in\GL_2(\R)$ 得到
\[
\frac{\partial\phi_f^\sharp}{\partial\bar z}(g,z)=\det(g)j(g,z)^{-k-2}\frac{\partial f}{\partial\bar z}(gz)
\]
（用到了函数 $z\mapsto gz$ 的 $\partial/\partial z$-导数为 $\det(g)/j(g,z)^2$）。
\item 从而 $f$ 全纯当且仅当
\[
\frac{\partial\phi_f^\sharp}{\partial\bar z}(g,i)=0
\]
对一切 $g\in\GL_2(\R)$ 成立。
\end{itemize}
% ===== PDF p.028 =====
\begin{itemize}
\item 对任意道路 $\gamma:(-\epsilon,\epsilon)\to U$，$\gamma(0)=1$，可以算出
\begin{align*}
\phi_f(g\gamma(t))
&=\phi_f^\sharp(g\gamma(t),i)\\
&=j(g\gamma(t),i)^{-k}f(g\gamma(t)i)\\
&=j(g,\gamma(t)i)^{-k}j(\gamma(t),i)^{-k}f(g\gamma(t)i)\\
&=\phi_f^\sharp(g,\gamma(t)i),
\end{align*}
因为 $j(\gamma(t),i)=1$。
\item 令
\[
X_\gamma:=\frac{\partial\gamma}{\partial t}(0)\in\Lie(U).
\]
\item 则
\begin{align*}
X_\gamma*\phi_f
&=\frac{\partial}{\partial t}\Bigl(\phi_f(g\gamma(t))\Bigr)\Big|_{t=0}\\
&=\frac{\partial\phi_f^\sharp}{\partial z}(g,i)\cdot\frac{\partial}{\partial t}\bigl(\gamma(t)\cdot i\bigr)\Big|_{t=0}
+\frac{\partial\phi_f^\sharp}{\partial\bar z}(g,i)\cdot\frac{\partial}{\partial t}\bigl(\overline{\gamma(t)\cdot i}\bigr)\Big|_{t=0},
\end{align*}
这里用了 Wirtinger 导数的链式法则\znote{原文此处（以及第 27、29 页）印刷为 $\partial/\partial z$，但按链式法则第二项应为 $\partial/\partial\bar z$，且后文「$Y*\phi_f=2i\partial\phi_f^\sharp/\partial\bar z$」与「$f$ 全纯 $\iff$ $Y*\phi_f=0$」只有在 $\partial/\partial\bar z$ 时才自洽（提取文本系统性地丢了 $\bar z$ 上的横线，参见第 13 页 $dx\wedge dy=\frac i2\,dz\wedge d\bar z$）。译文按 $\partial/\partial\bar z$ 处理。[来源：原文第 27--29 页]}。
\item 更具体地，先取第一条道路
\[
\gamma_1:(-\epsilon,\epsilon)\to U,\qquad t\mapsto\begin{pmatrix}1+t&0\\0&1\end{pmatrix}.
\]
则 $\gamma_1(t)\cdot i=(1+t)i$，于是
\[
\begin{pmatrix}1&0\\0&0\end{pmatrix}*\phi_f(g)
=i\frac{\partial\phi_f^\sharp}{\partial z}(g,i)-i\frac{\partial\phi_f^\sharp}{\partial\bar z}(g,i).
\]
\item 再取第二条道路
\[
\gamma_2:(-\epsilon,\epsilon)\to U,\qquad t\mapsto\begin{pmatrix}1&t\\0&1\end{pmatrix},
\]
则 $\gamma_2(t)\cdot i=i+t$\znote{原文此处误作 $\gamma_1(t)\cdot i=i+t$；由定义 $\gamma_1(t)\cdot i=(1+t)i$，而 $\gamma_2(t)\cdot i=i+t$，故应为 $\gamma_2$。[来源：原文第 28 页]}，于是
\[
\begin{pmatrix}0&1\\0&0\end{pmatrix}*\phi_f(g)
=\frac{\partial\phi_f^\sharp}{\partial z}(g,i)+\frac{\partial\phi_f^\sharp}{\partial\bar z}(g,i).
\]
\item 特别地，
\[
i\begin{pmatrix}1&0\\0&0\end{pmatrix}*\phi_f(g)+\begin{pmatrix}0&1\\0&0\end{pmatrix}*\phi_f(g)
=2\frac{\partial\phi_f^\sharp}{\partial\bar z}(g,i).
\]
\item 因此
\[
\begin{pmatrix}i&1\\0&0\end{pmatrix}\in\g_\C
\]
可以作为 $Y$ 的候选，即 $f$ 全纯当且仅当
\[
\begin{pmatrix}i&1\\0&0\end{pmatrix}*\phi_f=0.
\]
\item 但我们还可以做得更好。也就是说，定义
\[
Y:=\frac12\begin{pmatrix}-1&i\\i&1\end{pmatrix}
=\frac12\begin{pmatrix}1&0\\0&1\end{pmatrix}+\frac{1}{2i}\begin{pmatrix}0&-1\\1&0\end{pmatrix}
+i\begin{pmatrix}i&1\\0&0\end{pmatrix}.
\]
\znote{原文这个分解式与左边的矩阵不符：按右边逐项相加得到的是 $\frac12\begin{pmatrix}-1&i\\-i&1\end{pmatrix}$，即 $(2,1)$-项的符号相反（把 $+\frac{1}{2i}$ 改为 $-\frac{1}{2i}$ 才与左边的矩阵一致；下页对 $Y*\phi_f$ 的计算实际用的也是与左边矩阵一致的那个元素）。译文照左边的矩阵与原文的分解式直译，仅在此说明该分解式有误，不另作改动。[来源：原文第 28--29 页；参见 work/疑点清单.tsv 第 28 页条目]}
\end{itemize}

% ===== PDF p.029 =====
于是
\[
Y*\phi_f(g)
=\underbrace{-\frac k2\phi_f(g)+i\Bigl(-\frac{ki}{2}\Bigr)\phi_f(g)}_{=0}
+2i\frac{\partial\phi_f^\sharp}{\partial\bar z}(g,i)
=2i\frac{\partial\phi_f^\sharp}{\partial\bar z}(g,i),
\]
从而 $f$ 全纯当且仅当 $Y*\phi_f=0$。
\begin{itemize}
\item 为什么这个选择更好？
\begin{itemize}
\item 令
\[
H:=\begin{pmatrix}0&-i\\i&0\end{pmatrix},\qquad
X:=\frac12\begin{pmatrix}-1&-i\\-i&1\end{pmatrix}.
\]
\item 则 $H,X,Y$ 是 $\g_\C$ 中的一个 $\mathfrak{sl}_2$-三元组。
\item 确切地说，在基
\[
\begin{pmatrix}-i\\1\end{pmatrix},\qquad\begin{pmatrix}i\\1\end{pmatrix}
\]
下，它由
\[
\begin{pmatrix}1&0\\0&-1\end{pmatrix},\qquad
\begin{pmatrix}0&1\\0&0\end{pmatrix},\qquad
\begin{pmatrix}0&0\\1&0\end{pmatrix}
\]
给出。
\item 注意 $H$ 张成 Lie 代数 $i\Lie(\SO_2(\R))\subseteq\g_\C$。
\end{itemize}
\end{itemize}

我们可以转向下一个问题了。

\textbf{$f_\phi$ 何时在尖点处全纯？}
\begin{itemize}
\item 考虑具有 Fourier 展开的函数
\[
f(z)=\sum_{n\in\Z}a_ne^{2\pi izn},
\]
它定义在
\[
\{z\in\C\mid|\Re(z)|<1,\ \Im(z)>1\}
\]
上。
\item 回忆
\[
|e^{2\pi iz}|=e^{-2\pi y},
\]
其中 $y=\Im(z)$。
\item 则 $a_n=0$ 对一切 $n<0$ 成立，当且仅当
\[
|f(x+iy)|\le Cy^N
\]
对某 $C,N\in\R_{>0}$ 与 $y\to\infty$ 成立，因为当 $n<0$ 时指数函数 $e^{-2\pi ny}$ 的增长快于任何多项式函数。
\item 按定义，这就是说 $f$ 具有\textbf{适度增长}（或缓慢增长）；这一概念可以推广到 $\Spec(\Q)$ 上的一切约化群。
\item 我们还要指出：进一步地，$a_n=0$ 对一切 $n\le0$（即包括 $n=0$）成立，当且仅当对一切 $N\in\R_{>0}$，
\[
|f(x+iy)|\le Cy^{-N}
\]
对某 $C\in\R_{>0}$ 与 $y\to\infty$ 成立。
\item 换言之，$f$ 是\textbf{速降}的\znote{「适度增长」与「速降」只差常数项 $a_0$ 是否为零：前者允许 $y$ 的多项式增长，后者要求对任意 $N$ 都有 $y^{-N}$ 级别的衰减。这正是第 3 讲中「尖点处全纯」与「尖点形式」的区别在分析层面的体现：去掉 $n\le0$ 的 Fourier 系数就得到速降。[来源：原文第 29 页；参见第 3 讲第 14、22 页]}。
\end{itemize}

\textbf{一般约化群的适度增长：}对 $\Q$ 上任意的约化群 $G$，适度增长条件可以这样表述。
\begin{itemize}
\item 对每个嵌入 $\rho:G\hookrightarrow\SL_{n,\Q}$（不是到 $\GL_{n,\Q}$），我们得到 $G(\A)$ 上的范数
\[
|g|:=\sup_v\max_{i,j=1,\dots,n}\{|\rho(g)_{i,j}|_v\},
\]
其中 $v$ 遍历一切素数与 $\infty$。
% ===== PDF p.030 =====
\item 函数 $\phi:G(\A)\to\C$ 称为\textbf{适度增长}（或缓慢增长）的，若
\[
|\phi(g)|\le C|g|^N
\]
对适当的 $C,N\in\R_{>0}$ 成立，参见 \texttt{[GH19, Definition 6.4.]}。
\item 不难看出，这个定义不依赖嵌入 $\rho$（因为换一个 $\rho$ 只会在相差某个多项式函数的意义下改变范数）。
\end{itemize}

我们把它留作练习（参见 \texttt{[GH19, Lemma 6.3.1.]} 以及上面关于 Fourier 系数的注记）：检验一个（全纯）截面
\[
f\in H^0(\GL_2(\Q)\backslash(\GL_2(\Af)\times\HH^\pm),\omega^{\otimes k})
\]
在尖点处全纯，当且仅当函数 $\phi_f\in C^\infty(\GL_2(\Q)\backslash\GL_2(\A))$ 具有适度增长。

我们简短总结一下已完成的工作。

\textbf{要点：}
\begin{itemize}
\item 我们可以描述
\[
M_k\to C^\infty(\GL_2(\Q)\backslash\GL_2(\A)),\qquad f\mapsto\phi_f
\]
的像。
\item 也就是说，$\phi\in C^\infty(\GL_2(\Q)\backslash\GL_2(\A))$ 落在像中，当且仅当
\begin{itemize}
\item $\phi(g,g_\infty z)=z^{-k}\phi(g,g_\infty)$ 对 $(g,g_\infty)\in\GL_2(\A)$、$z\in\C^\times$ 成立。
\item $\phi$ 具有适度增长。
\item 对每个 $g\in\GL_2(\Af)$，
\[
Y*\phi(g,-)=0,
\]
其中 $Y\in\g_\C=\mathfrak{gl}_2(\R)_\C$ 是前面构造的元素。
\end{itemize}
\item 注意，按我们的讨论，这些条件恰好对应模形式定义中的「模性、在尖点处全纯、全纯」。
\item 对 $k\ge1$，不难看出存在唯一的（在同构意义下）不可约 $(\g_\C,O_2(\R))$-模 $D'_{k-1}$，它含非零元素 $v$，使得 $Y*v=0$，且 $\C^\times$ 通过特征 $z\mapsto z^{-k}$ 作用在 $v$ 上，参见 \texttt{[Del73, Section 2.1]}（其中描述了 $D'_{k-1}$ 的对偶）\znote{$D'_{k-1}$ 是 $(\g_\C,O_2(\R))$-模：$O_2(\R)$ 给出 $K_\infty$-类型，$\g_\C$ 给出 Lie 代数作用；$Y*v=0$ 即最低权向量条件，$z\mapsto z^{-k}$ 则固定了权。于是上面的要点把「权 $k$ 的模形式」识别为从一个极小的表示 $D'_{k-1}$ 到适度增长函数空间的 $(\g_\C,O_2(\R))$-模同态，这正是把模形式纳入自守表示框架的关键一步；$D'_{k-1}$ 的构造在原文第 11 讲讨论。[来源：原文第 30 页]}。
\item 利用 $D'_{k-1}$，我们的结果可以改写为
\[
M_k\cong\Hom_{(\g_\C,O_2(\R))}\bigl(D'_{k-1},C^\infty_{\mathrm{mg}}(\GL_2(\Q)\backslash\GL_2(\A))\bigr),
\]
其中 $C^\infty_{\mathrm{mg}}(\GL_2(\Q)\backslash\GL_2(\A))\subseteq C^\infty(\GL_2(\Q)\backslash\GL_2(\A))$ 表示具有适度增长的函数构成的子集。
\item 更多内容将在我们讨论自守形式与 $(\g_\C,K_\infty)$-模时出现，参见第 11 讲。
\end{itemize}

我们以概述如何把尖点性纳入进来来结束本讲。

\textbf{尖点性的一个条件：}
\begin{itemize}
\item 设 $\phi=\phi_f\in C^\infty(\GL_2(\Q)\backslash\GL_2(\A))$，其中 $f\in M_k$，并把 $\phi$ 视为满足模性的函数 $\GL_2(\Af)\times\HH^\pm\to\C$。
\item 则 $f\in S_k$ 当且仅当积分
\[
\phi_B(g):=\int_{N(\Q)\backslash N(\A)}\phi(ng)\,dn
\]
% ===== PDF p.031 =====
对一切 $g\in\GL_2(\A)$ 为零；这里
\[
N=\begin{pmatrix}1&*\\0&1\end{pmatrix}
\]
是标准 Borel 子群
\[
B=\begin{pmatrix}*&*\\0&*\end{pmatrix}
\]
中的幺幂根基，$dn$ 是 $N(\Q)\backslash N(\A)$ 上的 $N(\A)$-不变测度\znote{把 $\phi$ 沿幺幂根基 $N$ 积分再令其为零，就是经典意义下「在所有尖点处常数 Fourier 系数为零」的表示论写法；下文的计算表明 $\phi_B(1)=\mu(m\hat\Z)a_0$，于是 $\phi_B\equiv0$ 恰好等价于各尖点处 $a_0=0$，即 $f$ 是尖点形式。[来源：原文第 30--31 页]}。
\item 也就是说：
\begin{itemize}
\item 可以假设 $f$（从而 $\phi=\phi_f$）在某个主同余子群 $K(m)\subseteq\GL_2(\Af)$ 下不变。
\item 取 $g=1$。利用 $N(\Q)\backslash N(\A)\cong(m\hat\Z\times\R)/m\Z$：
\begin{align*}
\int_{N(\Q)\backslash N(\A)}\phi_f(n)\,dn
&=\int_{N(\Q)\backslash N(\A)}j(n_\infty,i)^{-k}f(n,n_\infty i)\,d(n,n_\infty)\\
&=\int_{\R/m\Z}\int_{m\hat\Z}f(a,i+x)\,da\,dx\\
&=\mu(m\hat\Z)\int_{\R/m\Z}f(1,i+x)\,dx\\
&=\mu(m\hat\Z)a_0,
\end{align*}
其中 $\mu(m\hat\Z)$ 是 $m\hat\Z$ 的体积，而
\[
f(1,z)=\sum_{n=0}^\infty a_nq^{n/m}
\]
是模形式 $f(1,-)$ 关于 $\Gamma(m)$ 在尖点 $\infty$ 处的 Fourier 展开。
\item 确实：关键在于
\[
\int_{\R/m\Z}e^{2\pi i\frac nmx}\,dx=0
\]
当 $n>0$ 时成立（这是 Fourier 理论的一部分；更一般地，它可由「对任意紧 abelian 群 $A$ 与任意非平凡特征 $\chi:A\to\C^\times$，有 $\int_A\chi(g)\,dg=0$」推出）。
\item 于是条件「对一切 $g\in\GL_2(\A)$ 有 $\phi_B(g)=0$」等价于所有尖点处的常数 Fourier 系数消失。
\item 细节留作练习，参见 \texttt{[Del73, Rappel 1.3.2.]} 与 \texttt{[Gel75, Proposition 3.1.(vii)]}。
\end{itemize}
\item 注：因为 $\phi$ 对左平移是 $\GL_2(\Q)$-不变的，我们可以把 $B$ 换成任意共轭 $gBg^{-1}$，$g\in\GL_2(\Q)$。
\end{itemize}

我们要指出，我们把函数与 $\GL_2(\A)$ 相联系的方式在表述上不同于标准文献，例如 \texttt{[Gel75]}，后者通常限制到
\[
\Gamma_1(m):=\Bigl\{\begin{pmatrix}a&b\\c&d\end{pmatrix}\in\SL_2(\Z)\ \Bigm|\ c\equiv0,\ d\equiv1\ \mathrm{mod}\ m\Bigr\}
\]
% ===== PDF p.032 =====
上的模形式，或限制到关于
\[
\Gamma_0(m):=\Bigl\{\begin{pmatrix}a&b\\c&d\end{pmatrix}\in\SL_2(\Z)\ \Bigm|\ c\equiv0\ \mathrm{mod}\ m\Bigr\}
\]
带旁特征\idx{旁特征} $\chi:(\Z/m)^\times\to\C^\times$ 的模形式。尽管这些抓住了由模形式生成的一切自守表示，我们仍觉得在此刻把自己限制到这些同余子群（或它们在 $\GL_2(\Af)$ 中的对应物 $K_1(m),K_0(m)$）有点缺乏动机。下一讲我们将澄清我们的方法与经典方法之间的关系。

\chapter{从模形式到自守表示 III}

% ===== PDF p.033 =====
\textbf{上次}：
\begin{itemize}
\item 设 $\phi\in C^\infty(\GL_2(\Q)\backslash\GL_2(\A))$。
\item 则 $\phi=\phi_f$（某个 $f\in M_k$）当且仅当
\begin{itemize}
\item $\phi(g,g_\infty z)=z^{-k}\phi(g,g_\infty)$，对 $(g,g_\infty)\in\GL_2(\A)$、$z\in\C^\times$；
\item $\phi$ 适度增长；
\item 对每个 $g\in\GL_2(\Af)$ 有 $Y*\phi(g,-)=0$，其中 $Y\in\g_\C=\mathfrak{gl}_2(\R)_\C$ 是适当构造的自然元素。
\end{itemize}
\item $f\in M_k$ 是尖点形式，当且仅当对每个真抛物子群 $B\subseteq\GL_{2,\Q}$（即 Borel 子群），以 $N\subseteq\GL_{2,\Q}$ 记其幺幂根基，有
\[
\int_{N(\Q)\backslash N(\A)}\phi_f(ng)\,\mathrm{d}n=0
\]
对一切 $g\in\GL_2(\A)$ 成立。
\end{itemize}

\textbf{今天}：
\begin{itemize}
\item 把 $S_k$ 描述为 $\GL_2(\Af)$-表示。
\item 由此得到自守表示的主要例子来源。\znoteplain{原文脚注：更精确地说，$\widetilde\Phi(S_k)\subseteq L^2([\GL_2])$ 的闭包的不可约子表示，正是与尖点形式相伴的（$L^2$ 意义下的）自守表示。}
\end{itemize}

考察整个空间 $M_k$、$S_k$（而不只是某个紧开子群 $K\subseteq\GL_2(\Af)$ 对应的 $M_k(K)$、$S_k(K)$，也不只是同余子群 $\Gamma\subseteq\SL_2(\Z)$ 对应的 $M_k(\Gamma)$、$S_k(\Gamma)$）的好处是：现在可以使用局部 profinite 群 $\GL_2(\Af)$ 的表示论。就此而言，下面的观察很重要。

\textbf{观察}：
\begin{itemize}
\item $M_k$ 与 $S_k$ 是光滑且容许的 $\GL_2(\Af)$-表示（在下述意义下）。
\end{itemize}

\noindent\textbf{定义 5.1}\phantomsection\label{def:5.1}\quad 设 $G$ 为局部 profinite 群（例如 $\GL_2(\Af)$）。$G$ 在 $\C$-向量空间 $V$ 上的表示称为\textbf{光滑的}（smooth），若
\[
V=\bigcup_{K\subseteq G}V^K,
\]
其中 $K$ 遍历 $G$ 的紧开子群，且
\[
V^K:=\{v\in V\mid kv=v\ \text{对一切}\ k\in K\}.
\]
等价地，$V$ 光滑当且仅当作用态射 $G\times V\to V$ 连续，其中 $V$ 取离散拓扑。记 $\mathrm{Rep}^\infty_\C G$ 为 $G$ 在 $\C$-向量空间上的光滑表示构成的（abelian）范畴。

% ===== PDF p.034 =====
\noindent\textbf{定义 5.2}\phantomsection\label{def:5.2}\quad 设 $G$ 为局部 profinite 群。$G$ 的光滑表示 $V$ 称为\textbf{容许的}（admissible），若
\[
\dim_\C V^K<\infty
\]
对一切紧开子群 $K\subseteq G$ 成立。\znoteplain{光滑性与容许性是两件不同的事：光滑只要求每个向量都被某个紧开子群固定，容许性进一步要求每个 $V^K$ 有限维。对 $\GL_2(\Af)$ 而言 $M_k(K)$ 有限维，所以 $M_k$、$S_k$ 两者兼备；$C^\infty(G(\Q)\backslash G(\A))$ 光滑但不容许；$L^2([G])$ 连光滑都不是。}

看几个例子。
\begin{itemize}
\item $M_k=\bigcup_{K\subseteq\GL_2(\Af)}M_k^K$，其中 $M_k^K=M_k(K)$，它是光滑且容许的，因为各 $M_k(K)$ 都有限维。
\item 同理，$S_k$ 是 $\GL_2(\Af)$ 的容许表示。
\item 若 $G$ 在 $\Q$ 上约化，则 $C^\infty(G(\Q)\backslash G(\A))$ 是光滑的 $G(\Af)$-表示，但不容许。
\item $L^2([G])$ 不是光滑的 $G(\Af)$-表示。
\end{itemize}

在第 11 讲中，我们将对 $\Q$ 上的约化群 $G$ 引入自守形式空间。它可作为 $M_k$ 的替代（实际上是极大的推广），并且仍是光滑的 $G(\Af)$-表示。我们将会看到，光滑表示比 $L^2([G])$ 这样的酉表示在性质上更代数，这使得局部 profinite 群的光滑表示更便于操作。

在模形式的经典理论中，Hecke 算子（其定义看起来可能有些 ad hoc）很重要。从表示论的观点看，它们来自一个 Hecke 代数\idx{Hecke 代数}的作用。

\textbf{Hecke 代数}（局部 profinite 群的）：
\begin{itemize}
\item $G$ 为局部 profinite 群。
\begin{itemize}
\item 例如 $\GL_2(\Af)$、$\GL_n(\Q_p)$、$\Gal(\overline\Q/\Q)$、$\dots$
\end{itemize}
\item 对紧开子群 $K\subseteq G$，存在唯一的左不变（相应地，唯一的右不变）Haar 测度 $\mu$，使 $\mu(K)=1$。
\item 具体地：
\begin{itemize}
\item 由平移不变性，只需对紧开子群 $K'\subseteq G$ 定义 $\mu$。
\item 令 $K'':=K'\cap K$，它紧开。则 $K''$ 在 $K$ 中、相应地也在 $K'$ 中指数有限。
\item 令
\[
\mu(K'):=\frac{[K':K'']}{[K:K'']}.
\]
\item 这定义了所要的测度，且这是唯一可能的选择。
\item 注意 $\mu$ 的取值实际上在 $\Q$ 中。这是「局部 profinite 群的光滑表示论更代数」的第一个迹象。
\end{itemize}
\item 今后固定 $G$ 上一个左不变 Haar 测度 $\mu$。
\item 定义相应的「Hecke 代数」
\[
H(G):=C_c^\infty(G),
\]
即局部常值、紧支集函数 $G\to\C$ 之集。
\item $H(G)$ 通过卷积成为代数：
\[
f*g(x):=\int_G f(xy)g(y^{-1})\,\mathrm{d}y
\]
% ===== PDF p.035 =====
对 $f,g\in H(G)$。注意 $H(G)$ 隐含依赖于所选测度。不过可以把 $H(G)$ 内蕴地定义为紧支分布之空间，乘法由分布的卷积给出，参见 \texttt{[BRa, Definition 7]}。
\item $H(G)$ 中的乘法是纯代数运算：
\begin{itemize}
\item 每个 $f\in H(G)$ 都是有限和
\[
f=\sum_{i=1}^m a_i\chi_{g_iK_i},
\]
其中 $g_i\in G$，$K_i\subseteq G$ 为紧开子群。
\item 不妨设 $K_1=\dots=K_m$（通过缩小）。
\item 设 $K\subseteq G$ 紧开，$g_1,g_2\in G$。则存在 $h_j\in G$ 使
\[
g_1Kg_2K=\bigsqcup_{j=1}^m h_jK,\qquad
\chi_{g_1K}*\chi_{g_2K}=\sum_{j=1}^m\mu(K)\chi_{h_jK}.
\]
\end{itemize}
\item 代数 $H(G)$ 有很多幂等元，即
\[
e_K:=\frac{1}{\mu(K)}\chi_K,
\]
其中 $K\subseteq G$ 遍历紧开子群，$\chi_K$ 为 $K$ 的特征函数。但一般 $H(G)$ 没有单位元（除非 $G$ 紧）。
\end{itemize}

\textbf{Hecke 代数的模}（参见 \texttt{[BRa]}）：
\begin{itemize}
\item $H(G)$ 是群代数 $\C[G]$ 的替代品，它更适合用来研究 $G$ 的光滑表示。
\item 若 $V\in\mathrm{Rep}^\infty_\C G$，则
\[
f*v:=\int_G f(h)hv\,\mathrm{d}h
\]
（$f\in H(G)$，$v\in V$）定义了 $H(G)$ 在 $V$ 上的作用。
\item 具体地：若 $K\subseteq G$ 是紧开子群且充分小，使 $v\in V$ 被 $K$ 固定，而 $f=\sum_{i=1}^m a_i\chi_{g_iK}$（$a_i\in\C$，$g_i\in G$），则
\[
f*v=\sum_{i=1}^m a_ig_iv.
\]
\item 设 $K\subseteq G$ 为紧开子群。定义 $G$ 相对于 $K$ 的 Hecke 代数
\[
H(G,K):=e_KH(G)e_K.
\]
则 $H(G,K)$ 是以 $e_K$ 为单位的代数。
\item 设 $V$ 为 $G$ 的任意光滑表示。则不难看出
\[
V^K=e_K*V
\]
对任意紧开子群 $K\subseteq G$ 成立。
% ===== PDF p.036 =====
\item 把它应用于光滑 $G$-表示 $H(G)$（经左、右乘法），得到
\[
H(G,K)\cong C_c^\infty(K\backslash G/K),
\]
即 $H(G,K)$ 等同于 $G$ 上 $K$-双不变、紧支集函数之集。
\end{itemize}

下面这个小命题在随后讨论 $\GL_2(\Af)$-表示 $S_k$ 时很重要。
\begin{itemize}
\item 若 $V\in\mathrm{Rep}^\infty_\C G$ 不可约且 $V^K\neq0$，则 $V^K$ 作为 $H(G,K)$-模不可约。
\begin{itemize}
\item 事实上：若 $M\subseteq V^K$ 是非零 $H(G,K)$-子模，则由 $V$ 的不可约性有 $V=H(G)*M$，从而
\[
V^K=e_K(H(G)*M)=e_KH(G)e_KM=M,
\]
因为 $e_K*M=M$。
\end{itemize}
\item 反之可以证明：若对一切紧开子群 $K\subseteq G$，$V^K$ 不可约或为零，则 $V$ 不可约（当 $V\neq0$ 时）。
\end{itemize}

也就是说，这个性质使我们能把 Hecke 特征系统\idx{Hecke 特征系统}联系到不可约容许的 $\GL_2(\Af)$-表示上。

\textbf{从不可约容许 $\GL_2(\Af)$-表示到 Hecke 特征系统}：
\begin{itemize}
\item 设 $V$ 为不可约、容许的 $\GL_2(\Af)$-表示。
\item 回忆紧开子群
\[
K(m)=\ker(\GL_2(\widehat\Z)\to\GL_2(\Z/m)),\qquad m\ge0,
\]
构成 $\GL_2(\Af)$ 中单位元的紧开邻域基。
\item 由上述重要性质与容许性（后者容易推出 Schur 引理的一个版本）知
\begin{itemize}
\item 当 $m\gg0$ 时，$V^{K(m)}$ 是不可约 $H(\GL_2(\Af),K(m))$-模；
\item $\End_{H(\GL_2(\Af),K(m))}(V^{K(m)})\cong\C$。
\end{itemize}
\item 下面我们用到：对每个素数 $p$，光滑 $\GL_2(\Af)$-表示沿包含 $\GL_2(\Q_p)\subseteq\GL_2(\Af)$ 限制，得到光滑的 $\GL_2(\Q_p)$-表示。
\item 令
\[
S:=\bigcup_{m\ge0,\,V^{K(m)}\neq0}\{\,p\ \text{素数}:p\mid m\,\}.
\]
\item 注意 $p\in S$ 当且仅当 $V^{\GL_2(\Z_p)}=0$。
\item 若 $p\notin S$，即 $V^{\GL_2(\Z_p)}\neq0$，则称素数 $p$ 对 $V$ \textbf{非分歧\idx{非分歧}}。
\item 下面的观察非常重要：对每个素数 $p$，代数 $H(\GL_2(\Q_p),\GL_2(\Z_p))$ 交换！
\item 更强的是：取 $m\ge1$ 与 $p\nmid m$。则 $H(\GL_2(\Q_p),\GL_2(\Z_p))$ 在 $H(\GL_2(\Af),K(m))$ 中中心。更精确地说，对某个 $A\in\GL_2(\Q_p)$ 的双陪集，可以联系到元素
\[
(1,\dots,1,A,1,\dots,1)\in\GL_2(\Af)
\]
（第 $p$ 个分量为 $A$，其余分量为单位元）的双陪集，而所得的嵌入
\[
H(\GL_2(\Q_p),\GL_2(\Z_p))\longrightarrow H(\GL_2(\Af),K(m))
\]
% ===== PDF p.037 =====
其像落在 $H(\GL_2(\Af),K(m))$ 的中心。因为容易看出：由 $(1,\dots,1,A,1,\dots,1)$（$A\in\GL_2(\Q_p)$）给出的双陪集与 $B\in\GL_2(\Af)$ 的双陪集交换，只要 $B$ 在 $p$ 处的分量是单位矩阵。于是 $H(\GL_2(\Q_p),\GL_2(\Z_p))$ 的中心性实际上由它的交换性推出。
\item 为证交换性，我们用 Gelfand 技巧。即考虑反对合
\[
\sigma:\GL_2(\Q_p)\to\GL_2(\Q_p),\qquad g\mapsto g^{tr},
\]
它具有以下性质：
\begin{itemize}
\item $\sigma$ 保持 $\GL_2(\widehat\Z_p)$；
\item $\sigma(f_1*f_2)=\sigma(f_2)*\sigma(f_1)$，对 $f_1,f_2\in H(\GL_2(\Q_p),\GL_2(\Z_p))$；
\item 双陪集 $\GL_2(\Z_p)g\GL_2(\Z_p)$（$g\in\GL_2(\Q_p)$ 为对角矩阵）张成 $H(\GL_2(\Q_p),\GL_2(\Z_p))$（由初等因子定理）；
\item 现在可以证明交换性了。因为 $\sigma(f)=f$ 对一切 $f\in H(\GL_2(\Q_p),\GL_2(\Z_p))$ 成立，故
\[
f_1*f_2=\sigma(f_1*f_2)=\sigma(f_2)*\sigma(f_1)=f_2*f_1,
\]
如所欲。
\end{itemize}
\item $H(\GL_2(\Q_p),\GL_2(\Z_p))$ 的中心性很重要，因为它允许定义 $V$ 的 Hecke 特征系统（回忆 $V$ 是不可约容许的 $\GL_2(\Af)$-表示）。
\item 回忆 $S$ 记 $V$ 的分歧素数集，即那些使 $V^{\GL_2(\Z_p)}=0$ 的素数 $p$。
\item 对 $p\notin S$，我们下面构造一个 $\C$-代数态射
\[
\sigma_{V,p}:H(\GL_2(\Q_p),\GL_2(\Z_p))\to\C,
\]
它自然联系于 $V$。
\item 为此取 $p\notin S$，并选 $v\in V$ 被 $\GL_2(\Z_p)$ 固定。则由 $\GL_2(\Af)$-作用的光滑性，$v$ 被某个 $K(m)$ 固定，其中 $m\ge0$ 且 $p\nmid m$。
\item 如前所述，$V$ 的不可约性蕴涵 $V^{K(m)}$ 不可约，从而由 Schur 引理
\[
\End_{H(\GL_2(\Af),K(m))}(V^{K(m)})\cong\C.
\]
\item 由于 $H(\GL_2(\Q_p),\GL_2(\Z_p))$ 在 $H(\GL_2(\Af),K(m))$ 中中心，它通过 $H(\GL_2(\Af),K(m))$-模同态作用在 $V^{K(m)}$ 上。
\item 于是得到态射
\[
\sigma_{V,p}:H(\GL_2(\Q_p),\GL_2(\Z_p))\to\End_{H(\GL_2(\Af),K(m))}(V^{K(m)})\cong\C,
\]
它先验依赖于 $v$ 与 $m$。但不难看出 $\sigma_{V,p}$ 在放大 $m$ 时不变。这意味着 $\sigma_{V,p}$ 实际上与 $v$、$m$ 无关，从而典范地联系于 $V$ 与 $p$。
\end{itemize}

\noindent\textbf{定义 5.3}\phantomsection\label{def:5.3}\quad 系统
\[
\{\sigma_{V,p}:H(\GL_2(\Q_p),\GL_2(\Z_p))\to\C\}_{p\notin S}
\]
称为 $V$ 的\textbf{Hecke 特征系统}（system of Hecke eigenvalues）。

% ===== PDF p.038 =====
我们给出 Hecke 特征系统的另一种视角。设 $S$ 为某个有限素数集，定义「远离 $S$ 的 Hecke 代数」为
\[
T_S:=H(\GL_2(\Af^S),\GL_2(\widehat\Z^S)),
\]
其中
\[
\widehat\Z^S:=\prod_{p\notin S}\Z_p
\qquad\text{且}\qquad
\Af^S:=\prod_{p\notin S}\Q_p.
\]
\znote{原文把 $\Af^S$ 写成 $\Q\otimes_\Z\widehat\Z^S$；受限张量积的正确刻画是 $\Af^S=\prod_{p\notin S}\Q_p$（即几乎处处落在 $\Z_p$ 中的族），张量积一般不保持无穷乘积，故此处按正解译。[来源：原文 p.038 紧随其后的 $\widehat\Z^S$ 定义与 $T_S$ 的受限张量积式，参见 \texttt{[Fla79]}。]}

利用 Gelfand 技巧可以看出 $T_S$ 交换。此外，若展开适当的定义（参见 \texttt{[Fla79]}），则「受限张量积」给出
\[
T_S\cong\sideset{}{'}\bigotimes_{p\notin S}H(\GL_2(\Q_p),\GL_2(\Z_p)).
\]

现在还可以作如下定义。

\noindent\textbf{定义 5.4}\phantomsection\label{def:5.4}\quad \textbf{Hecke 特征系统}（Hecke eigensystem，或 system of Hecke eigenvalues）是 $T_S$ 中的一个极大理想，其中 $S$ 为有限素数集，$T_S:=H(\GL_2(\Af^S),\GL_2(\widehat\Z^S))$。

我们可以收集以下观察。
\begin{itemize}
\item $T_S$ 的维数可数，因此每个 Hecke 特征系统的剩余域为 $\C$。
\item $S$ 的一个 Hecke 特征系统等价于一族 $\C$-代数同态
\[
\sigma_p:H(\GL_2(\Q_p),\GL_2(\Z_p))\to\C\qquad(p\notin S),
\]
如前。
\item 若 $S\subseteq S'$ 为两个有限素数集，则每个 $S$ 的 Hecke 特征系统诱导 $S'$ 的一个。两个 Hecke 特征系统 $\{\sigma_p\}_{p\notin S}$ 与 $\{\sigma'_p\}_{p\notin S'}$ 称为\textbf{等价的}，若它们在某个包含 $S$ 与 $S'$ 的有限素数集之外一致。
\item 如前所见，每个不可约容许 $\GL_2(\Af)$-表示都通过上述构造给出一个 Hecke 特征系统。
\end{itemize}

我们将会看到，在 Langlands 纲领中，与模形式相伴的 Hecke 特征系统被期望具有算术意义。

为了对 Hecke 特征系统获得更有用的视角，我们需要 $H(\GL_2(\Q_p),\GL_2(\Z_p))$ 的完整描述。

\textbf{$H(\GL_2(\Q_p),\GL_2(\Z_p))$ 的完整描述}：
\begin{itemize}
\item 暂令 $G:=\GL_2(\Q_p)$，$K=\GL_2(\Z_p)$。
\item 归一化 $G$ 上的 Haar 测度，使 $\mu(K)=1$。
\item 由初等因子定理，$\C$-向量空间 $H(G,K)$ 以
\[
K\begin{pmatrix}p^i&0\\0&p^j\end{pmatrix}K,\qquad
(i,j)\in\Z^2_+:=\{(i,j)\in\Z^2\mid i\ge j\}
\]
为基。
% ===== PDF p.039 =====
\item 我们断言
\[
f:\C[X_1,X_2^{\pm1}]\xrightarrow{\ \sim\ }H(G,K)
\]
由
\[
X_1\mapsto K\begin{pmatrix}p&0\\0&1\end{pmatrix}K,\qquad
X_2\mapsto K\begin{pmatrix}p&0\\0&p\end{pmatrix}K
\]
给出。
\item 这里（以及今后）我们把双陪集与其特征函数等同。
\item 我们来验证这个断言（不借助 Gelfand 技巧）：
\begin{itemize}
\item 元素 $K\begin{pmatrix}p&0\\0&p\end{pmatrix}K$ 在 $H(G,K)$ 中中心，因为对一切 $g\in G$ 与群论中心中的 $z\in G$ 有
\[
KzK\cdot KgK=KgK\cdot KzK.
\]
\item 这说明 $f$ 良定义（注意 $\C[X_1,X_2^{\pm1}]$ 交换，而 $H(\GL_2(\Q_p),\GL_2(\Z_p))$ 先验不交换）。
\item 再验证 $f$ 满射。乘以 $f(X_2)$ 之后，只需说明每个
\[
K\begin{pmatrix}p^i&0\\0&1\end{pmatrix}K\qquad(i\ge1)
\]
都在 $f$ 的像中。
\item 用归纳法可见：对 $n\ge1$，$n$ 重乘积
\[
K\begin{pmatrix}p&0\\0&1\end{pmatrix}K\cdot\ldots\cdot
K\begin{pmatrix}p&0\\0&1\end{pmatrix}K
\]
与 $\mathrm{Mat}_{2,2}(\Z_p)$ 中行列式赋值等于 $n$ 的矩阵之集 $\mathrm{Mat}_{2,2}(\Z_p)^{\mathrm{val}=n}$ 一致。
\item 由 $H(G,K)$ 中乘法的定义可以推出，上述乘积就是 $\mathrm{Mat}_{2,2}(\Z_p)^{\mathrm{val}=n}$ 的特征函数 $\chi$。
\item 由于这个特征函数 $\chi$ 以 $K\begin{pmatrix}p^n&0\\0&1\end{pmatrix}K$ 的特征函数为系数 $1$ 的求和项，我们可以（用一个小归纳）断定它落在 $f$ 的像中，因为 $\mathrm{Mat}_{2,2}(\Z_p)^{\mathrm{val}=n}$ 中出现的其它双陪集都有形如
\[
\begin{pmatrix}p^j&0\\0&p^i\end{pmatrix},\qquad i,j\ge1,\ i+j=n
\]
的矩阵作为代表元。
\item 特别地，我们看到 $H(G,K)$ 必定交换，而无需借助 Gelfand 技巧。
% ===== PDF p.040 =====
\item 现在不难断定 $f$ 必为同构。事实上，
\[
\Spec(H(G,K))\subseteq\Spec(\C[X_1,X_2^{\pm1}])
\]
定义了一个子概形，它到 $\Spec(\C[X_2^{\pm1}])$ 的投影是无限秩自由的；而这对 $\Spec(\C[X_1,X_2^{\pm1}])$ 自身才可能（利用后者的整性）。
\end{itemize}
\end{itemize}

我们总结关于 Hecke 特征系统的讨论。

\textbf{要点}：
\begin{itemize}
\item 每个不可约容许 $\GL_2(\Af)$-表示 $V$ 都典范地给出一个 Hecke 特征系统
\[
\tilde a_p\in\C,\qquad \tilde b_p\in\C^\times
\]
对 $V$ 的分歧素数有限集 $S$ 之外的一切素数 $p$。
\item 具体地：若 $v\in V$ 被 $K(m)$ 固定且 $p\nmid m$，则
\[
\tilde T_p(v):=\GL_2(\Z_p)\begin{pmatrix}p&0\\0&1\end{pmatrix}\GL_2(\Z_p)*v=\tilde a_pv,
\]
其中 $\tilde T_p$ 是「$p$ 处的 Hecke 算子」，且
\[
\tilde S_p(v):=\GL_2(\Z_p)\begin{pmatrix}p&0\\0&p\end{pmatrix}\GL_2(\Z_p)*v=\tilde b_pv.
\]
\item 注意 Hecke 特征系统完全由复数族 $\{\tilde a_p,\tilde b_p\}_{p\notin S}$ 决定。
\item 令 $Z\subseteq\GL_2$ 为中心，即标量矩阵构成的子群。
\item 令 $\psi:\Af^\times\cong Z(\Af)\to\C^\times$ 为 $V$ 的中心特征（由 Schur 引理存在）。则
\[
\tilde b_p=\psi((1,\dots,1,p,1,\dots,1))
\]
对满足 $Z_p^\times\subseteq\Af^\times$ 在 $V$ 上平凡作用的素数 $p$ 成立。
\item 中心特征 $\psi$ 给出不可约容许 $\GL_2(\Af)$-表示的另一个不变量。
\item 注意 $\psi$ 一般不由 Hecke 特征系统决定，因为缺少有限多个 $p$，且没有要求 $V$ 在 $\Q^\times$ 上平凡，即没有要求「自守性」。
\end{itemize}

我们现在用上述讨论得到光滑容许 $\GL_2(\Af)$-表示 $S_k$ 的分解。先收集一些一般性观察，再得到精确的分解。

\textbf{关于 $\GL_2(\Af)$-表示 $S_k$ 的一般结果}：
\begin{itemize}
\item $\GL_2(\Af)$-表示 $S_k$ 分解为不可约容许表示 $V_i$ 的直和
\[
S_k\cong\bigoplus_{i\in I}V_i.
\]
\item 事实上：
% ===== PDF p.041 =====
\begin{itemize}
\item 在乘以特征 $|\det|_{\mathrm{ad\acute elique}}^{k/2}|_{\GL_2(\Af)}$ 的意义下，有一个 $\GL_2(\Af)$-等变嵌入 $S_k\subseteq L^2([\GL_2])$，从而 $\C$-向量空间 $S_k\otimes|\det|_{\mathrm{ad\acute elique}}^{k/2}$ 带有非退化的 $\GL_2(\Af)$-等变 Hermite 配对
\[
(\cdot,\cdot):S_k\times S_k\to\C,\qquad
(f,g)\mapsto\int_{[\GL_2]}\tilde\phi_f\overline{\tilde\phi_g}.
\]
\item 由此，通过取正交补，每个不可约 $\GL_2(\Af)$-子商其实已经是直和项。
\item 由 Zorn 引理，每个非零光滑 $\GL_2(\Af)$-表示 $V$ 都有不可约子商。事实上，不妨设 $V$ 由某个 $v\in V$ 生成；再对不含 $v$ 的子模之集合应用 Zorn 引理。
\end{itemize}
\end{itemize}

我们需要 Shalika 与 Piatetski-Shapiro 的下述结果。

\noindent\textbf{定理 5.5}（Shalika/Piatetski-Shapiro）\phantomsection\label{thm:5.5}\quad 对每个有限素数集 $S$ 与每个 Hecke 特征系统 $\{\tilde a_p,\tilde b_p\}_{p\notin S}$，至多存在一个具有该（等价类）Hecke 特征系统的不可约子表示 $V\subseteq S_k$。

\begin{itemize}
\item 这是 $\GL_n$ 尖点自守表示的「强重数一定理」的特例，参见 \texttt{[GH19, Theorem 11.7.2]}。\znoteplain{原文脚注：更精确地说，我们通过尖点自守形式的稠密性 \texttt{[GH19, Theorem 6.5.1.]} 与嵌入 $S_k\otimes|\det|_{\mathrm{ad\acute elique}}^{k/2}$，从 $L^2_{\mathrm{cusp}}([\GL_2])$ 的上述更强结果推出它。}
\item 我们可以记录下述推论：设 $0\neq f\in S_k$ 被 $K(m)$ 固定，并且是 $p\nmid m$ 时 $\tilde T_p$、$\tilde S_p$ 算子的本征向量。则由 $f$ 生成的 $\GL_2(\Af)$-表示 $V:=\langle f\rangle_{\GL_2(\Af)}$ 不可约。\znoteplain{「强重数一」这个名字容易让人误解：它并不是说每个不可约表示的重数为 $1$，而是说给定的 Hecke 特征系统（等价类）至多对应一个不可约子表示。正因如此，下面的定理 5.6 才能把 $S_k$ 的分解按新形式逐个列出，而不出现重数。}
\begin{itemize}
\item 因为每个非零 $v\in V$ 都给出一个 Hecke 特征系统，它与 $f$ 的等价；由强重数一\idx{强重数一}，它们必须落在同一子空间中。
\end{itemize}
\end{itemize}

现在我们把 Hecke 算子 $\tilde T_p$、$\tilde S_p$ 与它们在模形式上的经典对应物联系起来，并通过 Atkin/Lehner 的新形式\idx{新形式}描述 $S_k$ 的分解。

\textbf{以新形式为指标的 $S_k$ 分解}：
\begin{itemize}
\item 由 Atkin--Lehner/Casselman 的一个结果（参见 \texttt{[Cas73]}），对每个不可约 $\GL_2(\Af)$-子表示 $V\subseteq S_k$，存在某个 $m\in\N$ 与某个 $0\neq v\in V$，使 $v$ 被子群
\[
K_1(m):=\Bigl\{\begin{pmatrix}a&b\\c&d\end{pmatrix}\Bigm|c\equiv0,\ d\equiv1\ (\mathrm{mod}\ m)\Bigr\}\subseteq\GL_2(\widehat\Z)
\]
固定。
\item 令 $\Gamma_1(m):=\GL_2(\Z)\cap K_1(m)$。
\item 注记：
\begin{itemize}
% ===== PDF p.042 =====
\item 我们的记号与许多文献不同，那里 $\Gamma_1(m)\subseteq\SL_2(\Z)$。这一记号冲突是无害的，因为
\[
\Gamma_1(m)\backslash\HH^\pm\cong(\Gamma_1(m)\cap\SL_2(\Z))\backslash\HH.
\]
\item 使这样的非零 $v\in V$ 存在的最小 $m$ 称为 $V$ 的\textbf{导子}（conductor）。此时 $v$ 在相差一个标量的意义下唯一（称为该表示的新形式）。\znoteplain{导子是「级」的表示论版本：它记录表示在哪些素数处真正分歧，而不是记录一个具体的同余子群。Atkin--Lehner 理论的内容正是：每个不可约子表示都恰有一个极小级，且该级上的新形式（相差标量）唯一；其余同级的本征形式都由更低级的形式「旧形式」诱导而来。}
\item 更多细节见 \texttt{[Del73]}。
\end{itemize}
\item 我们已经注意到，态射
\[
\Gamma_1(m)\backslash\HH^\pm\xrightarrow{\ \sim\ }\GL_2(\Q)\backslash(\GL_2(\Af)/K_1(m)\times\HH^\pm),\qquad
[z]\mapsto[(1,z)]
\]
在第 3 讲中是同构。当然，该同构与 $\omega^{\otimes k}$ 相容。
\item 特别地，
\[
M_k(\Gamma_1(m))\cong M_k(K_1(m)),
\]
并且两者都与 $\Gamma_1(m)\cap\SL_2(\Z)$ 的经典模形式空间一致。
\item 让我们把这个同构写得更明确，取 $\tilde f\in M_k(K_1(m))$。
\item 设 $f\in M_k(\Gamma_1(m))$ 是对应于 $\tilde f$ 的截面。
\item 则
\[
\tilde f(g,z)=\tilde f(\gamma,z)=j(\gamma^{-1},z)^{-k}f(\gamma^{-1}z),
\]
其中 $g=\gamma h$，$\gamma\in\GL_2(\Q)$，$h\in K_1(m)$。
\item 注意这里用到了 $\GL_2(\Q)K_1(m)=\GL_2(\Af)$。
\item 固定某个满足 $p\nmid m$ 的素数 $p$。
\item 两个算子 $\tilde T_p$、$\tilde S_p$ 作用在 $M_k(K_1(m))$ 上。于是由结构的搬移，它们作用在 $M_k(\Gamma_1(m))$ 上。
\item 我们把 $\tilde S_p$ 在 $M_k(\Gamma_1(m))$ 上的这个作用写明确。
\item 对矩阵 $\gamma\in\GL_2(\Q)$，记 $\gamma_p\in\GL_2(\Af)$ 为第 $p$ 个分量为 $\gamma\in\GL_2(\Q)\subseteq\GL_2(\Q_p)$、其余分量为单位矩阵的元素。不要把此元素与 $\gamma_{\mathrm{diag}}\in\GL_2(\Af)$ 混淆，后者指由对角嵌入 $\GL_2(\Q)\subseteq\GL_2(\Af)$ 得到的元素。
\item $\tilde S_p$ 在 $M_k(K_1(m))$ 上的作用容易写出，但要在 $M_k(\Gamma_1(m))$ 上描述它，需找表达式
\[
\begin{pmatrix}p&0\\0&p\end{pmatrix}_p=\gamma_{\mathrm{diag}}h,\qquad
\gamma\in\GL_2(\Q),\ h\in K_1(m).
\]
\item 因为 $p$、$m$ 互素，存在 $a,b\in\Z$ 使
\[
A:=\begin{pmatrix}a&b\\m&p\end{pmatrix}\in\GL_2(\Z).
\]
\item 则
\[
A_{\mathrm{diag}}\cdot\begin{pmatrix}p^{-1}&0\\0&p^{-1}\end{pmatrix}_{\mathrm{diag}}\cdot
\begin{pmatrix}p&0\\0&p\end{pmatrix}_p\in K_1(m),
\]
% ===== PDF p.043 =====
其中用了前面引入的记号 $(-)_{\mathrm{diag}}$、$(-)_p$。
\item 换言之，可以取
\[
\gamma:=\begin{pmatrix}p&0\\0&p\end{pmatrix}A^{-1}.
\]
\item 取此 $\gamma$，算子 $\tilde S_p$ 通过
\[
f\mapsto\bigl(z\mapsto j(\gamma^{-1},z)^{-k}f(\gamma^{-1}z)\bigr)
\]
作用在 $M_k(\Gamma_1(m))$ 上。
\item 群
\[
\Gamma_0(m):=\Bigl\{\begin{pmatrix}a&b\\c&d\end{pmatrix}\Bigm|c\equiv0\ (\mathrm{mod}\ m)\Bigr\}
\]
正规化 $\Gamma_1(m)$，于是有 $\Gamma_0(m)/\Gamma_1(m)\cong(\Z/m)^\times$ 的作用：
\[
(B,f)\mapsto j(B,z)^{-k}f(Bz),\qquad B\in\Gamma_0(m),\ f\in M_k(\Gamma_1(m)).
\]
\znote{原文把 $\Gamma_0(m)$ 写成 $c\equiv1\ (\mathrm{mod}\ m)$，应为 $c\equiv0$：只有 $c\equiv0$ 时 $\Gamma_0(m)$ 才包含 $\Gamma_1(m)$（后者要求 $c\equiv0$、$d\equiv1$），$\Gamma_0(m)/\Gamma_1(m)\cong(\Z/m)^\times$ 与旁特征 $\chi(d)$ 的作用才自洽。此处按 $c\equiv0$ 译。[来源：原文 p.042 的 $\Gamma_1(m)$ 定义与 p.043 旁特征作用的相容性；同一笔误亦见于 p.032。]}
\item 于是得到分解
\[
M_k(\Gamma_1(m))=\bigoplus_{\chi:(\Z/m)^\times\to\C^\times}M_k(\Gamma_0(m),\chi),
\]
其中 $M_k(\Gamma_0(m),\chi)$ 按定义是 $\Gamma_0(m)$ 上旁特征为 $\chi$ 的模形式空间，即满足
\[
f\Bigl(\frac{az+b}{cz+d}\Bigr)=\chi(d)(cz+d)^kf(z)
\]
对 $\begin{pmatrix}a&b\\c&d\end{pmatrix}\in\Gamma_0(m)$ 成立的 $\Gamma_1(m)$-模形式之空间。
\item 对每个 Dirichlet 特征 $\chi:(\Z/m)^\times\to\C^\times$，存在唯一的 adélic 特征
\[
\psi_\chi:\Q^\times\backslash\A^\times\bigl/((1+m\widehat\Z)\cap\widehat\Z^\times\bigr)\to\C^\times,
\]
使 $\psi_\chi(r)=r^{-k}$ 对 $r\in\R_{>0}$ 成立。
\item 我们得到
\[
M_k(\Gamma_0(m),\chi)\cong M_k(K_1(m),\psi_\chi),
\]
其中右端指 $\Af^\times$（$\GL_2(\Af)$ 的中心）在 $M_k(K_1(m))$ 上作用的 $\psi_\chi$-本征空间。
\item 用这套术语可以完成对算子 $\tilde S_p$ 在 $M_k(\Gamma_1(m))$ 上的描述。
\item 即：$\tilde S_p$ 在子空间 $M_k(\Gamma_0(m),\chi)\subseteq M_k(\Gamma_1(m))$ 上的作用，由
\[
\gamma^{-1}=\begin{pmatrix}p^{-1}&0\\0&p^{-1}\end{pmatrix}A
\]
的作用给出；而 $\begin{pmatrix}p^{-1}&0\\0&p^{-1}\end{pmatrix}$ 作用为 $p^k$，$A\in\Gamma_0(m)$ 作用为 $\chi(p)$。
\item 现在描述 $\tilde T_p$ 在 $M_k(\Gamma_0(m),\chi)$ 上的作用（由结构的搬移）。
\item 有
\[
\GL_2(\Z_p)\begin{pmatrix}p&0\\0&1\end{pmatrix}\GL_2(\Z_p)
=\bigsqcup_{j=0}^{p-1}\begin{pmatrix}p&j\\0&1\end{pmatrix}\GL_2(\Z_p)
\sqcup\begin{pmatrix}1&0\\0&p\end{pmatrix}\GL_2(\Z_p).
\]
% ===== PDF p.044 =====
\item 我们同样要把（上述代表的）双陪集分解为
\[
\gamma\cdot h,\qquad \gamma\in\GL_2(\Q),\ h\in K_1(m).
\]
\item 可以取分解
\[
\begin{pmatrix}p&j\\0&1\end{pmatrix}_p
=\begin{pmatrix}p&j\\0&1\end{pmatrix}_{\mathrm{diag}}\cdot h
\]
对 $j=0,\dots,p-1$，其中 $h\in K_1(m)$，即
\[
\gamma^{-1}=\begin{pmatrix}p^{-1}&-p^{-1}j\\0&1\end{pmatrix}.
\]
\item 对 $\begin{pmatrix}1&0\\0&p\end{pmatrix}$，与上面的论证类似，可得
\[
\gamma^{-1}=\begin{pmatrix}a&b\\m&p\end{pmatrix}\cdot\begin{pmatrix}1&0\\0&p^{-1}\end{pmatrix}
=\begin{pmatrix}1&0\\0&p^{-1}\end{pmatrix}\cdot\begin{pmatrix}a&p^{-1}b\\pm&p\end{pmatrix},
\]
如所欲。
\item 把这些放在一起，得到
\[
\tilde T_p(f)=p^k\chi(p)f(pz)+\sum_{j=0}^{p-1}f\Bigl(\frac{z-j}{p}\Bigr).
\]
\item 若把带旁特征的模形式 $f\in M_k(\Gamma_0(m),\chi)$ 写成 Fourier 展开
\[
f(z)=\sum_{n=0}^\infty a_nq^n,\qquad q=e^{2\pi iz},
\]
则更具体地
\[
\tilde T_p(f)(z)=\sum_{n=0}^\infty pa_{pn}q^n+\sum_{n=0}^\infty p^k\chi(p)a_nq^{pn}.
\]
\item 这由下述事实推出：
\[
\sum_{j=0}^{p-1}e^{2\pi i\frac{-jn}{p}}=
\begin{cases}
p,&p\mid n,\\
0,&\text{否则}.
\end{cases}
\]
\item 特别地，若 $f\in M_k(\Gamma_0(m),\chi)$ 是 $\tilde T_p$ 的本征向量且 $a_1=1$，则（考察 $q^1$ 前的系数）
\[
\tilde T_p(f)=pa_pf.
\]
这促使我们定义
\[
T_p=\frac1p\tilde T_p,
\]
它就是作用在模形式上的经典 Hecke 算子（如同 \texttt{[DS05, Section 5.2.]}）。换言之，$T_p$ 算子的 Hecke 特征值就是归一化本征形式的 Fourier 系数。\znoteplain{这里必须留意两套归一化：$\tilde T_p$ 是双陪集的特征函数（不除体积）给出的算子，而经典 $T_p$ 是它在 $\GL_2(\Af)$ 侧的「平均」版本，两者相差因子 $p$（正好是 $\tilde T_p$ 中 $p$ 个斜陪集的个数）。因此同一个本征形式满足 $\tilde T_p f=pa_pf$ 与 $T_pf=a_pf$，本征值 $\tilde a_p=pa_p$ 与 $a_p$ 只在差一个 $p$ 的意义下相同；比较两种文献时若忘了这个因子，就会得到相差 $p$ 的「矛盾」。}
\end{itemize}

% ===== PDF p.045 =====
我们现在可以证明关于 $\GL_2(\Af)$-表示 $S_k$ 分解的定理。

\noindent\textbf{定理 5.6}\phantomsection\label{thm:5.6}\quad $\GL_2(\Af)$-表示 $S_k$ 分解为不可约表示
\[
S_k\cong\bigoplus_f\langle f\rangle_{\GL_2(\Af)},
\]
其中 $f$ 遍历 $\Gamma_0(N)$ 上、旁特征为 $\chi$ 的（归一化）新形式，$N$ 与 Dirichlet 特征 $\chi:(\Z/N)^\times\to\C^\times$ 都在变。

我们不给出新形式的精确定义，参见 \texttt{[DS05, Section 5.8.]}。粗略地说，新形式是在所有具有等价 Hecke 特征系统的本征形式中实现极小 $N$ 的本征形式。

\textbf{证明.} 经过上述准备，我们现在可以引用 \texttt{[DS05, Proposition 5.8.4]}、\texttt{[DS05, Theorem 5.8.2]}，它们蕴涵：对 $S_k$ 中出现的每个 Hecke 特征系统，存在唯一的归一化（即 Fourier 展开中 $a_1=1$）新形式，其 Hecke 特征系统与之等价。\hfill$\square$

在其它教材中，例如 \texttt{[Gel75, Proposition 3.1]}，自守表示直接联系到某个 $f\in M_k(\Gamma_0(m),\chi)$，即不引入空间 $M_k$，而用 $\GL_2(\Af)$-等变嵌入
\[
M_k\to C^\infty(\GL_2(\Q)\backslash\GL_2(\A)),\qquad
f\mapsto\phi_f(g,g_\infty)=j(g_\infty,i)^{-k}f(g,g_\infty i),
\]
以及
\[
S_k\otimes|\det|_{\mathrm{ad\acute elique}}^{k/2}\to L^2([\GL_2]),\qquad
f\mapsto\tilde\phi_f(g,g_\infty)=|\det(g,g_\infty)|_{\mathrm{ad\acute elique}}^{k/2}\phi_f.
\]
我们指出，在上述同构 $M_k(\Gamma_1(m))\cong M_k(K_1(m))$ 下，我们的公式特殊化为 \texttt{[Gel75, Proposition 3.1.]} 中的公式。此外，\texttt{[Gel75, Lemma 3.7.]} 中出现的看似意外的因子 $p^{k/2-1}$ 有一个简单的解释：出现了 $|\det|_{\mathrm{ad\acute elique}}^{k/2}$ 的扭，以及从（我们的）$\tilde T_p$ 到经典 Hecke 算子 $T_p=\frac1p\tilde T_p$ 的过渡。

\chapter{新形式的 Langlands 互反}

% ===== PDF p.046 =====
\textbf{上次}：
\begin{itemize}
\item 我们证明了分解
\[
S_k\cong\bigoplus_f\langle f\rangle_{\GL_2(\Af)}
\]
为不可约 $\GL_2(\Af)$-表示。
\item 这里 $f$ 遍历权 $k$、$\Gamma_0(N)$ 上、旁特征为 $\chi$ 的（归一化）新形式，$N$ 在变，Dirichlet 特征 $\chi:(\Z/N)^\times\to\C^\times$ 也在变。
\item 当然，在证明该分解的过程中，我们对每个不可约容许 $\GL_2(\Af)$-表示 $\pi$ 联系了一个 Hecke 特征系统
\[
\{\tilde a_p(\pi),\tilde b_p(\pi)\}_{p\notin S},
\]
其中 $S:=\{\,p\ \text{素数}:\pi^{\GL_2(\Z_p)}=0\,\}$。事实上，$\tilde a_p(\pi)$（相应地 $\tilde b_p(\pi)$）定义为双陪集
\[
\GL_2(\Z_p)\begin{pmatrix}p&0\\0&1\end{pmatrix}\GL_2(\Z_p)
\quad\text{相应地}\quad
\GL_2(\Z_p)\begin{pmatrix}p&0\\0&p\end{pmatrix}\GL_2(\Z_p)
\]
作用在由 $\GL_2(\Z_p)$ 固定的非零向量上的本征值。
\item 若 $\pi=\langle f\rangle_{\GL_2(\Af)}$，其中 $f\in S_k(\Gamma_0(N),\chi)$ 为归一化新形式，则 $S=\{\,N\ \text{的素因子}\,\}$，且
\[
\tilde a_p(\pi)=pa_p,\qquad \tilde b_p(\pi)=p^k\chi(p)
\]
对 $p\notin S$ 成立，其中 $f(q)=\sum_{n=1}^\infty a_nq^n$ 是 $f$ 的 Fourier 展开。
\end{itemize}

本讲我们将引入 $\ell$-进 Galois 表示的 Frobenius 的迹（它们是 Hecke 特征系统的「Galois 论」对应物），然后陈述 Langlands 纲领对 $S_k$ 的预测。

\textbf{$\ell$-进表示的 Frobenius 的迹}：
\begin{itemize}
\item 设 $\ell$ 为某个素数。
\item
\[
\overline\Q_\ell=\varinjlim_{E/\Q_\ell}E
\]
通过余极限拓扑成为拓扑域。
\item 设 $W$ 为有限维 $\overline\Q_\ell$-向量空间。
\item 设 $\sigma:G_\Q:=\Gal(\overline\Q/\Q)\to\GL(W)$ 为连续表示，也称为「$\ell$-进 Galois 表示」。
\item 回忆对每个素数 $p$ 存在嵌入
\[
G_{\Q_p}:=\Gal(\overline\Q_p/\Q_p)\hookrightarrow G_\Q,
\]
它在 $G_\Q$ 中相差共轭的意义下良定义。
\item 局部绝对 Galois 群 $G_{\Q_p}$ 处于正合列
\[
1\to I_p\to G_{\Q_p}\to\Gal(\overline\F_p/\F_p)=\langle\Frob_p^{\mathrm{geom}}\rangle^{\widehat\Z}\to1,
\]
其中 $I_p\subseteq G_{\Q_p}$ 是惯性子群，而
\[
\Frob_p^{\mathrm{arith}}:\overline\F_p\to\overline\F_p,\qquad x\mapsto x^p
\]
% ===== PDF p.047 =====
是 $p$ 处的算术 Frobenius。
\item 逆元 $\Frob_p^{\mathrm{geom}}:=(\Frob_p^{\mathrm{arith}})^{-1}$ 称为\textbf{几何 Frobenius}。
\item 设 $S$ 为有限素数集。称 $\sigma$ 在 $S$ 外\textbf{非分歧}，若 $\sigma(I_p)=1$ 对一切 $p\notin S$ 成立。
\item 若 $\sigma$ 在 $S$ 外非分歧，$n:=\dim_{\overline\Q_\ell}W$，则可以对 $\sigma$ 联系一个系统 $\{c_{1,p},\dots,c_{n,p}\}_{p\notin S}$，$c_{i,p}\in\overline\Q_\ell$，即
\[
c_{i,p}:=\Tr\bigl(\sigma(\Frob_p^{\mathrm{arith}})\mid\Lambda^iW\bigr).
\]
\item 注：$c_{i,p}$ 良定义，因为 $\sigma(I_p)=1$。
\item 注：$c_{i,p}$ 只依赖于 $p$，而不依赖于嵌入 $G_{\Q_p}\subseteq G_\Q$，因为迹在共轭下不变。\znote{原文此处第二条「注」把 $c_{i,p}$ 写成 $\Tr(\sigma(\Frob_p^{\mathrm{geom}})\mid\Lambda^iW)$，但 $c_{i,p}$ 的定义用的是算术 Frobenius，且 p.050 的页边注明确说用几何 Frobenius 是错的。故此处按算术 Frobenius 读，该注只保留「不依赖嵌入的选取」这一层意思。[来源：原文 p.047 的定义与 p.050 页边注。]}
\item 若 $W$ 半单，则 $\sigma$ 由 $\{c_{1,p},\dots,c_{n,p}\}_{p\notin S}$ 决定（对每个使 $W$ 在 $S$ 外非分歧的有限素数集 $S$），参见 \texttt{[Ser97, I-10]}。
\item 事实上，族 $\{c_{1,p}\}_{p\notin S}$ 已经足够（再加上 $W$ 半单且在 $S$ 外非分歧的假设）。
\item 确实：
\begin{itemize}
\item 考虑群代数 $\Lambda:=\overline\Q_\ell[G_{\Q,S}]$，其中 $G_{\Q,S}$ 是 $G_\Q$ 对由 $I_p$（$p\notin S$）生成的子群的闭包所取之商。
\item 设 $W,W'$ 为两个半单连续 $G_{\Q,S}$-表示，且
\[
\Tr(\Frob_p^{\mathrm{arith}}|W)=\Tr(\Frob_p^{\mathrm{arith}}|W')
\]
对一切 $p\notin S$ 成立。
\item 由 Chebotarev 密度定理与连续性，这蕴涵
\[
\Tr(\lambda|W)=\Tr(\lambda|W')
\]
对一切 $\lambda\in\Lambda$ 成立。
\item 对两两不同构的不可约 $\Lambda$-模 $W_1,\dots,W_m$ 与每个 $i=1,\dots,m$，存在 $\mu_i\in\Lambda$，使 $\mu_i=1$ 在 $W_i$ 上、而 $\mu_i=0$ 在 $W_j$ 上（$j\neq i$）（这是中国剩余定理的一个版本）。
\item 把
\[
W=\bigoplus_{i=1}^mW_i^{\oplus n_i},\qquad
W'=\bigoplus_{i=1}^mW_i^{\oplus n'_i}
\]
写成同型分量，其中 $W_1,\dots,W_m$ 不可约、两两不同构，$n_i,n'_i\in\N$（可能为零）。
\item 则
\[
n_i\dim_{\overline\Q_\ell}W_i=\Tr(\mu_i|W)=\Tr(\mu_i|W')=n'_i\dim_{\overline\Q_\ell}W_i
\]
对一切 $i=1,\dots,m$ 成立，即 $n_i=n'_i$，故 $W\cong W'$，如所欲。
\end{itemize}
\item 这个陈述可以看作上次提到的 $\GL_n$ 尖点自守表示「强重数一」的类比。
% ===== PDF p.048 =====
\item 重要的是我们考虑的是特征 $0$ 系数的表示；也就是说，「仅凭迹 $c_{1,p}$ 就决定（半单）表示」这一说法是错误的。\znoteplain{原文此处写作 $c_{i,p}$；但按下一句（完整的系统 $\{c_{1,p},\dots,c_{n,p}\}$ 才唯一决定）与上下文，这里指的应是单个迹 $c_{1,p}$，故译作 $c_{1,p}$。}由 Brauer--Nesbitt 定理，完整的系统 $\{c_{1,p},\dots,c_{n,p}\}_{p\notin S}$ 仍然唯一决定有限维半单表示。\znoteplain{特征 $0$ 时半单表示由迹函数唯一决定，这就是 Brauer--Nesbitt 定理的通常形式；正特征时只知道 $\Tr$ 不够，必须知道全部外幂 $\Lambda^iW$ 上的迹，这正是上面要把 $\{c_{1,p},\dots,c_{n,p}\}$ 整族留下来的原因。}
\end{itemize}

Frobenius 元的迹含有高度的算术信息。下面的讨论对理解 Langlands 纲领的算术蕴涵非常重要。

\textbf{Frobenius 的迹的算术意义}：
\begin{itemize}
\item 设 $\sigma:G_\Q\to\GL(W)$ 为 $\ell$-进表示，在 $S$ 外非分歧。
\item 迹
\[
c_{1,p}:=\Tr\bigl(\sigma(\Frob_p^{\mathrm{arith}})\mid W\bigr)
\]
编码重要的算术信息。
\item 例如，设 $F/\Q$ 为有限 Galois 扩张，Galois 群 $\Gal(F/\Q)\cong S_3$，且
\[
\sigma:G_\Q\twoheadrightarrow\Gal(F/\Q)\to\GL_2(\overline\Q_\ell)
\]
是由 $S_3$ 的唯一二维不可约表示 inflation 得到的。令 $S:=\{\,p\ \text{在}\ F\ \text{中分歧}\,\}$，并对 $p\notin S$ 记 $\Frob_{F,p}^{\mathrm{arith}}$ 为 $\Frob_p^{\mathrm{arith}}\in G_{\Q,S}$ 在 $\Gal(F/\Q)\cong S_3$ 中的像。
\item 则：
\begin{itemize}
\item $\Tr(\Frob_p^{\mathrm{arith}}|W)=2$ $\Leftrightarrow$ $\Frob_{F,p}^{\mathrm{arith}}\in\{1\}$ $\Leftrightarrow$ $p$ 在 $\OO_F$ 中完全分裂；
\item $\Tr(\Frob_p^{\mathrm{arith}}|W)=0$ $\Leftrightarrow$ $\Frob_{F,p}^{\mathrm{arith}}\in\{(1\,2),(1\,3),(2\,3)\}$ $\Leftrightarrow$ $p$ 在 $\OO_F$ 中分裂为三个不同素数；
\item $\Tr(\Frob_p^{\mathrm{arith}}|W)=-1$ $\Leftrightarrow$ $\Frob_{F,p}^{\mathrm{arith}}\in\{(1\,2\,3),(1\,3\,2)\}$ $\Leftrightarrow$ $p$ 在 $\OO_F$ 中分裂为两个不同素数。
\end{itemize}
\item 因此：知道一切（有限像的）$\ell$-进 Galois 表示的 Frobenius 的迹，就蕴涵知道一切有限扩张中非分歧素数的分解。
\item $\ell$-进表示的其它例子是平展上同调群
\[
H^i_{\text{ét}}(X_{\overline\Q},\overline\Q_\ell)
\]
对 $\Q$ 上 proper、光滑的概形 $X$ 与 $i\ge0$。
\item $G_\Q$-作用由 $H^i_{\text{ét}}(-,\overline\Q_\ell)$ 的函子性，从 $G_\Q$ 在
\[
X_{\overline\Q}=X\times_{\Spec(\Q)}\Spec(\overline\Q)
\]
上的（右）作用诱导。
% ===== PDF p.049 =====
\item proper 光滑基变换与 Grothendieck--Lefschetz 迹公式蕴涵：若 $X$ 在 $p$ 处有好约化且 $p\neq\ell$，则 $H^i_{\text{ét}}(X_{\overline\Q},\overline\Q_\ell)$ 在 $p$ 处非分歧（$i\ge0$），且
\[
\sum_{i=0}^{2\dim(X)}(-1)^i\Tr\bigl(\Frob_p^{\mathrm{geom}}\mid H^i_{\text{ét}}(X_{\overline\Q},\overline\Q_\ell)\bigr)=\sharp X(\F_p),
\]
其中 $X\to\Spec(\Z_{(p)})$ 是 $X$ 在 $p$ 处的 proper 光滑模型。
\end{itemize}

从算术的观点看，现在可以陈述 Langlands 纲领的一个要点。

\textbf{Langlands 互反\idx{Langlands 互反}的全部要点}：
\begin{itemize}
\item Frobenius 的迹应当与 Hecke 特征值相匹配！
\item 于是，算术信息被期望编码在自守信息中，而后者可能更易把握。
\item 下例取自 Matthew Emerton 的一篇帖子。他更有洞见的帖子可在其主页找到：\url{http://www.math.uchicago.edu/~emerton/}。
\begin{itemize}
\item 设 $F$ 为 $x^3-x-1$ 的分裂域。
\item 则 $F$ 是 $\Q(\sqrt{-23})$ 的 Hilbert 类域，且 $F/\Q$ 是 Galois 扩张，Galois 群为 $S_3$。
\item 像上面一样考虑相伴的二维不可约 Galois 表示 $\sigma:G_\Q\to\Gal(F/\Q)\hookrightarrow\GL_2(\overline\Q_\ell)$。
\item Hecke 证明了存在归一化新形式 $f(q)=\sum_{n=1}^\infty a_nq^n$，使
\[
a_p=\Tr\bigl(\sigma(\Frob_p^{\mathrm{arith}})\bigr)
\]
对一切满足 $p\nmid23$ 的素数 $p$ 成立。
\item 可以把 $f$ 写得更明确：考虑分歧，$f$ 必落在 $S_1(\Gamma_0(23),\chi)$ 中，其中 $\chi:(\Z/23)^\times\to\C^\times$ 是决定二次子域 $\Q(\sqrt{-23})\subseteq\Q(\zeta_{23})$ 的二次特征。
\item 熟悉模形式的人（或使用 LMFDB 之类数据库的人）会告诉我们必有
\[
f(q)=q\prod_{i=1}^\infty(1-q^i)(1-q^{23i}).
\]
\item 现在可以计算 $f$ 的展开（或在 LMFDB 中查它）：
\[
f(q)=q-q^2-q^3+\ldots+q^{58}+2q^{59}+\ldots
\]
并推断素数在 $F$ 中如何分解，例如 $59$ 完全分裂！
\end{itemize}
\item 也许 Frobenius 的迹的模性最著名的应用是 Fermat 大定理的证明。
\item 即：Fermat 定理由 Wiles/Taylor 用下述策略证明（策略由 Frey 提出，Serre、Ribet 补全）：
\begin{itemize}
\item 设 $u^p+v^p+w^p=0$，$u,v,w\in\Q$，$uvw\neq0$，$p\ge3$。
\item 考虑 $\Q$ 上由（仿射）Weierstraß 方程
\[
y^2=x(x+u^p)(x-v^p)
\]
% ===== PDF p.050 =====
（在把 $u,v,w$ 稍作调整之后）给出的椭圆曲线\idx{椭圆曲线} $E$。
\item 考虑 Galois 表示
\[
G_\Q\to\GL\bigl(H^1_{\text{ét}}(E_{\overline\Q},\overline\Q_\ell)\bigr)\cong\GL_2(\overline\Q_\ell)
\]
的对偶 $\sigma$。
\item 证明 $\sigma$ 是模的，即存在权 $2$ 的归一化新形式 $f_1(q)=\sum_{n=1}^\infty a_nq^n$，使
\[
\Tr\bigl(\sigma(\Frob_p^{\mathrm{arith}})\bigr)=a_p
\]
对几乎所有素数 $p$ 成立。
\item 用 Ribet 的一个定理可推出 $f_1$ 必与
\[
S_2(\Gamma_0(2))
\]
中的某个归一化新形式 $f_2$ 同余，即 $f_1,f_2$ 的 Fourier 系数是代数整数，且在某个素数模下同余。
\item 但 $S_2(\Gamma_0(2))=0$，故 $f_2$ 不存在。
\item 于是得到所要的矛盾。
\item 更多细节参见 \texttt{[Wil95]}、\texttt{[Rib90]}。
\end{itemize}
\end{itemize}

现在陈述关于 Galois 表示与尖点形式之间关系的预期。

\textbf{新形式的 Langlands 互反}：
\begin{itemize}
\item 固定素数 $\ell$ 与同构
\[
\iota:\overline\Q_\ell\xrightarrow{\ \sim\ }\C.
\]
\item 对新形式，Langlands 纲领结合 Fontaine--Mazur 猜想预测存在双射
\[
\mathrm{LL}:\cA_{\mathrm{mod}}\xrightarrow{\ \sim\ }\mathcal G_{\mathrm{mod}},
\]
从集合
\[
\cA_{\mathrm{mod}}:=\Bigl\{\ \text{不可约}\ \GL_2(\Af)\text{-子表示}\ \pi\subseteq\bigoplus_{k\ge1}S_k\ \Bigr\}
\]
到某个集合 $\mathcal G_{\mathrm{mod}}$，后者由不可约二维 Galois 表示
\[
\sigma:G_\Q\to\GL_2(\overline\Q_\ell)
\]
组成。
\item $\mathrm{LL}$ 应满足：对一切 $\pi\in\cA_{\mathrm{mod}}$，令 $\sigma:=\mathrm{LL}(\pi)$，有
\[
\iota\bigl(\Tr(\sigma(\Frob_p^{\mathrm{arith}}))\bigr)=a_p(\pi),\qquad
\iota\bigl(\det(\sigma(\Frob_p^{\mathrm{arith}}))\bigr)=b_p(\pi)
\]
对某个指定的有限素数集 $S$ 之外的 $p$ 成立。\znote{原文此处页边注「In a previous version I considered the geometric Frobenius, which was wrong」重复排印了两次（同一句话上下各一遍），译文只保留一次。这句话针对的正是上面的匹配公式：在二维情形 $\Frob_p^{\mathrm{geom}}=(\Frob_p^{\mathrm{arith}})^{-1}$，而 $\Tr(A^{-1})=\Tr(A)/\det(A)$，若误用几何 Frobenius，得到的将是 $a_p/b_p$ 而不是 $a_p$。[来源：原文 p.050 页面的页边注与 p.047 的 Frobenius 定义。]}
\item 这里使用的是本征值 $a_p(\pi)$（相应地 $b_p(\pi)$），而不是 $\tilde a_p(\pi)$、$\tilde b_p(\pi)$。\znoteplain{固定 $\iota:\overline\Q_\ell\cong\C$ 的作用是把 $\ell$-进数（Frobenius 的迹与行列式）搬到复数域，才能与 Fourier 系数 $a_p(\pi),b_p(\pi)\in\C$ 比较。换一个 $\iota$ 会把所有这些代数数同时换成它们的共轭，双射 $\mathrm{LL}$ 也随之改变，所以 $\iota$ 必须一开始就固定。}
\item 注：这样的双射 $\mathrm{LL}$ 若存在，则由相应的重数一定理唯一决定。
\item 猜想地，集合 $S$ 有明确的描述。回忆素数 $p$ 对 $\sigma$（相应地 $\pi$）非分歧，若
\[
\sigma(I_p)=1
\qquad\text{相应地}\qquad
\pi^{\GL_2(\Z_p)}\neq0.
\]
% ===== PDF p.051 =====
\item 则猜想地，$p\neq\ell$ 对 $\pi$ 非分歧当且仅当它对 $\sigma:=\mathrm{LL}(\pi)$ 非分歧。
\item Hecke 特征值的匹配更精确地是说：$\sigma(\Frob_p^{\mathrm{arith}})$ 的特征多项式应为
\[
X^2-\iota^{-1}(a_p(\pi))X+\iota^{-1}(b_p(\pi)),
\]
对 $\pi$ 的每个非分歧素数 $p\neq\ell$。
\item 课程的下一目标是说明 $\mathrm{LL}$ 可以如何构造，即如何把 Galois 表示联系到归一化新形式。为此我们遵循 Deligne 的结果（参见 \texttt{[Del71b]}）与 Deligne/Serre 的结果（参见 \texttt{[DS74]}）。
\end{itemize}

下面更精确地说明预期与一般新形式相伴的是哪个 Galois 表示集合。

\textbf{集合 $\mathcal G_{\mathrm{mod}}$}：
\begin{itemize}
\item 集合 $\mathcal G_{\mathrm{mod}}$ 由 Fontaine/Mazur 在他们关于几何 Galois 表示的著名猜想中定义，参见 \texttt{[FM95]}。
\item 设 $\sigma:G_\Q\to\GL_2(\overline\Q_\ell)$ 为不可约二维 $\ell$-进表示。
\item 则按定义 $\sigma\in\mathcal G_{\mathrm{mod}}$ 当且仅当
\begin{itemize}
\item $\sigma$ 在几乎所有 $p$ 处非分歧，即 $\sigma(I_p)=1$ 对某个有限素数集 $S$ 之外的 $p$ 成立；
\item $\sigma$ 是\textbf{奇的}（odd），即对每个复共轭 $c\in\Gal(\C/\R)\subseteq G_\Q$ 有
\[
\det(\sigma(c))=-1;
\]
\item 限制
\[
\sigma_\ell:=\sigma|_{G_{\Q_\ell}}:G_{\Q_\ell}\to\GL_2(\overline\Q_\ell)
\]
是 de Rham 的，Hodge--Tate 权\idx{Hodge--Tate 权}为 $0,\omega$（某个 $\omega\in\N$；注意 $\omega=0$ 是允许的）。
\end{itemize}
\item 若 $\sigma=\mathrm{LL}(\pi)$ 且 $\pi\subseteq S_k$，则猜想地 $\omega=k-1$。
\item 利用许多人的工作（Wiles、Kisin、Emerton、Skinner--Wiles、Pan、$\dots$），任何 Hodge--Tate 权互不相同的 $\sigma\in\mathcal G_{\mathrm{mod}}$ 在 $\ell\ge5$ 时是模的，参见 \texttt{[Pan19, Theorem 1.0.4.]}。
\item 若 $\sigma\in\mathcal G_{\mathrm{mod}}$ 有有限像，则 $\sigma$ 是模的，参见 \texttt{[PS16]}、\texttt{[GH19, Section 13.4.]}（以及其中的参考文献）。
\item 猜想地：全部 Hodge--Tate 权为 $0$ $\Leftrightarrow$ $\sigma$ 的像有限。
\end{itemize}

我们必须对 de Rham 这一重要概念作一段简短的离题讨论。

\textbf{de Rham 表示\idx{de Rham 表示}}：
\begin{itemize}
\item 设 $p$ 为素数（即上面的 $\ell$）。
\item 设 $K/\Q_p$ 为离散赋值、非 Archimedes 扩张，剩余域完美（例如 $K/\Q_p$ 有限）。
\item 设 $\sigma:G_K:=\Gal(\overline K/K)\to\GL_m(\Q_p)$ 为连续表示。
\item 则 $\sigma$ 的系数落在某个有限扩张 $E/\Q_p$ 中（因为 $G_K$ 紧，而这里 $\Q_p$ 带余极限拓扑）。
\item 于是得到 $\Q_p$-线性表示，即「局部 $p$-进 Galois 表示」
\[
\rho:G_K\to\GL_m(E)\subseteq\GL_n(\Q_p),\qquad n:=m\cdot\dim_{\Q_p}E.
\]
% ===== PDF p.052 =====
\item 于是可以改变设置，令 $V$ 为有限维 $\Q_p$-向量空间，$\rho:G_K\to\GL(V)$ 为连续表示。
\item 则 $V$ 称为 \textbf{de Rham 的}，若
\[
\dim_K\bigl(B_{\mathrm{dR}}\otimes_{\Q_p}V\bigr)^{G_K}=\dim_{\Q_p}V
\]
对某个带 $G_K$-作用的域扩张 $B_{\mathrm{dR}}$ 成立，参见 \texttt{[BC09, Section 6]}。
\item 即：$B_{\mathrm{dR}}$ 是 Fontaine 的 $p$-进周期域，参见 \texttt{[Fon94]}、\texttt{[BC09, Definition 4.4.7]}。抽象地，
\[
B_{\mathrm{dR}}\cong C_K((t)),
\]
其中 $C_K$ 是 $K$ 对其 $p$-进赋值的完备化（这个拓扑比 $K$ 上的余极限拓扑更粗）。
\item $B_{\mathrm{dR}}$ 是离散赋值域，剩余域为 $C_K$；由此得到的递减滤过 $\mathrm{Fil}^\bullet B_{\mathrm{dR}}$ 是 $G_K$-稳定的，且作为 $C_K$-半线性 $G_K$-表示
\[
\mathrm{gr}^\bullet B_{\mathrm{dR}}\cong\bigoplus_{j\in\Z}C_K(j),
\]
其中 $C_K(j):=C_K\otimes_{\Z_p}\Z_p(j)$，而
\[
\Z_p(1):=T_p\mu_{p^\infty}(C_K)
\]
是与分圆特征 $\chi_{\mathrm{cyc}}:G_K\to\Z_p^\times\cong\Aut(\mu_{p^\infty}(C_K))$ 相伴的 $p$-进 Galois 表示。
\item Tate 的一个重要定理（参见 \texttt{[BC09, Theorem 2.2.7]}）说：设 $\chi:G_K\to\Q_p^\times$ 为特征。则
\[
H^0\bigl(G_K,C_K(\chi)\bigr)=
\begin{cases}
K,&\chi|_{I_K}\ \text{的像有限},\\
0,&\text{否则},
\end{cases}
\]
这里 $I_K\subseteq G_K$ 是惯性子群。
\item 特别地：
\[
H^0\bigl(G_K,C_K(j)\bigr)=0,\qquad j\neq0,
\]
因为 $j\neq0$ 时 $\chi_{\mathrm{cyc}}^j|_{I_K}$ 的像无限。
\item 结论：若 $\dim_{\Q_p}V=1$，则 $V$ 是 de Rham 的当且仅当 $V\cong\chi\cdot\chi_{\mathrm{cyc}}^j$，其中 $\chi:G_K\to\Q_p^\times$ 是某个使 $I_K$ 的像有限的特征，$j\in\Z$。特别地，存在许多非 de Rham 的表示，例如特征 $(\chi_{\mathrm{cyc}}^{p-1})^a$，$a\in\Z_p\setminus\Z$。\znote{原文作 $\dim_\C V=1$；这里 $V$ 是 $G_K$ 的 $\Q_p$-表示，而 de Rham 判据比较的是 $\dim_K(B_{\mathrm{dR}}\otimes_{\Q_p}V)^{G_K}$ 与 $\dim_{\Q_p}V$，故应读作 $\dim_{\Q_p}V=1$，否则整个陈述没有意义。[来源：原文 p.052 上一条 de Rham 定义。]}
\item 我们提一些关于 de Rham 表示的一般事实。
\begin{itemize}
\item 若存在 $K$ 的有限扩张 $K'$ 使 $\rho|_{G_{K'}}$ 是 de Rham 的，则 $\rho:G_K\to\GL(V)$ 是 de Rham 的，参见 \texttt{[BC09, Proposition 6.3.8]}。
\item de Rham 表示对子商、对偶与张量积稳定，参见 \texttt{[BC09, Section 6.1]}。
\end{itemize}
\item 设 $V$ 为 $n$ 维 de Rham 表示。则
\[
V\otimes_{\Q_p}C_K\cong C_K(j_1)\oplus\ldots\oplus C_K(j_n)
\]
作为 $C_K$-半线性 $G_K$-表示，其中 $j_1,\dots,j_n\in\Z$。无序族 $(j_1,\dots,j_n)$ 称为 $V$ 的 \textbf{Hodge--Tate 权}之族。
\end{itemize}

% ===== PDF p.053 =====
因此，de Rham 条件给 $p$-进表示加上了某种「整性」条件。

我们以说明 $\mathrm{LL}$ 如何构造（当 $k=2$ 时）来结束本讲。

\textbf{构造 $\mathrm{LL}$ 的策略（粗略）}：
\begin{itemize}
\item 设 $\pi\in\cA_{\mathrm{mod}}$，相伴的 Hecke 特征系统为 $\{a_p,b_p\}_{p\notin S}$。
\item 为简单起见设 $\pi\subseteq S_2$，即 $\pi$ 由某个权 $2$ 的新形式生成。
\item 则我们将构造（接近）一个 $\GL_2(\Af)$-等变嵌入
\[
S_2\hookrightarrow H^1\bigl(\GL_2(\Q)\backslash(\GL_2(\Af)\times\HH^\pm),\C\bigr)
:=\varinjlim_{K}H^1\bigl(\GL_2(\Q)\backslash(\GL_2(\Af)/K\times\HH^\pm),\C\bigr),
\]
把 $S_2$ 嵌入上同调。
\item 给定的同构 $\iota:\overline\Q_\ell\cong\C$ 给出上同调群的同构
\[
H^1\bigl(\GL_2(\Q)\backslash(\GL_2(\Af)\times\HH^\pm),\C\bigr)
\cong H^1\bigl(\GL_2(\Q)\backslash(\GL_2(\Af)\times\HH^\pm),\overline\Q_\ell\bigr).
\]
\item 如前所述，对（充分小的）紧开子群 $K\subseteq\GL_2(\Af)$，复流形
\[
\GL_2(\Q)\backslash(\GL_2(\Af)/K\times\HH^\pm)
\]
实际上是代数的，即由某个拟射影概形 $\widetilde X_K\to\Spec(\C)$ 的 $\widetilde X_K(\C)$ 给出。
\item 平展比较定理蕴涵
\[
H^1\bigl(\GL_2(\Q)\backslash(\GL_2(\Af)/K\times\HH^\pm),\overline\Q_\ell\bigr)
\cong H^1_{\text{ét}}(\widetilde X_K,\overline\Q_\ell).
\]
\item 现在奇迹发生了：对（充分小的）紧开子群 $K\subseteq\GL_2(\Af)$，拟射影概形 $\widetilde X_K\to\Spec(\C)$ 典范地定义在 $\Spec(\Q)$ 上！换言之，$\widetilde X_K\cong X_K\times_{\Spec(\Q)}\Spec(\C)$，其中 $X_K\to\Spec(\Q)$ 是（相容的）拟射影概形。
\item 特别地，由平展上同调在代数闭基域变换下的不变性，
\[
H^1_{\text{ét}}(\widetilde X_K,\overline\Q_\ell)\cong H^1_{\text{ét}}(X_{K,\overline\Q},\overline\Q_\ell).
\]
但现在右端带有 $G_\Q$ 的作用！
\item 现在看 $G_\Q\times\GL_2(\Af)$-模
\[
H^1_{\text{ét}}(X_{\overline\Q},\overline\Q_\ell)
:=\varinjlim_{K}H^1_{\text{ét}}(X_{K,\overline\Q},\overline\Q_\ell).
\]
\item 则我们将验证 $G_\Q$-模
\[
\mathrm{LL}(\pi):=\Hom_{\GL_2(\Af)}\bigl(\pi,H^1_{\text{ét}}(X_{\overline\Q},\overline\Q_\ell)\bigr)
\]
是二维的，并且可以用 Hecke 特征值表达 Frobenius 的迹。这个二维表示的对偶将是所要的表示。
\item 把常值层 $\overline\Q_\ell$ 换成某个局部系统，类似的策略对权 $k\ge2$ 也有效。\znote{原文 "Replacing $\overline\Q_\ell$ by some local system"：这里换掉的是系数层，即把常值层 $\overline\Q_\ell$ 换成与 $\omega^{\otimes k}$ 相应的局部系统（而不是把表示 $\sigma$ 换掉）；只有这样从 $H^1$ 中切出的才是权 $k$ 的表示。[来源：原文 p.053 与第 9 讲对 $\omega^{\otimes k}$ 的用法。]}
\end{itemize}

% ===== PDF p.054 =====
\begin{itemize}
\item 然而这对 $k=1$ 没有帮助。这里必须利用：权 $1$ 的模形式可以与某个权 $k\ge2$ 的模形式同余，然后使用前一种情形。
\end{itemize}

% 注：main.tex 中译注宏为 \znote / \znoteplain（任务书写作 \zhnote），此处按 main.tex 实际定义使用。
\chapter{模曲线作为椭圆曲线的模空间}

% ===== PDF p.055 =====
\noindent{\small\emph{（本讲由 Ben Heuer 撰写）}}

\textbf{上次}：
\begin{itemize}
\item 朗兰兹互反 $\mathrm{LL}\colon\cA_{\mathrm{mod}}\xrightarrow{\ \simeq\ }\mathcal{G}_{\mathrm{mod}}$ 的构造的初步草图：
\[
\cA_{\mathrm{mod}}:=\Bigl\{\pi\subseteq\bigoplus_{k\ge1}S_k\ \Bigm|\ \pi\ \text{是不可约的}\ \GL_2(\Af)\text{-子表示}\Bigr\},
\]
\[
\mathcal{G}_{\mathrm{mod}}:=\Bigl\{\sigma\colon G_\Q\to\GL_2(\Ql)\ \Bigm|\ \text{在有限集}\ S\ \text{之外非分歧、奇、在}\ \ell\ \text{上方 de Rham}\Bigr\}.
\]
\item 构造的关键点：对充分小的紧开子群 $K\subseteq\GL_2(\Af)$，复流形
\[
\GL_2(\Q)\backslash\bigl(\GL_2(\Af)/K\times\HH^{\pm}\bigr)
\]
(1) 在 $\Spec(\C)$ 上是代数的，(2) 有 $\Spec(\Q)$ 上的典范模型 $X_K$。由此得到 $\ell$-进平展上同调 $H^1_{\mathrm{\acute{e}t}}(X_{K,\overline{\Q}},\Q_\ell)$ 上的 Galois 作用。\znoteplain{「$K$ 充分小」是用来消去椭圆曲线的非平凡自同构的：它使商 $K\backslash\HH^{\pm}$ 成为光滑复曲线；再由 Baily--Borel 定理，这个商其实是拟射影代数簇，因而带有 $\Spec(\C)$ 上的代数结构。}
\end{itemize}

\textbf{今天}：
\begin{itemize}
\item 解释代数模型 $X_K$ 为何存在。
\item 把此前定义的概念用代数几何的语言重新解释：模形式、紧化、$q$-展开、adélic Hecke 作用，等等。
\end{itemize}

\medskip
\noindent\textbf{椭圆曲线的回顾}

\noindent 参考文献：[Silverman: The Arithmetic...]
\begin{itemize}
\item 设 $K$ 为任意域。椭圆曲线有如下相互等价的定义。
\end{itemize}

\noindent\textbf{定义 7.1}\phantomsection\label{def:7.1}\quad 域 $K$ 上的椭圆曲线等价地是下列对象之一：
\begin{itemize}
\item 亏格为 $1$ 的连通光滑射影曲线 $E|K$，连带一个选定的点 $O\in E(K)$；
\item $1$ 维连通光滑射影代数群 $E|K$；
\item 由 Weierstraß 方程
\[
E\colon y^2=x^3+ax+b,\qquad a,b\in K
\]
定义的非奇异平面三次曲线（若 $\operatorname{char}K\notin\{2,3\}$ 即为此式；否则需要更多项，参见 [Del75]）。
\end{itemize}

\begin{itemize}
\item \textbf{事实}：上述群概形自动是交换的（！）。
\item \textbf{事实}：非奇异性可以写成判别式的条件：
\[
E\ \text{非奇异}\iff\Delta(a,b)=-16(4a^3+27b^2)\ne0\qquad(\text{当}\ \operatorname{char}K\notin\{2,3\}).
\]
\item \textbf{事实}：给定 $E$，Weierstraß 方程并不唯一。但 $j$-不变量
\[
j(E):=j(a,b)=1728\,\frac{4a^3}{4a^3+27b^2}
\]
与 Weierstraß 方程的选取无关。
\item \textbf{事实}：若 $K$ 代数闭，则
\[
E\cong E'\iff j(E)=j(E').
\]
\item 更一般地，可以把 $\Spec(K)$ 换成任意基概形 $S$。
\end{itemize}

\noindent\textbf{定义 7.2}\phantomsection\label{def:7.2}\quad $S$ 上的椭圆曲线是指一个固有光滑曲线 $E\to S$，其几何纤维连通且亏格为 $1$，并带有一个点 $0\colon S\to E$。

% ===== PDF p.056 =====
\begin{itemize}
\item 可以把它看作以 $S$ 为参数的椭圆曲线族。
\item 它同样自动带有群结构：
\end{itemize}

\noindent\textbf{定理}（Abel）\quad 给椭圆曲线 $E\to S$ 赋予交换 $S$-群概形的结构、使得 $0$ 为单位截面，这样的方式唯一。

\begin{itemize}
\item 在 $S$ 上 Zariski 局部地，可以用一个 Weierstraß 方程描述 $E$。
\end{itemize}

\noindent\textbf{定义 7.3}\phantomsection\label{def:7.3}\quad 椭圆曲线之间的态射 $E\to E'$（在 $S$ 上）是指 $S$-群概形的同态。\textbf{同源\idx{同源}}是指 fppf-满射的同态，且其核有限平坦。

\noindent\textbf{例 7.4}\phantomsection\label{ex:7.4}\quad 对任意 $N\in\Z$，群概形结构给出「乘以 $N$」的态射，记作 $[N]\colon E\to E$。
\begin{itemize}
\item \textbf{事实}：$[N]$ 是秩为 $N^2$ 的有限平坦态射。若 $N$ 在 $S$ 上可逆，则它是有限平展的。
\item 特别地，$[N]$ 是同源。
\end{itemize}

\noindent\textbf{定义 7.5}\phantomsection\label{def:7.5}\quad 记 $E[N]:=\ker[N]\subseteq E$，即 $N$-挠子群。这是有限平坦的闭子群概形；当 $N$ 在 $S$ 上可逆时，它是有限平展的。
\begin{itemize}
\item 更一般地，对任意指数为 $N$ 的有限平坦子群概形 $D\subseteq E[N]$，存在唯一的椭圆曲线 $E/D$ 以及核为 $D$ 的同源 $\phi\colon E\to E/D$。
\item $[N]$ 诱导出对偶同源 $\phi^\vee\colon E/D\to E$，其核为 $E[N]/D$。
\end{itemize}

\medskip
\noindent\textbf{$K=\C$ 上的椭圆曲线}

\noindent 参考文献：[Diamond--Shurman: A first course...]
\begin{itemize}
\item 设 $E|\C$ 为椭圆曲线。则 $E(\C)$ 具有紧复 Lie 群的结构。作为这样的对象，它可以被典范地单值化：
\item 记 $\Lie E$ 为 $0\in E(\C)$ 处的切空间。则 $\Lie E\cong\C$，且有复 Lie 群的同构
\[
E(\C)=\Lie E/H_1(E,\Z)\cong\C/\Lambda,\qquad \Lambda\cong\Z^2 .
\]
\item 对任意 $N\in\Z$，有
\[
E(\C)[N]=\bigl(\tfrac1N\Lambda\bigr)/\Lambda\subseteq\C/\Lambda .
\]
\item \textbf{注}：$E(\C)[N]=\bigl(\tfrac1N\Z/\Z\bigr)^2=(\Z/N)^2$ 的秩为 $N^2$，与上述相符。
\item \textbf{要点}：$E(\C)$ 是 $1$ 维复环面：即 $\C$ 对一个格 $\Z^2\cong\Lambda\subseteq\C$（满足 $\Lambda\otimes_\Z\R=\C$）的商。
\item 复环面的态射 $:=$ 复 Lie 群的态射。
\end{itemize}

\noindent\textbf{定理}\quad 有范畴等价
\[
\{\C\ \text{上的椭圆曲线}\}\longrightarrow\{1\ \text{维复环面}\}.
\]
\begin{itemize}
\item 本质满性：设 $\Lambda\subseteq\C$ 是如上之格。定义 Weierstraß $\wp$ 函数
\[
\wp(z,\Lambda)=\frac1{z^2}+\sum_{w\in\Lambda,\ w\ne0}\Bigl(\frac1{(z-w)^2}-\frac1{w^2}\Bigr).
\]
\znoteplain{底本提取文本此处把括号内写作 $1/z^2-1/(z-w)^2$；求和时该式在 $|w|\to\infty$ 发散；Weierstraß $\wp$ 函数的标准定义是 $1/(z-w)^2-1/w^2$，减去主项后和式才对 $|w|\to\infty$ 收敛，译文从后者。}
\item 于是 $\wp(z)=\wp(z,\Lambda)$ 在 $\C\backslash\Lambda$ 上全纯，其复导数为
\[
\wp'(z)=-2\sum_{w\in\Lambda}\frac1{(z-w)^3}.
\]
\end{itemize}

% ===== PDF p.057 =====
\begin{itemize}
\item 它显然以 $\Lambda$ 为周期，$\rightsquigarrow$ $\wp$ 与 $\wp'$ 是 $\C/\Lambda$ 上的亚纯函数。
\end{itemize}

\noindent\textbf{定理}\quad $\wp$ 与 $\wp'$ 由一个 Weierstraß 方程相联系：
\[
E_\Lambda\colon \wp'^2=4(\wp^3+g_2(\Lambda)\wp+g_3(\Lambda)),
\]
其中 $g_2(\Lambda),g_3(\Lambda)\in\C$ 有显式表达式。于是得到 Lie 群的同构
\[
(\wp,\wp')\colon \C/\Lambda\xrightarrow{\ \sim\ }E_\Lambda(\C).
\]
\begin{itemize}
\item 把 $\C/\Lambda\mapsto E_\Lambda$ 送过去，定义了从 $\{1\ \text{维复环面}\}$ 到 $\{\C\ \text{上的椭圆曲线}\}$ 的拟逆。
\end{itemize}

\medskip
\noindent\textbf{椭圆曲线的复模空间}
\begin{itemize}
\item 现在我们可以把复椭圆曲线按同构加以参数化。
\item 为此只需确定何时两个格 $\Lambda_1,\Lambda_2$ 满足 $\C/\Lambda_1\cong\C/\Lambda_2$。
\item \textbf{事实}：复环面之间的任何同态都是乘以某个 $a\in\C$。事实上，它可以提升为 Lie 代数上的全纯同态。
\item 于是这发生当且仅当 $\Lambda_1=a\Lambda_2$ 对某个 $a\in\C^\times$ 成立。
\item 现在可以如下单值化：写 $\Lambda=w_1\Z\oplus w_2\Z$。
\item 条件 $\Lambda\otimes_\Z\R=\C$ 蕴含 $w_1/w_2\in\C\backslash\R=\HH^\pm$。
\item 从而对每个复环面 $E$，存在 $\tau\in\HH^\pm$ 使得
\[
E\cong E_\tau:=\C/\Lambda_\tau,\qquad \Lambda_\tau=\Z\oplus\tau\Z .
\]
\item 最后，对 $\tau,\tau'\in\HH^\pm$，有
\[
\Z+\tau\Z=\Z+\tau'\Z\iff\exists\,\gamma=\begin{pmatrix}a&b\\c&d\end{pmatrix}\in\GL_2(\Z):\ \gamma\tau=\frac{a\tau+b}{c\tau+d}=\tau'.
\]
\item 这表明：$\tau\mapsto E_\tau$ 定义了
\[
\GL_2(\Z)\backslash\HH^\pm\cong\{\text{复环面}\}/\!\sim\ \cong\{\C\ \text{上的椭圆曲线}\}/\!\sim .
\]
\item 这给我们早先的空间
\[
\SL_2(\Z)\backslash\HH=\GL_2(\Z)\backslash\HH^\pm=\GL_2(\Q)\backslash\bigl(\GL_2(\Af)/K\times\HH^\pm\bigr),\qquad K=\GL_2(\widehat{\Z})
\]
以「模空间」的解释（目前还是弱意义上的）：
\[
\SL_2(\Z)\backslash\HH=\{\C\ \text{上的椭圆曲线}\}/\!\sim .
\]
\item 由于 $\C$ 代数闭，$j$-不变量给出双射（事实上是双全纯）
\[
j\colon \SL_2(\Z)\backslash\HH\xrightarrow{\ \sim\ }\C .
\]
\item 通过覆叠映射 $\HH^\pm\to\GL_2(\Z)\backslash\HH^\pm$，由上述得到：
\[
\HH^\pm=\{\text{带有序基}\ \alpha\colon\Z^2\xrightarrow{\sim}\Lambda\ \text{的复环面}\}/\!\sim\ =\{\C\ \text{上带}\ \alpha\colon\Z^2\xrightarrow{\sim}H_1(E,\Z)\ \text{的椭圆曲线}\ E\}/\!\sim .
\]
\item 其他级的模解释呢？回忆对 $N\in\N$，我们曾有级 $K=K(N)$，
\[
X_K:=\GL_2(\Q)\backslash\bigl(\GL_2(\Af)/K\times\HH^\pm\bigr)=\coprod\Gamma(N)\backslash\HH^\pm,
\]
其中
\[
\Gamma(N)=\Bigl\{\gamma\in\GL_2(\Z)\ \Bigm|\ \gamma\equiv\begin{pmatrix}1&0\\0&1\end{pmatrix}\bmod N\Bigr\}.
\]
\end{itemize}

% ===== PDF p.058 =====
\begin{itemize}
\item 这些正是保持子群 $E_\Lambda[N]=\frac1N\Lambda/\Lambda\subseteq\C/\Lambda$ 的 $\Lambda$ 的自同构。
\item 于是有更一般的陈述
\[
\begin{aligned}
\Gamma(N)\backslash\HH^\pm&=\Bigl\{\text{带有序基}\ \alpha\colon(\Z/N)^2\xrightarrow{\sim}\tfrac1N\Lambda/\Lambda\\
&\text{的复环面}\Bigr\}/\!\sim\\
&=\Bigl\{\C\ \text{上带}\ \alpha\colon(\Z/N)^2\xrightarrow{\sim}E(\C)[N]\ \text{的椭圆曲线}\ E\Bigr\}/\!\sim .
\end{aligned}
\]
\item 在 $\GL_2(\Z)$ 与 $\Gamma(N)$ 之间，我们曾有群
\[
\begin{gathered}
\Gamma_1(N)=\Bigl\{\gamma\in\GL_2(\Z)\ \bigm|\ \gamma\equiv\left(\begin{smallmatrix}*&*\\[1.5pt]0&1\end{smallmatrix}\right)\bmod N\Bigr\},\\
\Gamma_0(N)=\Bigl\{\gamma\in\GL_2(\Z)\ \bigm|\ \gamma\equiv\left(\begin{smallmatrix}*&*\\[1.5pt]0&*\end{smallmatrix}\right)\bmod N\Bigr\}.
\end{gathered}
\]
\item 与它们相联系的模曲线有模解释：
\[
\Gamma_1(N)\backslash\HH^\pm=\{\C\ \text{上的椭圆曲线}\ E\ \text{连带一个恰为}\ N\ \text{阶的点}\ Q\in E(\C)[N]\}/\!\sim,
\]
\[
\Gamma_0(N)\backslash\HH^\pm=\{\C\ \text{上的椭圆曲线}\ E\ \text{连带一个秩为}\ N\ \text{的循环子群}\ C\subseteq E(\C)[N]\}/\!\sim .
\]
\end{itemize}

\medskip
\noindent\textbf{引导例：Legendre 族}
\begin{itemize}
\item 考虑变量 $\lambda$ 的复流形 $U=\Pone\backslash\{0,1,\infty\}$。
\item 在 $U$ 上，我们有一个复椭圆曲线族：
\item Legendre 族由 Weierstraß 方程
\[
E_\lambda\colon Y^2=X(X-1)(X-\lambda)
\]
从 $\mathbb{P}^2_U\to U$ 中割出。
\end{itemize}

\noindent\textbf{命题 7.6}\phantomsection\label{prop:7.6}
\begin{enumerate}
\item[(1)] 有 $E_\lambda[2]=\{O,(0,0),(1,0),(\lambda,0)\}$。特别地，有自然同构 $E_\lambda[2]=(\Z/2\Z)^2$。
\item[(2)] 对 $\C$ 上任意椭圆曲线 $E$ 连同平凡化 $\alpha\colon(\Z/2\Z)^2\xrightarrow{\sim}E[2]$，存在唯一的点 $x\in U$ 使得 $E=(E_\lambda)_x$。
\end{enumerate}

\begin{itemize}
\item \textbf{要点}：$E_\lambda\to U$ 在某种意义下是带 $\Gamma(2)$-级的万有椭圆曲线。
\item 这一切都是代数的！可以把 $E_\lambda\to U$ 理解为 $\C$ 上概形的态射。
\item 更好的是：这些概形已经定义在 $\Q$ 上。
\item 可以把这一论证细化，证明
\[
\Gamma(2)\backslash\HH^\pm=:X_{\Gamma(2)}(\C)=\Pone\backslash\{0,1,\infty\}.
\]
\item \textbf{注}（NB）：特别地，紧化为 $X^*_{\Gamma(2)}=\Pone$。
\item 这表明 $S_2(\Gamma_0(2))=0$，这是我们在 Fermat 大定理中用到的。\footnote{（原注 11）我们隐含地用了 $\omega^{\otimes2}\cong\Omega^1$。}
\end{itemize}

\noindent\textbf{思路}：从 $\C$ 过渡到一般 $\Q$-概形 $S$ 上椭圆曲线的模，以得到 $\Gamma(N)\backslash\HH^\pm$ 在 $\Q$ 上的代数模型。

\begin{itemize}
\item \textbf{问题}：在 $\Q$ 上的概形范畴上，是否存在一个概形表示函子
\[
P_{\SL_2(\Z)}\colon S\mapsto\{S\ \text{上的椭圆曲线}\}/\!\sim\ ?
\]
这将是一个 $\Q$ 上的概形 $X$，其 $\C$-点为 $\SL_2(\Z)\backslash\HH$。\znoteplain{原文把「不带级结构」的模函子记作 $P_{\SL_2(\Z)}$（底本提取文本作 PSL2(Z)，下标丢失）；这里 $P_{-}$ 是模函子的记号，并不是群 $\mathrm{PSL}_2(\Z)$。}
\end{itemize}

% ===== PDF p.059 =====
\begin{itemize}
\item \textbf{回答}：不存在。看出这一点的一个简便途径是扭（twist）现象：
\item 可以有在 $\Q$ 上不同构、而在某个二次扩域 $K$ 上同构的两条椭圆曲线 $E\not\cong E'$。特别地，$j(E)=j(E')$。
\item 显然，$X$ 的两个 $\Q$-点相同当且仅当它们在 $K$ 上相同。所以这样的 $X$ 不可能存在！
\item 上面的 Legendre 族有同样的问题：它并不表示
\[
S\mapsto\{S\ \text{上的椭圆曲线}\ E\ \text{连同}\ \alpha\colon(\Z/2\Z)^2\xrightarrow{\sim}E[2]\}/\!\sim,
\]
因为它没有把二次扭考虑进去。
\item \textbf{事实}：域 $K$ 上椭圆曲线 $E$ 的扭对应于 1-余圈 $G_{\overline{K}|K}\to\Aut(E)$。二次扭来自对合 $[-1]\colon E\to E\in\Aut(E)$。
\item \textbf{口号}：非平凡的自同构对可表性是不利的。\znoteplain{原文脚注（12）：更确切地说，出问题的是自同构群的非平坦性（由椭圆曲线的例外同构造成）。例如，在许多情形下，参数化线丛同构类的 Picard 函子是可表的；但在诸如固有、几何整纤维的假设下，线丛的自同构群都是 $\Gm$，从而是平坦的。}
\item \textbf{要点}：为了得到可表的函子，需要额外的数据。
\end{itemize}

\medskip
\noindent\textbf{椭圆曲线的模问题}

\noindent [Katz--Mazur: Arithmetic moduli...]
\begin{itemize}
\item 考虑 $\Z[\frac1N]$ 上概形上的模函子：
\[
P_{\Gamma_0(N)}\colon S\mapsto\{(E|S,\ C\subseteq E[N]\ \text{为秩}\ N\ \text{的循环子群概形})\}/\!\sim,
\]
\[
P_{\Gamma_1(N)}\colon S\mapsto\{(E|S,\ Q\in E[N](S)\ \text{为恰}\ N\ \text{阶的点})\}/\!\sim,
\]
\[
P_{\Gamma(N)}\colon S\mapsto\{(E|S,\ \alpha\colon(\Z/N\Z)^2\xrightarrow{\sim}E[N]\ \text{为群概形的同构})\}/\!\sim .
\]
\item 也可以构造「混合的模问题」，例如
\[
P_{\Gamma_1(N)\cap\Gamma_0(p)}:=P_{\Gamma_1(N)}\times_{P_{\SL_2(\Z)}}P_{\Gamma_0(p)}\colon S\mapsto\{(E,Q,C)\}/\!\sim .
\]
\item 关键点：设 $(E,C)\in P_{\Gamma_0(N)}(S)$。自同构 $[-1]\colon E\to E$ 把 $C$ 送到 $C$。特别地，$[-1]\in\Aut(E,C)$。
\item 相反，设 $(E,Q)\in P_{\Gamma_1(N)}(S)$，则 $[-1]\colon E\to E$ 把 $Q\in E[N]$ 送到 $-Q\in E[N]$；当 $N\ge3$ 时二者不同。特别地，$[-1]\notin\Aut(E,Q)$。\znoteplain{原文脚注（13）：而且没有例外同构。}
\item $P_{\Gamma(N)}(S)$ 的情形类似。这解释了 Legendre 族的问题：$N=2$ 太小，对付不了自同构 $[-1]$！我们需要 $N\ge3$。
\item \textbf{要点}：当 $N\ge3$ 时，模问题 $P_{\Gamma(N)}$ 与 $P_{\Gamma_1(N)}$ 是「刚性的」（$:=$ 没有非平凡自同构），而 $P_{\Gamma_0(N)}$ 与 $P_{\Gamma(2)}$ 则不然。
\end{itemize}

\noindent\textbf{定理}\quad 设 $N\ge3$，$p$ 为任意素数。
\begin{enumerate}
\item[(1)] 模问题 $P_{\Gamma(N)}$ 与 $P_{\Gamma_1(N)}$ 分别可表为 $\Z[\frac1N]$ 上的光滑仿射曲线 $X_{\Gamma(N)}$ 与 $X_{\Gamma_1(N)}$。
\item[(2)] 对任意 $n\in\N$，模问题 $P_{\Gamma(N)\cap\Gamma_0(p^n)}$ 可表为 $\Z[\frac1N]$ 上的平坦仿射曲线 $X_{\Gamma_1(N)\cap\Gamma_0(p^n)}$，且该曲线在 $\Z[\frac1{pN}]$ 上光滑。\znoteplain{原文此处把模问题 $P_{\Gamma(N)\cap\Gamma_0(p^n)}$ 的可表曲线记作 $X_{\Gamma_1(N)\cap\Gamma_0(p^n)}$，下标与模问题不一致；按上下文二者应同为 $\Gamma(N)\cap\Gamma_0(p^n)$。此处照录原文记号。}
\end{enumerate}

% ===== PDF p.060 =====
\begin{itemize}
\item \textbf{注}：为了表示 $P_{\SL_2(\Z)}$，可以从概形过渡到更大的叠范畴。得到「椭圆曲线的模叠」$\cM_{1,1}$。
\item 在各级结构 $\Gamma(N)\subseteq\Gamma'\subseteq\Gamma$ 之间，有遗忘态射 $X_{\Gamma'}\to X_\Gamma$。它们在 $\Z[\frac1N]$ 上都是有限平展的。
\item $\GL_2(\Z/N\Z)$ 从右边作用在 $P_{\Gamma(N)}$ 上，通过预复合 $\alpha\mapsto\alpha\circ\gamma$。
\item 等价地，这是通过用 $\gamma^\vee=\gamma^{-1}\det\gamma$ 预复合得到的左作用。
\item 这诱导出 $X_{\Gamma(N)}$ 上的 $\GL_2(\Z/N\Z)$-作用。
\item 类似地，$X_{\Gamma_1(N)}$ 上有自然的 $(\Z/N\Z)^\times$-作用。
\item 于是可以定义模曲线
\[
X_{\Gamma_0(N)}:=X_{\Gamma_1(N)}/(\Z/N\Z)^\times .
\]
\item 它不表示 $P_{\Gamma_0(N)}$，但「尽可能地接近」。
\item 特别地，对代数闭域 $K$，有 $X_{\Gamma_0(N)}(K)=P_{\Gamma_0(N)}(K)$。
\end{itemize}

\medskip
\noindent\textbf{紧化}
\begin{itemize}
\item 由 Weierstraß 方程给出的 $j$-不变量定义了有限平坦函数
\[
j\colon X\to\A^1 .
\]
\item 把 $X$ 经 $j\colon X\to\A^1\hookrightarrow\Pone$ 看成 $\Pone$ 的开子概形，再取它在函数域 $K(X)$ 中的正规化，可以定义
\[
j\colon X^*\to\Pone,
\]
它仍是有限平坦态射。把 $X^*$ 想成 $X$ 添上一个有限除子。
\item \textbf{事实}：$X^*$ 光滑且固有。
\item \textbf{理由}：可以用 $\Z[\frac1N][\![q]\!]$ 上的 Tate 曲线看出。
\item $X^*(\C)$ 就是前面提到的 $\Gamma\backslash\HH^\pm$ 的光滑紧化 $X^*_\C$。
\item \textbf{注}：$X^*$ 是「广义椭圆曲线」的模空间，参见 [DR73]。
\end{itemize}

\medskip
\noindent\textbf{几何模形式\idx{几何模形式}}

\noindent 参考文献：[Katz]，[Loeffler: lectures notes]
\begin{itemize}
\item 固定 $N\ge3$。令 $X=X_{\Gamma_1(N)}$。
\item 目标：把模形式几何地重新解释为 $X$ 上某些层的截面：
\end{itemize}

\noindent\textbf{定义 7.7}\phantomsection\label{def:7.7}\quad 令 $\omega:=e^*\Omega^1_{E|X}$，其中 $e\colon X\to E$ 是单位截面。
\begin{itemize}
\item 由于 $E\to X$ 是 $1$ 维光滑态射，这是一个线丛。
\item \textbf{事实}：它唯一地延拓为 $X^*$ 上的线丛 $\omega$。\textbf{理由}：万有椭圆曲线延拓为群概形 $E^*\to X^*$，取 $e^*\Omega^1_{E^*|X^*}$。
\end{itemize}

\noindent\textbf{命题 7.8}\phantomsection\label{prop:7.8}\quad $\omega^{\otimes k}$ 的解析化自然同构于 $X^*(\C)$ 上权 $k$ 的模形式层 $\omega_k$。

\noindent\textbf{证明略述}．回忆：$\HH^\pm$ 是带有 $\alpha\colon\Z^2\xrightarrow{\sim}H_1(E,\Z)$ 的 $E|\C$ 的模空间。沿 $\alpha(1,0)$ 积分定义了同构 $e^*\Omega^1_{E|\C}\xrightarrow{\sim}\C$。于是得到 $\omega$ 在 $\HH^\pm$ 上的典范平凡化。验证 $\GL_2(\Z)$-作用与之相符即可。\hfill$\square$

\begin{itemize}
\item 特别地，层 $\omega$ 已经有 $\Z[\frac1N]$ 上的代数模型。
\end{itemize}

% ===== PDF p.061 =====
\noindent\textbf{定义 7.9}\phantomsection\label{def:7.9}\quad 对任意 $\Z[\frac1N]$-代数 $R$，可以定义
\[
M_k(\Gamma_1(N),R)=\Gamma\bigl(X^*_{\Gamma_1(N)}\times_{\Spec(\Z[1/N])}\Spec(R),\ \omega^{\otimes k}\bigr),
\]
即系数在 $R$ 中的权 $k$ 模形式构成的 $R$-模。
\begin{itemize}
\item \textbf{事实}：它是有限自由的 $\Z[\frac1N]$-模。
\item \textbf{事实}：利用 Tate 曲线，可以得到 $R[[q]]$ 中的 $q$-展开。$\rightsquigarrow$ 由此可以定义尖点形式。
\item \textbf{事实}：若 $R,S$ 是平坦的 $\Z[\frac1N]$-代数（例如任何 $\Q$-代数），则
\[
M_k(\Gamma_1(N),R)\otimes_R S=M_k(\Gamma_1(N),S).
\]
\item \textbf{要点}：模形式显然是代数几何的对象！
\end{itemize}

% 注：main.tex 中译注宏为 \znote / \znoteplain（任务书写作 \zhnote），此处按 main.tex 实际定义使用。
\chapter{Eichler--Shimura 关系}

% ===== PDF p.062 =====
\noindent{\small\emph{（本讲由 Ben Heuer 撰写）}}

\textbf{上次}：
\begin{itemize}
\item 有 $\GL_2(\Z)$-等变的双射
\[
\HH^\pm\to\Bigl\{\C\ \text{上的椭圆曲线}\ E\ \text{连同}\ \alpha\colon\Z^2\xrightarrow{\sim}H_1(E,\Z)\Bigr\}/\!\sim,
\]
\[
\tau\mapsto\bigl(E_\tau:=\C/(\Z+\Z\tau),\ \alpha\colon\Z^2\to\Z+\Z\tau,\ (1,0)\mapsto1,\ (0,1)\mapsto\tau\bigr).
\]
\item \textbf{推论}：
\[
\GL_2(\Z)\backslash\bigl(\GL_2(\widehat{\Z})\times\HH^\pm\bigr)\cong\Bigl\{\C\ \text{上的椭圆曲线}\ E\ \text{连同}\ \alpha\colon\widehat{\Z}^2\xrightarrow{\sim}T_E\Bigr\}/\!\sim,
\]
其中 $T_E:=\varprojlim_{N\in\N}E[N]$ 是 $E$ 的完整 adélic Tate 模。
\item 这可以定义在 $\Q$ 上！
\item 令 $\widehat{X}\to\Spec(\Q)$ 为表示函子
\[
\Q\text{-}\mathbf{Sch}\to\mathbf{Sets},\qquad S\mapsto\Bigl\{S\ \text{上的椭圆曲线}\ E\ \text{连同}\ \alpha\colon\widehat{\Z}^2\xrightarrow{\sim}T_E\Bigr\}/\!\sim
\]
的模概形。这里 $T_E=\varprojlim_N E[N]$ 是 $S$ 上有限平展群概形的逆向极限。
\item $-1\in\GL_2(\widehat{\Z})$ 在 $\widehat{X}$ 上作用平凡。
\item 对满足 $K\subseteq K_1(3)$ 的紧开子群 $K\subseteq\GL_2(\widehat{\Z})$（注意特别地 $-1\notin K$），有光滑拟射影曲线 $X_K=\widehat{X}/K$。
\item 例如：若 $K=K(N)$、$N\ge3$，则 $X_{K(N)}$ 参数化椭圆曲线 $E$ 连同平凡化 $(\Z/N)^2\xrightarrow{\sim}E[N]$。
\item 其他例子：$K=K_1(N)$ 给出 $\Z/N\hookrightarrow E[N]$，$K=K_0(N)$ 给出子群。
\end{itemize}

\textbf{下一个目标}：
\begin{itemize}
\item 用模曲线几何地重新解释 Hecke 算子。
\item 由模曲线模 $p$ 的几何，导出 Galois 作用与 Hecke 作用之间的基本关系：Eichler--Shimura 关系\idx{Eichler--Shimura 关系}。
\end{itemize}

\textbf{adélic 作用}：
\begin{itemize}
\item 显然，$\GL_2(\widehat{\Z})$ 通过 $\alpha\colon\widehat{\Z}^2\xrightarrow{\sim}T_E$ 上的预复合作用在 $\widehat{X}$ 上。
\item 但是：为了得到 Hecke 算子，我们需要 $\GL_2(\Af)$-作用 $\Rightarrow$ 要在准同源的意义下处理椭圆曲线。
\item 回忆：$\Af=\widehat{\Z}\otimes_\Z\Q$，且
\[
\GL_2(\Z)\backslash\bigl(\GL_2(\widehat{\Z})\times\HH^\pm\bigr)\cong\GL_2(\Q)\backslash\bigl(\GL_2(\Af)\times\HH^\pm\bigr).
\]
\item 在椭圆曲线上可以这样解释这一点。\znoteplain{为什么要「到准同源为止」：$\GL_2(\Af)$ 中 $p$-进方向的分量会改变 $T_E$ 这个格，因而在 $E$ 的同构意义下不能直接作用；但在准同源意义下，有理 Tate 模 $VE=T_E\otimes_\Z\Q$ 只依赖于 $E$ 的同源类，于是 $\GL_2(\Af)$ 得以作用。}
\item 回忆左边几何化为 $\widehat{X}$，它表示函子
\[
P\colon\Q\text{-}\mathbf{Sch}\to\mathbf{Sets},\qquad S\mapsto\{(E|S,\ \alpha\colon\widehat{\Z}^2\xrightarrow{\sim}T_E)\}/\text{同构}.
\]
\item 令 $VE=T_E\otimes_\Z\Q$ 为有理 adélic Tate 模，并定义
\[
P'\colon\Q\text{-}\mathbf{Sch}\to\mathbf{Sets},\qquad S\mapsto\{(E|S,\ \alpha\colon\Af^2\xrightarrow{\sim}VE)\}/\text{准同源}.
\]
\znoteplain{原文此处有脚注 14，但底本的提取文本中没有该脚注的正文，暂付阙如，待对照原书补入。}
\end{itemize}

% ===== PDF p.063 =====
\noindent\textbf{定义 8.1}\phantomsection\label{def:8.1}\quad 回忆：同源是指具有有限平坦核的满同态 $E\to E'$。\textbf{准同源}是指
\[
\{E\to E'\ \text{同源}\}\otimes_\Z\Q
\]
中的元素。\textbf{$p$-准同源}是指
\[
\{E\to E'\ \text{同源，其次数为}\ p\ \text{的幂}\}\otimes_\Z\Z\bigl[\tfrac1p\bigr]
\]
中的元素。
\begin{itemize}
\item 于是 $P\cong P'$。确实：
\begin{itemize}
\item 显然有自然变换 $P\to P'$。逆方向的略述如下：
\item 取 $(E,\alpha\colon\Af^2\xrightarrow{\sim}VE)\in P'(S)$。乘以同源 $N\colon E\to E$，直到 $\alpha^{-1}$ 限制到 $T_E$ 上给出 $T_E\to\widehat{\Z}^2$。把两边模 $M\in\N$ 约化。令 $D$ 为诱导映射
\[
E[M]\to(\Z/M)^2
\]
的核；当 $M\gg0$ 时它稳定下来。对偶同源 $E/D\to E$ 诱导出 $\alpha\colon\widehat{\Z}^2\xrightarrow{\sim}T(E/D)$。\znoteplain{这里「乘以同源 $N$」的作用是消去分母：$\alpha^{-1}(T_E)$ 是 $VE$ 中与 $\widehat{\Z}^2$ 可公度的格，乘 $N$ 倍之后可使它落进 $\widehat{\Z}^2$ 内（相差一个有限指数）；随后 mod $M$ 的下降再把这个有限指数商掉。}
\item 这在 $P(S)$ 中定义了唯一的代表元。于是 $P(S)\twoheadrightarrow P'(S)$。
\end{itemize}
\item 这表明自然的 $\GL_2(\widehat{\Z})$-作用延拓为 $\GL_2(\Af)$-作用。
\item 但这不固定到任何 $X_K$ 上的投影，因为它改变了 $E$。
\item \textbf{例}：$\begin{pmatrix}p&0\\0&1\end{pmatrix}$ 的作用是送 $E\mapsto E/D$，其中 $D\subseteq E[p]$ 是由 $\alpha(1,0)$ 在 $T_pE\twoheadrightarrow E[p]$ 下的像所生成的子群。
\item \textbf{事实}：这一切在紧化之后仍然成立。
\end{itemize}

\medskip
\noindent\textbf{Hecke 算子的几何解释}
\begin{itemize}
\item 令 $V=C^\infty(\GL_2(\Af)\times\HH^\pm)$ 为光滑函数
\[
f\colon\GL_2(\Af)\times\HH^\pm\to\C,\qquad (g,z)\mapsto f(g,z)
\]
构成的光滑 $\GL_2(\Af)$-表示，即：$f$ 被 $\GL_2(\Af)$ 的某个紧开子群固定，且对任意 $g\in\GL_2(\Af)$，$z\mapsto f(g,z)$ 是光滑的。
\item 固定一个素数 $p$。
\item 令 $\phi\in\cH(\GL_2(\Q_p))$ 为 $p$ 处 Hecke 代数中的任意元素，即具有紧支集的局部常值函数 $\phi\colon\GL_2(\Q_p)\to\C$。
\item 回忆对 $f\in V$ 有
\[
\phi*f(g,z)=\int_{\GL_2(\Q_p)}\phi(h)f(gh,z)\,dh,
\]
其中我们在 $\GL_2(\Q_p)$ 上取使 $K_p=\GL_2(\Z_p)$ 的体积为 $1$ 的 Haar 测度。
\item 上面的积分可以改写如下。
\item 考虑图表
\[
\begin{tikzcd}[column sep=3.2em, row sep=2.4em]
& \GL_2(\Af)\times\HH^{\pm}\times\GL_2(\Q_p)\arrow[dl,"q_1"']\arrow[dr,"q_2"] &\\
\GL_2(\Af)\times\HH^{\pm} & & \GL_2(\Af)\times\HH^{\pm}
\end{tikzcd}
\]
\end{itemize}

% ===== PDF p.064 =====
\[
q_1(g,z,h)=(g,z),\qquad q_2(g,z,h)=(gh,z).
\]
\begin{itemize}
\item 这样的图表称为\textbf{对应}（correspondence）。
\item 还有到第三个因子的投影：
\[
\pi\colon\GL_2(\Af)\times\HH^\pm\times\GL_2(\Q_p)\to\GL_2(\Q_p).
\]
\end{itemize}

\noindent\textbf{引理 8.2}\phantomsection\label{lem:8.2}\quad 对 $f\colon\GL_2(\Af)\times\HH^\pm\to\C$（$f\in V$），有
\[
\phi*f=q_{1,!}\bigl(q_2^*(f)\cdot\pi^*(\phi)\bigr),
\]
其中 $(-)^*$ 表示函数的拉回，$(-)_!$ 表示沿纤维积分（由于 $\phi$ 有紧支集，这是有意义的）。

\noindent\textbf{证明}．$q_1$ 在 $(g,z)\in\GL_2(\Af)\times\HH^\pm$ 上的纤维是
\[
q_1^{-1}(g,z)=\{(g,z,h)\mid h\in\GL_2(\Q_p)\}\cong\GL_2(\Q_p),
\]
而 $q_2^*(f)\cdot\pi^*(\phi)$ 是函数 $(g,z,h)\mapsto f(gh,z)\phi(h)$。施加 $q_{1,!}$，即沿上述纤维积分，得到
\[
\int_{\GL_2(\Q_p)}\phi(h)f(gh,z)\,dh=\phi*f(g,z).
\]
注意这是良定义的，因为 $\phi$ 有紧支集。\hfill$\square$

\medskip
\noindent\textbf{过渡到有限级}
\begin{itemize}
\item 设 $N$ 与 $p$ 互素，并令 $K:=K_1(N)\subseteq\GL_2(\Af)$。则 $K_p=\GL_2(\Z_p)$。
\item 记 $K^p$ 为 $K$ 的 $p$-外部分，即 $K=K^pK_p$ 且 $K^p\cap K_p=1$。
\item 假设 $f$ 被 $K$ 固定，即 $f$ 由某个函数 $\tilde f\colon\GL_2(\Af)/K\times\HH^\pm\to\C$ 拉回得到。
\item 假设 $\phi\in\cH(\GL_2(\Q_p),K_p)$，即 $\phi$ 是 $K_p$-双不变的。
\item 于是 $\phi$ 是某个具有有限支集的函数 $\tilde\phi\colon K_p\backslash\GL_2(\Q_p)/K_p\to\C$ 的拉回。
\item 此时 $\phi*f$ 仍是 $K$-不变的，即它是某个函数 $\widetilde{\phi*f}\colon\GL_2(\Af)/K\times\HH^\pm\to\C$ 的拉回。
\item 于是在保持记号不变的情况下，可以把上面的图表改为
\[
\begin{tikzcd}[column sep=2.2em, row sep=2.8em]
& \bigl(\GL_2(\Af)/K^p\times\HH^{\pm}\bigr)\times_{K_p}\GL_2(\Q_p)/K_p\arrow[dl,"q_1"']\arrow[dr,"q_2"] &\\
\GL_2(\Af)/K\times\HH^{\pm} & & \GL_2(\Af)/K\times\HH^{\pm}
\end{tikzcd}
\]
这里 $-\times_{K_p}-$ 是第 3 节的收缩积（即对作用 $k\cdot(g,z,h):=(gk^{-1},z,kh)$ 取商），且
\[
q_2([g,z,h])=[gh,z],\qquad q_1([g,z,h])=[g,z],
\]
其中方括号表示等价类。
\end{itemize}

% ===== PDF p.065 =====
\begin{itemize}
\item 我们再次有投影
\[
\pi\colon\bigl(\GL_2(\Af)/K^p\times\HH^\pm\bigr)\times_{K_p}\GL_2(\Q_p)/K_p\longrightarrow K_p\backslash\GL_2(\Q_p)/K_p,\qquad [g,z,h]\mapsto[h].
\]
\item 接下来，在左边对 $\GL_2(\Q)$ 取商：回忆这给出复流形
\[
X_{\Gamma_1(N)}(\C)=\GL_2(\Q)\backslash\bigl(\GL_2(\Af)/K\times\HH^\pm\bigr).
\]
\item 然后需要把函数 $f$ 换成该流形上 $\omega_k$（$k\in\Z$）的截面。
\item 更精确地说，我们得到
\begin{equation}\tag{1}\label{eq:8.1}
\begin{tikzcd}[column sep=3.2em, row sep=2.8em]
& \begin{gathered}\GL_2(\Q)\backslash\bigl(\GL_2(\Af)/K^p\times\HH^{\pm}\bigr)\\\times_{K_p}\GL_2(\Q_p)/K_p\end{gathered}\arrow[dl,"q_1"']\arrow[dr,"q_2"] &\\
X_{\Gamma_1(N)}(\C) & & X_{\Gamma_1(N)}(\C)
\end{tikzcd}
\end{equation}
\item 并且对 $f\in H^0(X_{\Gamma_1(N)}(\C),\omega_k)$，即级为 $K_1(N)$ 的弱模形式，有
\[
\phi*f=q_{1,!}\bigl(q_2^*(f)\cdot\pi^*(\phi)\bigr).
\]
\item 为使之精确，要用到：有典范同构
\[
q_2^*(\omega_k)\cong q_1^*(\omega_k),
\]
因为 $\omega_k$ 是由 $\HH^\pm$ 上一个 $\GL_2(\Q)$-等变线丛的拉回来定义的。
\item 于是可以把 $q_2^*(f)\cdot\pi^*(\phi)\in q_2^*(\omega_k)$ 看作 $q_1^*(\omega_k)$ 的截面。
\item 于是可以沿纤维求和，得到截面
\[
q_{1,!}\bigl(q_2^*(f)\cdot\pi^*(\phi)\bigr)\in\Gamma\bigl(X_{\Gamma_1(N)}(\C),\omega_k\bigr).
\]
\item 这是模形式上 Hecke 作用的一个更几何的解释。
\end{itemize}

\medskip
\noindent\textbf{用模来解释}
\begin{itemize}
\item 在椭圆曲线上，式 (1) 对应一个 \textbf{Hecke 对应}：
\[
\begin{tikzcd}[column sep=4em, row sep=2.4em]
& \{(\iota\colon E_1\dashrightarrow E_2,\ \alpha)\}\arrow[dl,"q_1"']\arrow[dr,"q_2"] &\\
\{(E_1,\alpha)\} & & \{(E_2,\alpha')\}
\end{tikzcd}
\]
其中
\begin{itemize}
\item $(E_1,\alpha)$ 是 $\C$ 上的椭圆曲线带 $\Gamma_1(N)$-级结构 $\alpha\colon\Z/N\Z\hookrightarrow E_1[N]$；
\item 等价地，这是带一个恰为 $N$ 阶的点 $Q:=\alpha(1)\in E_1[N](\C)$ 的数据；
\item $(\iota\colon E_1\dashrightarrow E_2,\alpha)$ 是 $\C$ 上椭圆曲线 $E_1,E_2$ 连同 $E_1$ 上的 $\Gamma_1(N)$-级结构 $\alpha$ 以及一个 $p$-准同源 $\iota$ 的数据；
\item $q_1\colon(\iota\colon E_1\dashrightarrow E_2,\alpha)\mapsto(E_1,\alpha)$；
\item $q_2\colon(\iota\colon E_1\dashrightarrow E_2,\alpha)\mapsto(E_2,\alpha')$，其中 $\alpha'$ 是复合
\[
\Z/N\Z\xrightarrow{\ \alpha\ } E_1[N]\xrightarrow{\ \iota\ } E_2[N]
\]
（用到：由于 $(p,N)=1$，$\iota$ 在 $N$-挠部分上是同构）。等价地，它把 $Q$ 送到 $E_2[N]$ 中的像 $Q'$。\znoteplain{对应 $q_1,q_2$ 的用法：$q_2^*(f)\pi^*(\phi)$ 是中间空间上「带权」的函数，沿 $q_1$ 的纤维积分（纤维有限时就是求和）把它压回 $X_{\Gamma_1(N)}$；这正是 Hecke 算子 $\sum_{D}[E/D,Q+D]$ 的几何写法。}
\end{itemize}
\end{itemize}

% ===== PDF p.066 =====
\begin{itemize}
\item 映射 $\pi$ 把 $(\iota\colon E_1\dashrightarrow E_2,\alpha)$ 送到同构
\[
T_p(\iota)\colon T_pE_1\otimes_{\Z_p}\Q_p\xrightarrow{\sim}T_pE_2\otimes_{\Z_p}\Q_p
\]
的类。更具体地，任取 $T_pE_1$ 与 $T_pE_2$ 的基，则 $T_p(\iota)$ 由矩阵 $A\in\GL_2(\Q_p)$ 表示。$A$ 在 $K_p\backslash\GL_2(\Q_p)/K_p$ 中的类与基的选取无关。
\item \textbf{问题}：$p$-准同源构成的函子不可表（只能由归纳概形表示）。
\item 然而：我们主要关心由下列双陪集的特征函数给出的 $\phi$：
\[
K_p\begin{pmatrix}p&0\\0&p\end{pmatrix}K_p\ \longleftrightarrow\ \tilde S_p,\qquad
K_p\begin{pmatrix}p&0\\0&1\end{pmatrix}K_p\ \longleftrightarrow\ \tilde T_p .
\]
\item 这时定义积分中的项 $\phi(h)$ 的支集就在这个双陪集上。$\rightsquigarrow$ 于是 Hecke 对应的支集落在可由概形表示的子空间上！
\item 对 $\tilde S_p$，考虑形如
\[
([p]\colon E\to E,\ \alpha)
\]
的组 $(\iota\colon E_1\dashrightarrow E_2,\alpha)$，即 $E_1=E_2$ 且 $\iota\colon E_1\to E_2$ 是乘以 $[p]$。
\item 对 $\tilde T_p$，则需要考虑
\[
(\iota\colon E_1\to E_2,\ \alpha),
\]
其中 $\iota\colon E_1\to E_2$ 是次数为 $p$ 的同源。
\item 等价地，这是 $E_1$ 的一个秩为 $p$ 的循环子群 $\ker\iota$ 的数据。
\item \textbf{上次}：这些现在都可由概形表示，甚至可以在 $\Spec(\Q)$ 或适当的 $\Spec(\Z[1/N])$ 上表示！
\end{itemize}

\medskip
\noindent\textbf{模曲线上的 Hecke 作用——代数几何的观点}
\begin{itemize}
\item 现在我们过渡到有限级 $K=K_1(N)$，其中 $p\nmid N$。
\item 相应的模曲线 $X_{\Gamma_1(N)}$（在 $\Z[\frac1N]$ 上）表示 $\Z[\frac1N]$-概形 $S$ 上的对 $(E,Q)$：$S$ 上的椭圆曲线 $E$ 连同点 $Q\in E[N](S)$。
\item 我们从 $\langle p\rangle:=\tilde S_p$ 开始。由上述，只需把 $Q$ 乘以 $p$。
\item 为此，利用 $(\Z/N\Z)^\times$ 在 $X_{\Gamma_1(N)}$ 上的作用。
\item 由上述，$\tilde S_p$ 的对应现在成为
\[
\begin{tikzcd}[column sep=4em, row sep=2.4em]
& X_{\Gamma_1(N)}\arrow[dl,"\mathrm{id}"']\arrow[dr,"p\,\cdot"] &\\
X_{\Gamma_1(N)} & & X_{\Gamma_1(N)}
\end{tikzcd}
\]
\item \textbf{事实}：对任意 $a\in(\Z/N\Z)^\times$，有典范同构 $a^*\omega=\omega$。
\item 由于右边的映射是同构，沿纤维的积分是平凡的。
\item 一切都可以延拓到紧化上。
\item 综合上述，我们证明：
\end{itemize}

\noindent\textbf{命题 8.3}\phantomsection\label{prop:8.3}\quad 考虑由复合
\[
M_k(\Gamma_1(N),\Z[\tfrac1N])=H^0(X^*_{\Gamma_1(N)},\omega^{\otimes k})\xrightarrow{\ p^*\ }H^0(X^*_{\Gamma_1(N)},p^*\omega^{\otimes k})=M_k(\Gamma_1(N),\Z[\tfrac1N])
\]
定义的算子。则它基变换到 $\C$ 之后就是 Hecke 算子 $\tilde S_p$。

% ===== PDF p.067 =====
\begin{itemize}
\item 对 $\tilde T_p$，代替同源 $(\iota\colon E\to E',\alpha)$，我们参数化三元组 $(E,D,Q)$，其中 $D=\ker\iota$ 是秩为 $p$ 的有限循环子群。得到模函子的映射
\[
\pi_1\colon(E,D,Q)\mapsto(E,Q),\qquad \pi_2\colon(E,D,Q)\mapsto(E/D,Q'),
\]
\[
Q':=Q+D=\text{$Q$ 在}\ E[N]\to(E/D)[N]\ \text{下的像}.
\]
\item 它们诱导出 Hecke 对应
\[
\begin{tikzcd}[column sep=4em, row sep=2.4em]
& X_{\Gamma_0(p)\cap\Gamma_1(N)}\arrow[dl,"\pi_1"']\arrow[dr,"\pi_2"] &\\
X_{\Gamma_1(N)} & & X_{\Gamma_1(N)}
\end{tikzcd}
\]
\item \textbf{事实}：$\pi_1$ 与 $\pi_2$ 延拓到紧化 $X^*_{\Gamma_0(p)\cap\Gamma_1(N)}\to X^*_{\Gamma_1(N)}$。
\item \textbf{事实}：$\pi_1$ 与 $\pi_2$ 是次数为 $p+1$ 的有限平坦态射，在紧化上亦然。
\item \textbf{事实}：有典范同构 $\pi_1^*\omega=\pi_2^*\omega$（在 $\Spec(\Q)$ 上）。
\item 由于 $\pi_1$ 的纤维有限，Hecke 算子定义中出现的积分变成离散和，可以解释为迹映射
\[
\Tr_\pi\colon\pi_{1,!}\pi_1^*\omega\to\omega .
\]
\item 综合上述，我们证明：
\end{itemize}

\noindent\textbf{命题 8.4}\phantomsection\label{prop:8.4}\quad 把算子
\[
H^0(X^*_{\Gamma_1(N)},\omega^{\otimes k})\xrightarrow{\ \pi_2^*\ }H^0(X^*_{\Gamma_1(N)\cap\Gamma_0(p)},\omega^{\otimes k})\xrightarrow{\ \Tr_{\pi_1}\ }H^0(X^*_{\Gamma_1(N)},\omega^{\otimes k})
\]
基变换到 $\C$ 之后就是 Hecke 算子 $\tilde T_p$。

\begin{itemize}
\item \textbf{注}：这一切已经定义在 $\Q$ 上（甚至 $\Z[\frac1N]$ 上）！于是有
\end{itemize}

\noindent\textbf{推论 8.5}\phantomsection\label{cor:8.5}\quad $\tilde S_p,\tilde T_p$ 在 $M_k(\Gamma_1(N),\C)$ 上的特征值是代数的。

\begin{itemize}
\item 为给出 $\tilde T_p$ 更显式的描述，用除子重新解释。
\item 令 $\Div(X_{\Gamma_1(N)})$ 为 $X$ 上的除子集合（$\approx$ 点的形式和）。
\item 于是可以把 $\langle p\rangle:=\tilde S_p$ 与 $\tilde T_p$ 重新解释为 $\Div(X_{\Gamma_1(N)})$ 上的算子：
\[
\tilde S_p\colon\Div(X_{\Gamma_1(N)})\to\Div(X_{\Gamma_1(N)}),\qquad [E,Q]\mapsto[E,p\cdot Q],
\]
\[
\tilde T_p\colon\Div(X_{\Gamma_1(N)})\to\Div(X_{\Gamma_1(N)}),\qquad [E,Q]\mapsto\sum_{D\subseteq E[p]}[E/D,Q+D],
\]
其中 $E$ 是椭圆曲线，$Q\in E[N]$ 是 $N$ 阶的点。
\end{itemize}

\medskip
\noindent\textbf{特征 $p$ 中的模曲线}
\begin{itemize}
\item 设 $E$ 是特征 $p$ 的概形 $S$ 上的椭圆曲线。
\item 则当 $p\nmid N$ 时，$E[N]$ 是平展的。
\item 态射 $[p]\colon E\to E$ 分解为
\[
E\xrightarrow{\ F\ }E^{(p)}\xrightarrow{\ V\ }E,
\]
即 Frobenius 同源与 Verschiebung 同源。这里 $E^{(p)}=E\times_{S,F}S$。
\item 于是总有 $\ker F\subseteq E[p]$。它是连通的，即 $\ker F(S)=0$。
\end{itemize}

\noindent\textbf{定义 8.6}\phantomsection\label{def:8.6}\quad 在 $S=\F_p$ 上，有两种情形：

% ===== PDF p.068 =====
\begin{enumerate}
\item[(1)] $E[p](\Fq)\cong\Z/p\Z$。此时 $E$ 称为\textbf{普通的}（ordinary），且 $E[p]=\mu_p\times\Z/p\Z$。
\item[(2)] $E[p](\Fq)=0$。此时 $E$ 称为\textbf{超奇异的}（supersingular），且 $E[p]$ 连通。
\end{enumerate}
\znote{原文此定义（1）(2) 两条作 $E(\F_p)=\Z/p\Z$、$E(\F_p)=0$，按 $E[p](\Fq)$ 译：普通时 $E[p]\cong\mu_p\times\Z/p\Z$，超奇异时 $E[p]$ 连通。[来源：Goren, \emph{Hilbert Modular Forms}, Example 3.10，\url{https://www.wstein.org/people/goren/hilbert_modular_forms_book.pdf}（2026-11-26 访问）。]}

\begin{itemize}
\item 「超奇异」说的并不是光滑性。它只意味着「非常特殊」：
\item \textbf{事实}：$\Fq$ 上超奇异椭圆曲线只有有限多个同构类。其余都是普通的。
\item 回忆：$X_{\Gamma_1(N)}$ 定义在 $\Z[\frac1N]$ 上。
\item 因此可以约化到 $\F_p$，得到模曲线 $X_{\Gamma_1(N),\F_p}\to\Spec(\F_p)$。
\item 类似地，$X_{\Gamma_1(N)\cap\Gamma_0(p)}$ 表示 $(E,D,Q)$，其中 $D\subseteq E[p]$ 秩为 $p$。
\item Deligne--Rapoport：纤维 $X_{\Gamma_1(N)\cap\Gamma_0(p),\F_p}\to\Spec(\F_p)$ 具有如下形状
\[
\begin{tikzcd}[column sep=4.6em, row sep=2.6em]
& X_{\Gamma_1(N)\cap\Gamma_0(p),\F_p}\arrow[dl,"\mathrm{id}"']\arrow[dr,"F"] &\\
X_{\Gamma_1(N),\F_p}\ (D=\mu_p) & & X_{\Gamma_1(N),\F_p}\ (D=\Z/p\Z)
\end{tikzcd}
\]
\item 即两个 $X_{\Gamma_1(N),\F_p}$ 的拷贝，分别对应 $D=\mu_p$ 与 $D=\Z/p\Z$ 作为 $E[p]$ 的子群，它们在超奇异点处横截相交（图中记为黑点）。
\item 投影在一个拷贝上是恒等（$\deg=1$），在另一个上是 Frobenius（$\deg=p$）。
\end{itemize}

\medskip
\noindent\textbf{Eichler--Shimura 关系}

\noindent [Diamond--Shurman]，[Conrad: Appendix to Serre's...]
\begin{itemize}
\item Eichler--Shimura 关系把 $\tilde T_p$ 模 $p$ 的约化用 Frobenius 表达出来：
\item 基变换到代数闭域 $\Fq$。
\item 考虑态射 $\Div^0(X^*_{\Gamma_1(N),\Fq})\to\Pic^0(X^*_{\Gamma_1(N),\Fq})$。
\item \textbf{事实}：$\Pic^0(X^*_{\Gamma_1(N),\Fq})$ 是 $\Fq$ 上某个群概形的点。特别地，乘以 $p$ 通过 Frobenius $F$ 分解：存在 $V=F^\vee$ 使得 $p=V\circ F=F\circ V$。
\end{itemize}

\noindent\textbf{定理}（Eichler--Shimura 关系）\quad 在 $\Pic^0(X^*_{\Gamma_1(N),\Fq})$ 中，有
\[
\tilde T_p=F+\langle p\rangle V .
\]
\begin{itemize}
\item \textbf{注}：对 $X^*_{\Gamma_0(N),\Fq}$ 也有一个版本，众所周知它写作
\[
\tilde T_p=F+V .
\]
把上式乘以 $F$ 并用 $FF^\vee=p$，得到第二个重要的变体
\[
F^2-\tilde T_pF+\langle p\rangle p=0 .
\]
\znoteplain{这一变体的来历：由 $\tilde T_p=F+\langle p\rangle V$ 两边右乘 $F$，得 $\tilde T_pF=F^2+\langle p\rangle VF$；再用 $VF=p$ 即得上式。第 9 讲构造 Galois 表示时，特征多项式 $X^2-a_pX+\chi(p)p$ 正来源于此。}
\end{itemize}

\noindent\textbf{证明}（略述）．设 $E$ 是 $\breve{\Z}_p=W(\Fq)$ 上的普通椭圆曲线。令 $C\subseteq E[p]$ 为「典范子群」：它由
\[
E[p](\breve{\Z}_p)\to E[p](\Fq)
\]
的核生成。

% ===== PDF p.069 =====
\begin{itemize}
\item \textbf{关键事实}：对秩为 $p$ 的循环子群概形 $D\subseteq E[p]$，
\begin{enumerate}
\item[(1)] 同源 $E\to E/D$ 约化为 $F\colon E\to E^{(p)}$ $\iff$ $D=C$；
\item[(2)] 同源 $E\to E/D$ 约化为 $V\colon E\to E^{(p-1)}$ $\iff$ $D\ne C$。
\end{enumerate}
\item 这里 $E^{(p-1)}:=E\times_{\F_p,F^{-1}}\F_p$，即沿 Frobenius 的逆作基变换。
\item $F$ 把 $Q\in E[N]$ 送到基变换 $Q^{(p)}\in E^{(p)}[N]=E[N]^{(p)}$。
\item 由于 $V\circ F=p$ 作用在 $E[N]$ 上，这蕴含：$V$ 把 $Q$ 送到 $pQ^{(p-1)}\in E^{(p-1)}[N]$。
\item 回忆：$E[p]$ 在 $W(\Fq)$ 上有 $p+1$ 个不同的子群 $D$。
\item 因此 $\tilde T_p([E,Q])=\sum_{D\subseteq E[p]}[E/D,Q+D]$（在 $\Div(X_{\Gamma_1(N)})$ 中）约化为
\[
[E^{(p)},Q^{(p)}]+p[E^{(p-1)},pQ^{(p-1)}]
\]
（在 $\Div(X_{\Gamma_1(N),\Fq})$ 中）。
\item 第一个求和项是 $F[E,Q]$。
\item 按定义，$\langle p\rangle[E^{(p-1)},Q^{(p-1)}]=[E^{(p-1)},pQ^{(p-1)}]$。
\item 现在过渡到 $\Pic^0(X^*_{\Gamma_1(N),\Fq})$。这里 $p=VF$。于是有
\[
p[E^{(p-1)},Q^{(p-1)}]=pF^{-1}[E,Q]=V[E,Q].
\]
（严格地说，需要取 $[E,Q]-[E',Q']$ 才能落在 $\Div^0$ 中。）\hfill$\square$
\end{itemize}

\chapter{附着于新形式的 Galois 表示}

% ===== PDF p.070 =====
\noindent\textbf{下一个目标}：
\begin{itemize}
\item 在上同调中实现新形式。
\item 完成「附着于权 $k=2$ 的新形式的 Galois 表示」这一构造。
\item 对 $k\ge2$ 的情形给出如何继续的提示。
\end{itemize}

我们将借助微分形式，把（权 $2$ 的）模形式对应到奇异上同调中的类。为此需要回忆 de Rham 上同调的一些事实。

\noindent\textbf{de Rham 上同调}（参见 \texttt{[BT13]}）：
\begin{itemize}
\item 设 $X$ 为实流形。
\item 令 $\C^\infty(X)$ 为光滑函数 $f\colon X\to\C$ 构成的空间。
\item 令 $A^i(X)$ 为 $X$ 上光滑 $i$-形式构成的空间，$i\ge0$。
\item $X$ 的 \textbf{de Rham 复形}是复形
\[
A^\bullet(X)\colon\quad A^0(X)=\C^\infty(X)\xrightarrow{\ \mathrm{d}\ }A^1(X)\xrightarrow{\ \mathrm{d}\ }A^2(X)\to\cdots
\]
其中 $\mathrm{d}$ 为外微分。
\item $X$ 的 de Rham 上同调则定义为 de Rham 复形的上同调
\[
H^*_{\mathrm{dR}}(X):=H^*(A^\bullet(X)).
\]
\item 引人注目的是，de Rham 比较同构蕴涵
\[
H^*_{\mathrm{dR}}(X)\cong H^*(X,\C),
\]
右边是 $X$ 上常层 $\C$ 的层上同调。这有个有趣的推论：de Rham 上同调是 $X$ 的纯拓扑不变量。
\begin{itemize}
\item 证明略述：把开集 $U\subseteq X$ 送到 $A^\bullet(U)$，给出 $X$ 上的层复形 $\mathcal A^\bullet_X$，它是常层 $\C$ 的一个消解（由 Poincaré 引理）。用单位分解的论证（需要 $X$ 仿紧）可知，层 $\mathcal A^k_X$（$k\ge0$）都是 flasque（松弛）层。于是 $R\Gamma(X,\C)\cong\Gamma(X,\mathcal A^\bullet_X)$。
\end{itemize}
\item 此外，
\[
H^*(X,\C)\cong H^*_{\mathrm{sing}}(X,\C),
\]
右边是 $X$ 的以 $\C$ 为系数的奇异上同调。
\item 若离散群 $\Gamma$ 真不连续且自由地作用在 $X$ 上，则
\[
A^\bullet(\Gamma\backslash X)\cong A^\bullet(X)^\Gamma=\bigl(A^0(X)^\Gamma=\C^\infty(X)^\Gamma\xrightarrow{\ \mathrm{d}\ }A^1(X)^\Gamma\xrightarrow{\ \mathrm{d}\ }A^2(X)^\Gamma\to\cdots\bigr),
\]
也就是说，de Rham 复形中的不变量计算商 $\Gamma\backslash X$ 的上同调。注意：由于群上同调可能出现，这一陈述若放到上同调群层面便不成立。

% ===== PDF p.071 =====
\item 最后，我们令
\[
H^*_{\mathrm{dR}}\bigl(\GL_2(\Q)\backslash(\GL_2(\Af)\times\HH^\pm)\bigr):=\varinjlim_{K}H^*_{\mathrm{dR}}\bigl(\GL_2(\Q)\backslash(\GL_2(\Af)/K\times\HH^\pm)\bigr)
\]
以及
\[
H^*\bigl(\GL_2(\Q)\backslash(\GL_2(\Af)\times\HH^\pm),\C\bigr):=\varinjlim_{K}H^*\bigl(\GL_2(\Q)\backslash(\GL_2(\Af)/K\times\HH^\pm),\C\bigr),
\]
其中余极限取遍全部紧开子群 $K\subseteq\GL_2(\Af)$，过渡映射由拉回诱导。
\item 于是对 $i\ge0$，取余极限即得（$\GL_2(\Af)$-等变的）de Rham 比较
\[
H^i_{\mathrm{dR}}\bigl(\GL_2(\Q)\backslash(\GL_2(\Af)\times\HH^\pm)\bigr)\cong H^i\bigl(\GL_2(\Q)\backslash(\GL_2(\Af)\times\HH^\pm),\C\bigr).
\]
\znote{这里取余极限得到的大空间就是第 8 讲的 $\widehat X(\C)=\GL_2(\Q)\backslash(\GL_2(\Af)\times\HH^\pm)$，所以上面两条定义左边正是概形 $\widehat X$ 的上同调。下文把 $\C$ 换成 $\Q_\ell$、把 de Rham 上同调换成平展上同调，正是为了让 Galois 群在它上面作用。[来源：本讲义第 8 讲（PDF p.063）对 $\widehat X(\C)$ 的刻画。]}
\end{itemize}

\noindent\textbf{要点}：
\begin{itemize}
\item 可以用（闭的）微分形式构造 $H^*(X,\C)$ 中的类。
\end{itemize}

\noindent\textbf{附着于权 2 模形式的上同调类}：
\begin{itemize}
\item 注意
\[
\GL_2(\Q)\backslash(\GL_2(\Af)\times\HH^\pm)\cong\GL_2(\Q)^+\backslash(\GL_2(\Af)\times\HH),
\]
其中 $\GL_2(\Q)^+\subseteq\GL_2(\Q)$ 是行列式为正的元素构成的子群。
\item 取 $f\in H^0\bigl(\GL_2(\Q)^+\backslash(\GL_2(\Af)\times\HH),\omega^{\otimes k}\bigr)$，$k\in\Z$。
\item 我们把 $f$ 看作函数
\[
f\colon\GL_2(\Af)\times\HH\to\C,\qquad (g,z)\mapsto f(g,z),
\]
它满足
\[
f(\gamma g,\gamma z)=(cz+d)^kf(g,z),\qquad \gamma=\begin{pmatrix}a&b\\c&d\end{pmatrix}\in\GL_2(\Q)^+ .
\]
\item 假设 $f$ 的权为 $k=2$。
\item 此时 $\GL_2(\Af)\times\HH$ 上的微分形式 $f(g,z)\,\mathrm{d}z$ 是闭的，并且对 $\gamma\in\GL_2(\Q)^+$ 满足
\[
\gamma^*\bigl(f(g,z)\,\mathrm{d}z\bigr)=f(\gamma g,\gamma z)\,\gamma^*\mathrm{d}z=\det(\gamma)\,f(g,z)\,\mathrm{d}z .
\]
\item 确实：
\begin{itemize}
\item 闭性由全纯性得到，因为
\[
\mathrm{d}\bigl(f(g,z)\,\mathrm{d}z\bigr)=\frac{\partial f}{\partial z}\,\mathrm{d}z\wedge\mathrm{d}z-\frac{\partial f}{\partial\bar z}\,\mathrm{d}z\wedge\mathrm{d}\bar z=0 .
\]
\item 若 $\gamma=\begin{pmatrix}a&b\\c&d\end{pmatrix}\in\GL_2(\Q)$，则
\[
\gamma^*(\mathrm{d}z)=\frac{\det(\gamma)}{(cz+d)^2}\,\mathrm{d}z .
\]
\item 利用 $f$ 的权为 $2$，即得结论。
\end{itemize}

% ===== PDF p.072 =====
\item 令
\[
|\cdot|=|\cdot|_{\mathrm{ad\acute{e}lic}}\colon\Q^\times\backslash\A^\times\to\R_{>0}
\]
为 adélic 绝对值，即
\[
|(x_2,x_3,\dots,x_\infty)|:=\prod_p|x_p|_p\cdot|x_\infty|_\infty,
\]
其中 $|\cdot|_p\colon\Q_p\to\R_{\ge0}$，$x\mapsto p^{-v_p(x)}$ 是 $p$ 进绝对值，$|\cdot|_\infty\colon\R\to\R_{\ge0}$ 是实绝对值。
\item $|\cdot|$ 定义了 $\GL_2(\A)$ 的特征
\[
\chi_{no}:=|\cdot|\circ\det=|\det|_{\mathrm{ad\acute{e}lic}},
\]
它在 $\GL_2(\Q)$ 上平凡。
\item 于是对
\[
\tilde f(g,z):=\chi_{no}((g,1))f(g,z)
\]
（其中 $(g,1)\in\GL_2(\A)\cong\GL_2(\Af)\times\GL_2(\R)$），$\GL_2(\Af)\times\HH$ 上的形式
\[
\eta_f:=\tilde f(g,z)\,\mathrm{d}z
\]
是 $\GL_2(\Q)^+$-等变的。
\item 我们由此得到映射
\[
\alpha\colon M_2\longrightarrow H^1_{\mathrm{dR}}\bigl(\GL_2(\Q)^+\backslash(\GL_2(\Af)\times\HH)\bigr)\cong H^1\bigl(\GL_2(\Q)^+\backslash(\GL_2(\Af)\times\HH),\C\bigr),
\]
它把 $f$ 送到 $[\eta_f]$；这个映射不是 $\GL_2(\Af)$-等变的。
\item 对光滑 $\GL_2(\Af)$-表示 $V$ 与光滑特征 $\chi\colon\GL_2(\Af)\to\C^\times$，定义
\[
V(\chi):=V\otimes_\C\chi .
\]
\znoteplain{这里的扭是表示层面的一维扭：作为向量空间 $V(\chi)=V$，而 $g\in\GL_2(\Af)$ 的作用改为「先按 $V$ 作用、再乘 $\chi(g)$」。本讲此后反复用到这一记号，例如 $M_2(\chi_{no})$、$\pi(\chi_{no})$。}
\item 于是
\[
\alpha\colon M_2(\chi_{no})\longrightarrow H^1\bigl(\GL_2(\Q)^+\backslash(\GL_2(\Af)\times\HH),\C\bigr)
\]
是 $\GL_2(\Af)$-等变的。
\item 设 $K\subseteq\GL_2(\Af)$ 为紧开子群，并令
\[
X_K:=\GL_2(\Q)^+\backslash(\GL_2(\Af)/K\times\HH),
\]
其典范紧化为
\[
X^*_K:=\GL_2(\Q)^+\backslash(\GL_2(\Af)/K\times\HH^*),
\]
其中 $\HH^*:=\HH\cup\Pone(\Q)$（带 Satake 拓扑）。
\item 若 $f\in S_2(K)$，则 $\eta_f$ 延拓为 $X^*_K$ 上的全纯微分形式。
\begin{itemize}
\item 确实：若 $q=e^{2\pi iz}$，则
\[
\mathrm{d}q=2\pi iq\,\mathrm{d}z,\qquad\text{即}\qquad \mathrm{d}z=\frac{1}{2\pi iq}\,\mathrm{d}q .
\]
于是把 $f(g,z)$ 在每个尖点处写成 Fourier 展开，即写成 $q$ 的函数。
\end{itemize}

% ===== PDF p.073 =====
\item 下图交换：
\[
\begin{tikzcd}
H^1_{\mathrm{dR}}(X^*_K)\arrow[r,"\simeq"]\arrow[d,"\simeq"] & H^1_{\mathrm{dR}}(X_K)\arrow[d,"\simeq"]\\
H^1(X^*_K,\C)\arrow[r,"\simeq"] & H^1(X_K,\C)
\end{tikzcd}
\]
\item 此外，态射
\[
H^1_c(X_K,\C)\longrightarrow H^1(X^*_K,\C)=H^1_c(X^*_K,\C)
\]
是满射，因为其余核嵌入 $H^1(\{\text{尖点}\},\C)=0$。
\item 于是 $\alpha(S_2(K))$ 落在 $X_K$ 的内上同调
\[
\tilde H^1(X_K,\C)\cong H^1(X^*_K,\C)
\]
中；按定义，后者是紧支上同调 $H^1_c(X_K,\C)\to H^1(X_K,\C)$ 的像。
\item 注意上述讨论对反全纯模形式（即模形式的复共轭）同样适用。
\item 由此得到映射
\[
\bar\alpha\colon\bar S_2(K)\longrightarrow\tilde H^1(X_K,\C),\qquad f\mapsto[\chi_{no}f\,\mathrm{d}\bar z],
\]
其中 $\bar S_2(K)$ 表示反全纯模形式构成的 $\C$-线性空间。
\item 我们可以过渡到无穷级，得到态射
\[
\bar\alpha\oplus\alpha\colon\bar S_2\oplus S_2\longrightarrow\tilde H^1\bigl(\GL_2(\Q)\backslash(\GL_2(\Af)\times\HH^\pm),\C\bigr),
\]
其记号（希望）自明。
\end{itemize}

\noindent\textbf{定理 9.1}（Eichler--Shimura）\phantomsection\label{thm:9.1}\quad 映射
\[
\alpha\oplus\bar\alpha\colon S_2(\chi_{no})\oplus\bar S_2(\chi_{no})\longrightarrow\tilde H^1\bigl(\GL_2(\Q)\backslash(\GL_2(\Af)\times\HH^\pm),\C\bigr)
\]
是 $\GL_2(\Af)$-等变同构。

\begin{itemize}
\item 该陈述等价于它对所有（充分小的）紧开子群 $K\subseteq\GL_2(\Af)$ 的类似陈述。
\item 于是固定某个紧开子群 $K\subseteq\GL_2(\Af)$。
\item 证明将利用 $\cup$-积配对
\[
H^1_c(X_K,\C)\times H^1(X_K,\C)\longrightarrow H^2_c(X_K,\C)\xrightarrow{\ \int\ }\C,
\]
在 de Rham 比较之下，它由微分形式的配对
\[
(\eta_1,\eta_2)\mapsto\int_{X_K}\eta_1\wedge\eta_2
\]
诱导，其中 $\eta_1$ 具有紧支集（参见 \texttt{[BT13]}）。
\item 我们将用到如下观察：
\begin{itemize}
\item 紧支上同调可由 $H^*_c(X_K,\C)\cong H^*(A^\bullet_c(X_K))$ 计算，其中 $A^\bullet_c$ 表示紧支集微分形式。
\item 典范映射
\begin{equation}\tag{2}\label{eq:9.2}
A^\bullet_c(X_K)\longrightarrow A^\bullet_{\mathrm{rd}}(X_K)
\end{equation}
是拟同构，右边表示速降微分形式，参见 \texttt{[Bor80, Theorem 2]}。
\end{itemize}
\item 对 $f\in S_2$，形式 $\eta_f$ 速降。
\begin{itemize}
% ===== PDF p.074 =====
\item 利用 $|e^{2\pi inz}|=e^{-2\pi n\operatorname{Im}(z)}$ 当 $\operatorname{Im}(z)\to\infty$ 且 $n\ge1$ 时迅速衰减。
\end{itemize}
\item 于是 $\alpha$ 自然提升为映射
\[
\tilde\alpha\colon S_2(K)\longrightarrow H^1_c(X_K,\C).
\]
\item 取 $f\in S_2(K)$、$g\in\bar S_2(K)$，则：
\begin{itemize}
\item $\tilde\alpha(f)\cup\alpha(f)=0$，因为 $\mathrm{d}z\wedge\mathrm{d}z=0$。
\item 对 $g$ 同理。
\item 用 Stokes 定理可证
\[
\operatorname{tr}\bigl(\tilde\alpha(f)\cup\bar\alpha(g)\bigr)=\int_{X_K}\eta_f\wedge\eta_g .
\]
\item 由此可知，Poincaré 配对（相差一个标量）诱导出 $S_2$ 上的 Petersson 标量积，又称 $L^2$-配对 $\langle-,-\rangle$。
\end{itemize}
\item 现在可以证明 $\beta:=\alpha\oplus\bar\alpha$ 是单射。
\item 确实：
\begin{itemize}
\item 若 $\beta(f+g)=0$，则
\[
\beta(f+g)\cup\alpha(f)=\langle f,f\rangle=0,
\]
即 $f=0$。对 $g$ 同理。
\end{itemize}
\item 现在验证 $\tilde h^1(X_K,\C)=2\dim S_2(K)$。一旦成功，Eichler--Shimura 同构\idx{Eichler--Shimura 同构}的证明便完成了。
\item 即（我们用小写字母表示维数）：
\begin{itemize}
\item $\tilde h^1(X_K,\C)=h^1(X^*_K,\C)=2h^0(X^*_K,\C)-\chi_{\mathrm{top}}(X^*_K)$。
\item 这里 $\chi_{\mathrm{top}}$ 是拓扑 Euler 特征标。
\item 另一方面\znoteplain{原文脚注 15：我们隐含地使用了 $\omega^{\otimes2}(-\text{尖点})\cong\Omega^1$。原文作 $\omega^2$，按第 7 讲的记号应理解为 $\omega^{\otimes2}$。}：
\[
\begin{aligned}
\dim S_2(K)&=h^0\bigl(X^*_K,\Omega^1_{X^*_K}\bigr)\\
&\overset{\mathrm{RR}}{=}\deg\bigl(\Omega^1_{X^*_K}\bigr)+\chi_{\mathrm{hol}}(X^*_K)+h^1\bigl(X^*_K,\Omega^1_{X^*_K}\bigr)\\
&=-\chi_{\mathrm{hol}}(X^*_K)+h^0(X^*_K,\C).
\end{aligned}
\]
\znote{中间两步的来历：对曲线 $X^*_K$ 上的 $\Omega^1$ 用 Riemann--Roch 即得第一个等号后的三项；最后一步用 $\deg\Omega^1_{X^*_K}=2g-2=-2\chi_{\mathrm{hol}}(X^*_K)$ 与 Serre 对偶 $h^1(X^*_K,\Omega^1)=h^0(X^*_K,\OO)=1$，于是 $\dim S_2(K)=1-\chi_{\mathrm{hol}}(X^*_K)=-\chi_{\mathrm{hol}}(X^*_K)+h^0(X^*_K,\C)$。[来源：R.~Hartshorne, \emph{Algebraic Geometry}, GTM 52, IV.1.3（Riemann--Roch 定理）与 III.7.6（Serre 对偶定理）。]}
\item 这里 $\chi_{\mathrm{hol}}$ 是全纯 Euler 特征标。
\item $\chi_{\mathrm{top}}(X^*_K)=2\chi_{\mathrm{hol}}(X^*_K)$。
\end{itemize}
\end{itemize}

至此，我们完成了为新形式附着 Galois 表示的既定路线（至少在权 $2$ 的情形）。

\noindent\textbf{附着于新形式的 Galois 表示}：
\begin{itemize}
\item 回忆：前两讲我们引入了概形
\[
\widehat X\to\Spec(\Q),
\]
它带有 $\GL_2(\Af)$-作用，使得自然地
\[
\widehat X(\C)\cong\GL_2(\Q)\backslash(\GL_2(\Af)\times\HH^\pm).
\]
\item 对紧开（且充分小）的 $K\subseteq\GL_2(\Af)$，得到拟射影光滑曲线
\[
X_K\to\Spec(\Q),\qquad X_K(\C)\cong\GL_2(\Q)\backslash(\GL_2(\Af)/K\times\HH^\pm).
\]
% ===== PDF p.075 =====
（注意与上一节记号的冲突：那里 $X_K$ 就等于右边的商。）\znoteplain{第 8 讲的 $X_K$ 是 $\Spec(\Q)$ 上的曲线 $\widehat X/K$（概形），而本讲从第 72 页起 $X_K$ 指的是它的 $\C$-点集 $\GL_2(\Q)^+\backslash(\GL_2(\Af)/K\times\HH)$；两者是同一对象的两种看法，只是记法相撞。译文保留原记法：凡讲「曲线、概形」时指前者，凡讲「复流形、上同调」时指后者。}
\item 固定一个素数 $\ell$。
\item 对 $K\subseteq\GL_2(\Af)$，可以考虑内平展上同调
\[
\tilde H^1_{\mathrm{\acute{e}t}}\bigl(X_{K,\Q},\Q_\ell\bigr):=\operatorname{Im}\Bigl(H^1_{c,\mathrm{\acute{e}t}}\bigl(X_{K,\Q},\Q_\ell\bigr)\to H^1_{\mathrm{\acute{e}t}}\bigl(X_{K,\Q},\Q_\ell\bigr)\Bigr)\cong H^1_{\mathrm{\acute{e}t}}\bigl(X^*_{K,\Q},\Q_\ell\bigr).
\]
\item 取极限，得到大的上同调空间
\[
\tilde H^1_{\mathrm{\acute{e}t}}\bigl(\widehat X_{\Q},\Q_\ell\bigr):=\varinjlim_{K}\tilde H^1_{\mathrm{\acute{e}t}}\bigl(X_{K,\Q},\Q_\ell\bigr).
\]
\item 由 $H^1_{c,\mathrm{\acute{e}t}}$、$H^1_{\mathrm{\acute{e}t}}$ 的自然性，这个空间带有
\[
\Gal(\overline\Q/\Q)\times\GL_2(\Af)
\]
的作用，即 $\Gal(\overline\Q/\Q)$ 的一个作用与 $\GL_2(\Af)$ 的一个作用，且这两个作用交换。
\item 固定一个同构 $\iota\colon\Ql\xrightarrow{\ \sim\ }\C$。
\item 利用平展比较定理（以及 $\iota$），得到同构
\[
\tilde H^1_{\mathrm{\acute{e}t}}\bigl(\widehat X_{\Q},\Q_\ell\bigr)\cong H^1_{\mathrm{\acute{e}t}}\bigl(\widehat X_{\C},\Q_\ell\bigr)\cong\tilde H^1\bigl(\widehat X(\C),\Q_\ell\bigr)\otimes_{\Q_\ell,\iota}\C\cong\tilde H^1\bigl(\widehat X(\C),\C\bigr).
\]
\item 这些同构都是 $\GL_2(\Af)$-等变的。
\item 设 $\pi\subseteq S_2$ 为不可约 $\GL_2(\Af)$-表示。
\item 借助 $\iota$，我们把每个 $\C$ 上的光滑 $\GL_2(\Af)$-表示都看作 $\Ql$ 上的光滑表示。
\item 空间
\[
\tilde\rho_\pi:=\Hom_{\GL_2(\Af)}\bigl(\pi(\chi_{no}),\tilde H^1_{\mathrm{\acute{e}t}}(\widehat X_{\Q},\Q_\ell)\bigr)
\]
是二维的：这由 Eichler--Shimura 同构
\[
\tilde H^1(\widehat X(\C),\C)\cong S_2(\chi_{no})\oplus\bar S_2(\chi_{no})
\]
以及 $S_2$ 作为 $\GL_2(\Af)$-表示的重数一性质得到。
\item 总之，作为 $\Gal(\overline\Q/\Q)\times\GL_2(\Af)$-表示，
\begin{equation}\tag{3}\label{eq:9.3}
\tilde H^1_{\mathrm{\acute{e}t}}\bigl(\widehat X_{\Q},\Q_\ell\bigr)\cong\bigoplus_\pi\tilde\rho_\pi\otimes_{\Q_\ell}\pi(\chi_{no}),
\end{equation}
其中求和取遍全部不可约 $\GL_2(\Af)$-表示 $\pi\subseteq S_2$（注意从右边到左边有自然的态射）。
\end{itemize}

前面已试图说明：新形式的 Langlands 互反之所以重要，在于 Hecke 特征值与 Frobenius 的迹相合。下面就来讨论这一点。

\noindent\textbf{平展上同调中的 Eichler--Shimura 关系}：
\begin{quote}
我们需要说明：Hecke 特征值与 Frobenius 的迹相合。
\end{quote}

\begin{itemize}
\item 回忆：对 $\pi\subseteq S_k$ 及其 Hecke 特征系统
\[
\bigl\{a_p(\pi),\ b_p(\pi)=\chi(p)p^{k-1}\bigr\}_{p\notin S},
\]
我们想附着二维 $\ell$-进表示
\[
\rho_\pi\colon\Gal(\overline\Q/\Q)\longrightarrow\GL_2(\Q_\ell),
\]
使得对每个 $p\notin S$，$p$ 处的算术 Frobenius $\Frob_p^{\mathrm{arith}}$ 的特征多项式为
\[
X^2-a_p(\pi)X+\chi(p)p^{k-1}
\]
% ===== PDF p.076 =====
（以下我们省略 $\iota$ 与下标 $\pi$），即
\[
\bigl(\Frob_p^{\mathrm{arith}}\bigr)^2-a_p(\pi)\Frob_p^{\mathrm{arith}}+\chi(p)p^{k-1}=0 .
\]
\item 在（式 (3)）中取 $K_1(N)$-不变量，得到（对 $N\ge3$）
\[
H^1_{\mathrm{\acute{e}t}}\bigl(X^*_{\Gamma_1(N),\Q},\Q_\ell\bigr)\cong\bigoplus_\pi\tilde\rho_\pi\otimes_{\Q_\ell}\pi(\chi_{no})^{K_1(N)}
\]
（注意 $X_{K_1(N)}\cong X_{\Gamma_1(N)}$）。
\item 设 $p$ 为素数，$p\nmid\ell N$。
\item 固定 $\overline\Q$ 中位于 $p$ 之上的一个位。这确定了 $\F_p$ 的一个代数闭包 $\overline\F_p$。
\item 回忆：对域 $L$ 上的 abelian 簇 $A$ 与素数 $\ell$，记
\[
T_\ell A=\varprojlim_n A[\ell^n](\overline L),\qquad V_\ell A=T_\ell A\otimes_{\Z_\ell}\Q_\ell
\]
分别为 $\ell$-进 Tate 模与有理化的 $\ell$-进 Tate 模。
\end{itemize}

\noindent\textbf{引理 9.2}\phantomsection\label{lem:9.2}\quad 设 $p$ 为素数，$p\nmid\ell N$。则
\[
H^1_{\mathrm{\acute{e}t}}\bigl(X^*_{\Gamma_1(N),\Q},\Q_\ell\bigr)=V_\ell\Pic^0\bigl(X^*_{\Gamma_1(N),\Fq}\bigr)^\vee,
\]
其中 $(-)^\vee$ 表示 $\Q_\ell$-对偶。

\noindent\textbf{证明}\quad Kummer 序列 $1\to\mu_{\ell^n}\to\Gm\xrightarrow{\ell^n}\Gm\to1$ 蕴涵
\[
H^1_{\mathrm{\acute{e}t}}\bigl(X^*_{\Gamma_1(N),\overline\Q},\Q_\ell(1)\bigr)\cong V_\ell\Pic^0\bigl(X^*_{\Gamma_1(N),\overline\Q}\bigr).
\]
Weil 配对给出典范同构
\[
V_\ell\Pic^0\bigl(X^*_{\Gamma_1(N),\overline\Q}\bigr)^\vee\cong V_\ell\Pic^0\bigl(X^*_{\Gamma_1(N),\overline\Q}\bigr)(-1).
\]
最后，由于 $\ell\neq p$，
\[
V_\ell\Pic^0\bigl(X^*_{\Gamma_1(N),\overline\Q}\bigr)\cong V_\ell\Pic^0\bigl(X^*_{\Gamma_1(N),\overline\Q_p}\bigr)\cong V_\ell\Pic^0\bigl(X^*_{\Gamma_1(N),\Fq}\bigr),
\]
因为挠点可以提升。\hfill$\square$

\begin{itemize}
\item 特别地，Galois 表示
\[
H^1_{\mathrm{\acute{e}t}}\bigl(X^*_{\Gamma_1(N),\Q},\Q_\ell\bigr)
\]
在 $p$ 处非分歧。
\item 在 $V_\ell\Pic^0(X^*_{\Gamma_1(N),\Fq})$ 上有 Eichler--Shimura 关系，回忆它写作
\[
F^2-\tilde T_pF+p\tilde S_p=0,
\]
其中 $F$ 是算术 Frobenius。\znoteplain{为什么在 $\Pic^0$ 的 Tate 模上关系对算术 Frobenius 成立，而命题 9.3 在 $H^1_{\mathrm{\acute{e}t}}$ 上却对几何 Frobenius 成立？由引理 9.2，$H^1_{\mathrm{\acute{e}t}}$ 是 $V_\ell\Pic^0$ 的 $\Q_\ell$-对偶，对偶把算术 Frobenius 与几何 Frobenius 互换；而转置不改变特征多项式，故两处给出的是同一个多项式。}
\end{itemize}

\noindent\textbf{命题 9.3}\phantomsection\label{prop:9.3}\quad 设 $p$ 为素数，$p\nmid\ell N$。则在 $H^1_{\mathrm{\acute{e}t}}(X^*_{\Gamma_1(N),\Q},\Q_\ell)$ 上有
\[
\bigl(\Frob_p^{\mathrm{geom}}\bigr)^2-\tilde T_p\Frob_p^{\mathrm{geom}}+p\tilde S_p=0 .
\]

\begin{itemize}
\item 对归一化新形式 $f\in S_2(\Gamma_0(N),\chi)$，我们最终定义
\[
\rho_f:=(\tilde\rho_\pi)^\vee,
\]
其中 $\pi$ 是由 $f$ 生成的不可约 $\GL_2(\Af)$-表示，而
\[
\tilde\rho_\pi:=\Hom_{\GL_2(\Af)}\bigl(\pi(\chi_{no}),\tilde H^1_{\mathrm{\acute{e}t}}(\widehat X_{\Q},\Q_\ell)\bigr).
\]
\item 关于 $\rho_f$ 的一些分析，参见 \texttt{[Rib77]}：
\begin{itemize}
% ===== PDF p.077 =====
\item 把（归一化）新形式 $f\in S_2(\Gamma_0(N),\chi)$ 写成 Fourier 展开
\[
f(q)=\sum_{n\ge1}a_nq^n .
\]
则对 $p\nmid\ell N$，$\Frob_p^{\mathrm{arith}}$ 在 $\rho_f$ 上的特征多项式为
\[
X^2-a_pX+\chi(p)p .
\]
\item 确实：
\begin{itemize}
\item 由 $\Frob_p^{\mathrm{geom}}$ 的 Eichler--Shimura 关系（的对偶），在 $\rho_f$ 上得到
\[
\bigl(\Frob_p^{\mathrm{arith}}\bigr)^2-\tilde T_p\Frob_p^{\mathrm{arith}}+p\tilde S_p=0 .
\]
\item 回忆 $\tilde T_p$ 在 $\pi^{K_1(N)}$ 上的特征值是 $pa_p$，而 $\tilde S_p$ 的特征值是 $\chi(p)p^2$。
\item 我们曾把 $S_2$ 上的 $\GL_2(\Af)$-作用按 $\chi_{no}=|\det(-)|$ 扭过。因此 $\tilde T_p$ 在
\[
\tilde\rho_\pi\otimes_{\Q_\ell}\pi(\chi_{no})^{K_1(N)}\subseteq H^1_{\mathrm{\acute{e}t}}\bigl(X^*_{\Gamma_1(N),\Q},\Q_\ell\bigr)
\]
上的特征值是 $a_p$，而 $\tilde S_p$ 的特征值是
\[
\frac{1}{p^2}\chi(p)p^2=\chi(p).
\]
\end{itemize}
\item $\rho_f$ 的行列式为
\[
\psi\cdot\chi_{\mathrm{cyc}},
\]
其中
\[
\chi_{\mathrm{cyc}}\colon G_\Q\to\Q_\ell^\times
\]
表示分圆特征，而
\[
\psi\colon\Q^\times\R_{>0}\backslash\A^\times\to\Q_\ell^\times\cong_{\iota}\C^\times
\]
是由 $\chi$ 决定的（有限阶）adélic 特征。
\item 确实：
\begin{itemize}
\item 用 Chebotarev 验证：对 $p\nmid\ell N$，两边在算术 Frobenius 元上一致。
\end{itemize}
\item 特别地，$\rho_f$ 是奇的，即每个复共轭的行列式为 $-1$。
\begin{itemize}
\item 这用到：对一切 $k\ge1$，$S_k(\Gamma_0(N),\chi)$ 非平凡蕴涵 $\chi(-1)=(-1)^k$。
\end{itemize}
\item $\rho_f$ 不可约，参见 \texttt{[Rib77, Theorem (2.3)]}。
\item 由于 $\rho_f$ 实现为一个 $\Q$ 上 proper、光滑概形的平展上同调，$\ell$-进 Hodge 理论中的 de Rham 比较定理蕴涵 $\rho_f|_{G_{\Q_\ell}}$ 是 de Rham 表示。
\end{itemize}
\end{itemize}

\noindent\textbf{权 $\ge2$ 的构造概略}：
\begin{itemize}
\item 不用 $\tilde H^1_{\mathrm{\acute{e}t}}(\widehat X_{\Q},\Q_\ell)$，而用
\[
\tilde H^1_{\mathrm{\acute{e}t}}(\widehat X_{\Q},\cL_k),
\]
其中
\[
\cL_k:=\Sym^{k-1}R^1f_*(\Q_\ell)
\]
是局部系统\idx{局部系统}，而
% ===== PDF p.078 =====
$f\colon E\to\widehat X$ 为万有椭圆曲线。\znote{这正是第 6 讲（PDF p.053）预告的「把常数层 $\Q_\ell$ 换成某个局部系统」：$R^1f_*\Q_\ell$ 是秩 $2$ 的局部系统，其 $(k-1)$ 次对称幂 $\cL_k$ 秩为 $k$，于是 $\tilde H^1_{\mathrm{\acute{e}t}}(\widehat X_{\Q},\cL_k)$ 里装的正是权 $k$ 的新形式。[来源：本讲义第 6 讲（PDF p.053）的提示，以及 $R^1f_*\Q_\ell$ 作为椭圆曲线 Tate 模局部系统的标准描述。]}
\item 问题：$\cL_k$ 不能延拓为紧化上的局部系统。因此需要额外的论证，才能为基变换到 $\overline\F_p$ 辩护。
\item 用类似的策略可以得到 Eichler--Shimura 关系与同构的类比。
\item 并把 de Rham 比较、平展比较定理换为带系数的版本。
\item 更多细节见 \texttt{[Del71b]}。
\end{itemize}

\chapter{权 1 形式的 Galois 表示}
% ===== PDF p.079 =====
\noindent\textbf{（Ben Heuer 撰）}

\noindent\textbf{上次}：对权 $\ge 2$ 的新形式，我们构造了附着其上的 Galois 表示：

\noindent\textbf{定理 10.1}（Deligne，1971）\phantomsection\label{thm:10.1}\quad 设 $N\ge 3$、$k\ge 2$，$\varepsilon:(\Z/N\Z)^\times\to\C^\times$ 为 Dirichlet 特征。设 $f\in M_k(\Gamma_0(N),\varepsilon)$ 是权 $k$ 的新形式，$K$ 为由诸 $a_p(f)$ 与 $\varepsilon$ 生成的数域。则对任一素数 $\ell\nmid N$ 以及 $K$ 中位于 $\ell$ 之上的位 $\lambda$，存在连续、奇、不可约的 Galois 表示
\[
\rho:G_\Q\longrightarrow\GL_2(K_\lambda),
\]
它在 $N$ 外非分歧，且对一切 $p\nmid N\ell$，$\Frob_p$ 的特征多项式为
\[
X^2-a_pX+\varepsilon(p)p^{k-1}.
\]

\subsection*{附着于权 1 模形式的 Galois 表示}

\begin{itemize}
\item 今日目标：证明下面的定理。
\end{itemize}

\noindent\textbf{定理 10.2}（Deligne--Serre，1973）\phantomsection\label{thm:10.2}\quad 设 $N\ge 3$，$\varepsilon:(\Z/N\Z)^\times\to\C^\times$ 为 Dirichlet 特征，$f\in M_1(\Gamma_0(N),\varepsilon)$ 是权 $1$ 的新形式。则存在连续、奇、不可约的 Galois 表示
\[
\rho:G_\Q\longrightarrow\GL_2(\C),
\]
它在 $N$ 外非分歧，且对一切 $p\nmid N$，$\Frob_p$ 的特征多项式为
\[
X^2-a_pX+\varepsilon(p).
\]

\begin{itemize}
\item 至此完成了「新形式 $\rightsquigarrow$ Galois 表示」的构造。
\item 注意 $\rho$ 的像落在 $\GL_2(\C)$ 中，即 $\rho$ 是 Artin 表示；这与权 $\ge 2$ 的情形相反，那里得到的是 $\ell$-进表示。为什么？
\begin{itemize}
\item 简单事实：任何 Artin 表示的像都有限。
\item 事实上，$\rho$ 已经可以在数域上定义，因此也可以像前面那样嵌入 $p$-进域 $K_\lambda$；但用 $\C$ 是为了标示「像有限」。
\item 从 Galois 侧看，$\rho$ 只有在 $k=1$ 时才可能像有限：因为要求 $\det\rho=\varepsilon\chi_{\mathrm{cycl}}^{k-1}$，而 $\chi_{\mathrm{cycl}}$ 的像无限。
\end{itemize}
\item 若 $\varepsilon$ 为偶，则 $M_1(\Gamma_0(N),\varepsilon)=0$。因此 $\varepsilon$ 为奇，从而 $\rho$ 为奇。
\end{itemize}

\subsection*{构造的第一版草图}

与上次相比，这个证明的几何成分更少，Galois 表示论的成分更多：

\begin{itemize}
\item 把 $f$ 模 $p$ 约化，得到一个「模 $p$ 模形式」。
\item 利用模形式的同余，说明存在模 $p$ 的 Hecke 本征形式 $f'$，它与 $f$ 有相同的特征值，但权更高，$\ge 2$。
\item 用「Deligne--Serre 提升\idx{Deligne--Serre 提升}」把 $f'$ 提升回特征 $0$。
\item 用权 $\ge 2$ 的构造，附着 Galois 表示 $\rho_{f'}$。
\item 把 Galois 表示模 $\ell$ 约化。
\item 把 Galois 表示从 $\F_\ell$ 提升到 $\C$。
\end{itemize}
% ===== PDF p.080 =====
\subsection*{模 $p$ 模形式初识}

\begin{itemize}
\item 设 $N\ge 3$，$S$ 为任一 $\Z[1/N]$-概形。
\item 回忆：我们曾把 $S$ 上权 $k$ 的模形式定义为
\[
M_k(\Gamma_1(N);S):=\Gamma\bigl(X^{*}_{\Gamma_1(N),S},\omega^k\bigr).
\]
\item $\omega$ 是万有椭圆曲线上相对微分的丰沛线丛。
\item 可以定义 $q$-展开：设 $S=\Spec(A)$，则在 $\Z((q))\otimes_\Z A$ 上有 Tate 曲线 $T(q)$；在 $\Spec(\Z((q))\otimes_\Z A)$ 上，有典范同构
\[
\omega_{\Gm}=\omega_{T(q)}.
\]
\item 线丛 $\omega_{\Gm}$ 平凡：$\Gm$ 上的典范微分 $\mathrm{d}x/x$ 是可逆截面。
\item 借助这个平凡化，对 $X^{*}_{\Gamma_1(N),A}$ 的每个尖点都得到相应的 $q$-展开
\[
M_k(\Gamma_1(N);S)\longrightarrow A[[q]].
\]
\end{itemize}

\noindent\textbf{定理 10.3}（Katz 的 $q$-展开原理）\phantomsection\label{thm:10.3}\quad 对每个尖点、所有 $k$，该映射都是单射。

\begin{itemize}
\item 现在把这一切应用于 $S=\Spec(F)$，其中 $F$ 是特征 $p$ 的有限域：所得的模形式称为\textbf{模 $p$ 模形式}。
\item 下面的事实或许有助于体会模 $p$ 模形式：
\end{itemize}

\noindent\textbf{引理 10.4}\phantomsection\label{lem:10.4}\quad 对权 $k\ge 2$，任何模 $p$ 模形式都是某个数域上的模形式模 $p$ 的约化。特别地，此时
\[
M_k(\Gamma_1(N);F)=M_k\bigl(\Gamma_1(N);\Z[1/N]\bigr)\otimes_\Z F.
\]

\noindent\textbf{证明}\quad 用 Riemann--Roch 可以证明 $H^1\bigl(X^{*}_{\Gamma_1(N)},\omega^{\otimes k}\bigr)=0$。\hfill$\square$

\begin{itemize}
\item 然而，通常给出更内禀的定义会更好。例如：
\end{itemize}

\subsection*{模 $p$ 模形式的例子：Hasse 不变量\idx{Hasse 不变量}}

\begin{itemize}
\item 回忆：对特征 $p$ 的任一椭圆曲线，有 Verschiebung 同源 $V:E^{(p)}\to E$。
\item 特别地，对万有椭圆曲线也有。
\item 在相对微分上，这给出反方向的拉回映射 $V^*:\omega\to\omega^{\otimes p}$。
\item 与 $\omega^{-1}$ 张量，得
\[
V^*:\OO\longrightarrow\omega^{\otimes(p-1)},\qquad\text{即}\quad
V^*\in\Gamma\bigl(\omega^{\otimes(p-1)}\bigr)=M_{p-1}(\Gamma_1(N);F).
\]
\item 这就是 \textbf{Hasse 不变量}，记作 $\mathrm{Ha}$。
\item 它是权 $p-1$ 的模 $p$ 模形式。事实上：它是级 $1$ 的本征形式。
\end{itemize}

\noindent\textbf{命题 10.5}\phantomsection\label{prop:10.5}\quad Hasse 不变量的 $q$-展开为常数 $1\in\F_p[[q]]$。

\noindent\textbf{证明（略述）}\quad 利用 $\omega_{T(q)}=\omega_{\Gm}$：在 $\Gm$ 上有 $F=[p]$，从而 $V=1$。\hfill$\square$

\noindent\textbf{注}：$1\in M_0(\Gamma_1(N);\F_p)$ 也是 $q$-展开为 $1$ 的模 $p$ 模形式。

\noindent\textbf{注}：这与 $q$-展开原理并不矛盾，因为两者的权不同。\znote{标准定义如下：对特征 $p$ 的椭圆曲线 $E$ 与非零微分 $\eta$，由 $F^*\eta=A(E,\eta)\eta$ 得到权 $p-1$ 的模形式 $A$，其 $q$-展开恒为 $1$；$A$ 的零点恰为 $E$ 的超奇异点（即 $E[p]$ 为 $\alpha_p$ 的扩张之处）。因此乘以 $\mathrm{Ha}$ 的映射在超奇异点处退化，这正是要挖去这些点、转向 $p$-进模形式理论的原因。[来源：B 级讲义 \url{https://math.mit.edu/nt/Murilo2020-02-26.pdf}（2026-11-26 访问）。]}\par
% ===== PDF p.081 =====
\noindent\textbf{要点}：在特征 $p$ 下出现了一个新现象：

\begin{itemize}
\item 不同权的模形式可以有相同的 $q$-展开！
\item 例如，对任一权 $k$ 的模形式 $f$ 与任一 $n\in\N$，$\mathrm{Ha}^n\cdot f$ 是权 $k+n(p-1)$ 的模形式，且与 $f$ 有相同的 $q$-展开。
\item 特别地，这保持本征形式：Hecke 算子的同一特征值系统可以对应不同权的本征形式。
\item 乘法映射
\[
M_k(\Gamma_1(N);S)\xrightarrow{\ \cdot\,\mathrm{Ha}\ }M_{k+p-1}(\Gamma_1(N);S)
\]
是单射（$q$-展开原理），但不是满射！
\item 原因：$\mathrm{Ha}$ 有零点，且其零点恰为 $X^{*}_{\Gamma_1(N)}(\Fq)$ 的超奇异点。
\item 这是挖去这些点的动机之一；挖去之后便得到 $p$-进模形式的理论。
\end{itemize}

\subsection*{Eisenstein 级数}

\begin{itemize}
\item 对此还有一个更经典的视角，我们至少应当提一下：
\end{itemize}

\noindent\textbf{定义 10.6}\phantomsection\label{def:10.6}\quad 对任一偶 $k\ge 4$，令
\[
G_k:=\sum_{\substack{(n,m)\in\Z^2\\(n,m)\neq(0,0)}}\frac{1}{(m+n\tau)^k}.
\]
这是权 $k$、级 $1$ 的模本征形式，其 $q$-展开为
\[
2\zeta(k)\Bigl(1-\frac{2k}{B_k}\sum_{n=1}^{\infty}\sigma_{k-1}(n)q^n\Bigr),
\]
其中 $\sigma_{k-1}(n)=\sum_{d\mid n}d^{k-1}$，$B_k$ 是第 $k$ 个 Bernoulli 数。把归一化的 Eisenstein 级数定义为
\[
E_k:=\frac{1}{2\zeta(k)}G_k=1-\frac{2k}{B_k}\sum_{n=1}^{\infty}\sigma_{k-1}(n)q^n\in\Z_{(p)}[[q]].
\]
\znoteplain{原文（提取文本）此处作 $\frac{4}{B_k}$；按 $G_k$ 的标准 Fourier 展开应为 $\frac{2k}{B_k}$。核对：$k=4$ 时 $-\frac{2k}{B_4}=-\frac{8}{-1/30}=240$，正是 $E_4=1+240q+\cdots$ 的系数。译文按 $\frac{2k}{B_k}$ 处理。}

\noindent\textbf{引理 10.7}\phantomsection\label{lem:10.7}\quad 设 $p\ge 5$。则就 $q$-展开而言，有
\[
E_{p-1}\equiv 1\pmod p.
\]

\noindent\textbf{证明}\quad 这由 Bernoulli 数的 Kummer 同余得到。\hfill$\square$

\noindent\textbf{要点}：当 $p\ge 5$ 时，Eisenstein 级数 $E_{p-1}$ 模 $p$ 的约化就是 $\mathrm{Ha}$：
\[
E_{p-1}\equiv\mathrm{Ha}\pmod p,
\]
此即「Deligne 同余」。

\noindent\textbf{策略}（Deligne--Serre）：

\begin{itemize}
\item 给定权 $1$ 的本征形式，模 $p$ 约化。
\item 乘以 $\mathrm{Ha}$，使其成为权 $\ge 2$ 的本征形式。
\item 在此情形使用上一讲的构造。
\end{itemize}

\subsection*{附着于模 $p$ 模形式的 Galois 表示}
% ===== PDF p.082 =====
\noindent\textbf{定理 10.8}\phantomsection\label{thm:10.8}\quad 设 $k\ge 1$，$f\in M_k(\Gamma_1(N),\varepsilon;\F_p)$ 是模 $p$ 的尖点本征形式。设 $F_f$ 为由诸 $a_\ell(f)$ 与 $\varepsilon$ 在 $\F_p$ 上生成的（有限）域。则存在半单 Galois 表示
\[
\rho_f:G_\Q\longrightarrow\GL_2(F_f),
\]
它在 $Np$ 外非分歧，且对一切 $\ell\nmid Np$，
\[
\operatorname{charpoly}(\Frob_\ell\mid\rho_f)=X^2-a_\ell X+\varepsilon(\ell)\ell^{k-1}.
\]
\znote{原文此处 $\ell$ 与 $p$ 混用：先说「在 $N\ell$ 外非分歧」，又说「对一切 $\ell\nmid Np$」，且把特征多项式写成 $\Frob_p$ 的。按上下文，$p$ 是系数域 $\F_p$ 的特征，变动的素数应记作 $\ell$；译文统一为「在 $Np$ 外非分歧，且对一切 $\ell\nmid Np$」。相应地，$F_f$ 由诸 $a_\ell(f)$ 与 $\varepsilon$ 生成。}

\noindent\textbf{第一步：归约到 $k\ge 2$ 的情形。}若 $k=1$，乘以 $\mathrm{Ha}$。这不改变 $q$-展开，因此仍是本征形式，但其权 $\ge 2$。

\noindent\textbf{第二步：把 Hecke 特征系统提升到特征 $0$。}当权 $\ge 2$ 时这总是可行的，下面就来讨论。

\subsection*{Deligne--Serre 提升引理}

\begin{itemize}
\item 设 $R$ 是无零因子的 $\Z[1/N]$-代数。设
\[
\mathbb{T}_R\subseteq\End\bigl(M_k(\Gamma_1(N);R)\bigr)
\]
为 \textbf{Hecke 代数}，即由 Hecke 算子 $T_l,S_l$ 生成的 $R$-子代数。
\item \textbf{Hecke 特征系统}是环同态 $\psi:\mathbb{T}_R\to R$。Hecke 特征系统的典型例子如下：
\item 设 $g\in M_k(\Gamma_1(N);R)$ 是本征形式，它给出 Hecke 特征系统
\[
\psi_g:\mathbb{T}_R\longrightarrow R,\qquad T\longmapsto a_T,\quad\text{使得}\ T(g)=a_Tg.
\]
\end{itemize}

\noindent\textbf{引理 10.9}\phantomsection\label{lem:10.9}\quad 设 $R=K$ 是域。则任一 Hecke 特征系统 $\psi:\mathbb{T}_K\to K$ 都来自唯一一个 Hecke 本征形式。

\noindent\textbf{证明}。
\begin{itemize}
\item $\mathbb{T}_K$ 是 Artinian $K$-代数。
\item 取扩域后，它是 Artinian 局部环 $\mathbb{T}_i$ 的乘积。
\item 由于 $\mathbb{T}_K$ 忠实作用在 $M:=M_k(\Gamma_1(N);K)$ 上，子模 $M_i:=\mathbb{T}_i\cdot M$ 非零（交换代数事实）。
\item 在 $M_i$ 中任取本征向量，它的特征系统就是 $\psi$。
\item 唯一性：Hecke 特征值决定 $q$-展开。\hfill$\square$
\end{itemize}

\begin{itemize}
\item 设 $K$ 为数域，$\p$ 为 $K$ 中位于 $p\nmid N$ 之上的位，$\F_{\p}$ 为剩余域。
\item 记 $\OO_{K,(\p)}$ 为 $K$ 中 $\p$-整元素构成的赋值子环。
\end{itemize}

\noindent\textbf{引理 10.10}（Deligne--Serre 提升引理）\phantomsection\label{lem:10.10}\quad 设 $k\ge 2$，$g\in M_k(\Gamma_1(N);\F_{\p})$ 为本征形式。则存在有限扩张 $K'|K$、位扩张 $\p'|\p$ 以及本征形式
\[
\tilde g\in M_k\bigl(\Gamma_1(N);\OO_{K',(\p')}\bigr),
\]
使得 $\tilde g$ 模 $\p'$ 约化为 $g$。

\noindent\textbf{证明}。
\begin{itemize}
\item 考虑环同态
\[
\mathbb{T}_{\OO_K}\longrightarrow\mathbb{T}_{\F_{\p}}\xrightarrow{\ \psi_g\ }\F_{\p}.
\]
\item 这对应一个极大理想 $\m$。由于 $\Spec(\mathbb{T}_{\OO_K})\to\Spec(\OO_K)$ 局部自由，可以找到素理想 $\mathfrak{q}\subseteq\m$，使得 $\mathfrak{q}\cap\OO_K=(0)$。
% ===== PDF p.083 =====
\item 这定义了 $\mathbb{T}_K$ 的一个非零素理想，对应 Hecke 特征系统 $\mathbb{T}_K\to K'$。
\item 由引理 10.9，这对应一个本征形式。
\item 用引理 10.9 的唯一性可知它约化为 $g$。\hfill$\square$
\end{itemize}
% ===== PDF p.084 =====
\subsection*{定理 10.8 的证明}

\begin{itemize}
\item 设 $f$ 是权 $\ge 2$ 的模 $p$ 本征形式，定义在 $\F_p$ 的扩域 $F_f$ 上。
\item 把它提升为某数域 $K$ 上的本征形式 $\tilde f$；设 $\p$ 为 $K$ 中位于 $p$ 之上的位。
\item 与之相伴，有 Galois 表示
\[
\rho_{\tilde f}:G_\Q\longrightarrow\GL_2(K_{\p}).
\]
\item 由于像紧，总可以找到同构的表示
\[
\rho_{\tilde f}:G_\Q\longrightarrow\GL_2(\OO_{K,\p}).
\]
\item 把它模 $\p$ 约化，得到表示
\[
\rho_{\tilde f}:G_\Q\longrightarrow\GL_2(\F_{\p}).
\]
\item 它未必已经半单。于是把 $\rho_f$ 定义为半单化
\[
\rho_f:=\rho_{\tilde f}^{\mathrm{ss}}
\]
（即 Jordan--Hölder 因子的直和）。它与 $\rho_{\tilde f}$ 有相同的迹与行列式。
\item 由于 $\ker\rho_f\subseteq\ker\rho_{\tilde f}$，它在 $pN$ 外非分歧。
\item 还需说明：$\rho_f$ 已在 $F_f$ 上定义。为此只需证明：对一切 $\sigma\in\Gal(\overline{\F}_p|F_f)$，有 $\rho_f=\sigma\circ\rho_f$。
\item 当然，对一切 $\ell\nmid pN$，$\Frob_\ell$ 的特征多项式
\[
X^2-a_\ell X+\varepsilon(\ell)\ell^{k-1}
\]
都被 $\sigma$ 保持，因为按定义 $a_\ell,\varepsilon(\ell)\in F_f$。现在用：\znoteplain{原文此处写作 $X^2-a_\ell X+\varepsilon(\ell)$，漏掉了权 $k$ 的因子 $\ell^{k-1}$；由 $\det\rho_f=\varepsilon\chi_{\mathrm{cycl}}^{k-1}$ 可知应为 $X^2-a_\ell X+\varepsilon(\ell)\ell^{k-1}$。译文补全。}
\end{itemize}

\noindent\textbf{定理 10.11}（Brauer--Nesbitt）\phantomsection\label{thm:10.11}\quad 设 $F$ 是完美域，$G$ 是有限群，$V,W$ 是 $G$ 在 $F$ 上的两个半单有限维表示。则 $V\cong W$ 当且仅当对一切 $g\in G$ 都有
\[
\operatorname{charpoly}(g\mid V)=\operatorname{charpoly}(g\mid W).
\]

\begin{itemize}
\item 当系数域的特征为正时，这取代了第 6 节中关于迹的类似陈述。
\item 至此完成了「$f\rightsquigarrow\rho_f$」的构造。\hfill$\square$
\end{itemize}

\subsection*{Serre 猜想\idx{Serre 猜想}}

\begin{itemize}
\item 在继续讨论权 $1$ 形式之前，先作一段简短的插叙。
\item 1973 年，Serre 提出了一个猜想，恰好刻画了哪些模 $p$ Galois 表示以这种方式出现。它被称为
\end{itemize}

\noindent\textbf{猜想 10.12}（Serre 模性猜想，或称 Serre 猜想）\phantomsection\label{conj:10.12}\quad 设 $\rho:G_\Q\to\GL_2(\F_p)$ 是连续、奇、不可约的 Galois 表示。则存在模 $p$ 本征形式 $g$，使得
\[
\rho=\rho_g.
\]

1987 年，Serre 发表了一篇文章，细化了这个猜想：

\begin{itemize}
\item 他预测了 $g$ 的最优级（用 $\rho$ 的 Artin 导子\idx{导子}表出）。
\item 他预测了最优权（用 $\rho$ 在 $p$ 处的分歧给出配方）。
\item 他证明了：他的猜想将蕴含 Fermat 大定理。
\item Serre 猜想现已成为定理，由 Khare--Wintenberger（2008）证明。
\end{itemize}
% ===== PDF p.085 =====
\subsection*{把 $\rho_f$ 提升到 $\C$}

\begin{itemize}
\item 回到先前的设定：$f\in M_1(\Gamma_0(N),\varepsilon)$ 是本征形式。
\item $K$ 是由特征值与 $\varepsilon$ 生成的数域。
\item 令 $L$ 为在 $K$ 中完全分解的 $\Q$ 的素数之集。由 Chebotarev 密度定理，这是无限集。
\item 设 $\ell\in L$，$\lambda$ 为 $K$ 中位于 $\ell$ 之上的位，$\OO_{K,(\lambda)}\subseteq K$ 为 $\lambda$-整元素。
\item 现在把上一节应用于 $p$ 由变动中的 $\ell\in L$ 扮演的情形。
\item 对每个这样的 $\lambda$，考虑 $f$ 模 $\lambda$ 的约化 $\bar f$；这是模 $\ell$ 的本征形式。
\item 定理 10.8 给出半单 Galois 表示
\[
\rho_{f,\ell}:G_\Q\longrightarrow\GL_2(\F_\ell)
\]
（因 $\lambda$ 完全分裂，系数域就是 $\F_\ell$）；对任一 $p\nmid\ell N$，$\Frob_p$ 的特征多项式为
\[
\operatorname{charpoly}(\Frob_p\mid\rho_{f,\ell})=X^2-a_pX+\varepsilon(p).
\]
\item 它的零点是 $n$ 次单位根；由下一引理，可以假设它们全在 $\OO_K$ 中，从而在 $\F_\ell$ 中。于是该特征多项式 $=(X-a)(X-b)$，其中 $a,b\in\F_\ell^\times$。\znoteplain{这里用了两件事：其一，$K$ 已扩大为包含所有 $n\le A$ 次单位根，故 $\F_\ell^\times$ 中的 $n$ 次单位根在 $K$ 中的提升（Teichmüller 提升）落在 $K$ 内，且是 $\lambda$-整的，即在 $\OO_{K,(\lambda)}$ 中；其二，$\lambda$ 完全分裂，故剩余域就是 $\F_\ell$，这些单位根模 $\lambda$ 后落在 $\F_\ell$ 中。}
\end{itemize}

\noindent\textbf{定义 10.13}\phantomsection\label{def:10.13}\quad 令 $G_\ell:=\operatorname{im}(\rho_{f,\ell})\subseteq\GL_2(\F_\ell)$。这是有限群。

\noindent\textbf{引理 10.14}\phantomsection\label{lem:10.14}\quad 若 $f$ 的权为 $1$，则存在 $A\in\N$，使得对一切 $\ell\in L$ 都有 $|G_\ell|\le A$。

\noindent\textbf{证明}\quad 在用到关于权 $1$ 模形式系数的一个关键估计之后，这就变成纯群论问题：对 $\GL_2(\F_\ell)$ 的子群分类。我们不讨论它；若想看，见 \texttt{[DS74, Lemma 8.4]}。\hfill$\square$

\begin{itemize}
\item 扩大 $K$，使其包含所有 $n\le A$ 的 $n$ 次单位根。这会缩小 $L$。
\item 再把 $\le A$ 的素数从 $L$ 中排除；此时 $L$ 仍无限。
\end{itemize}

\noindent\textbf{引理 10.15}\phantomsection\label{lem:10.15}\quad 设 $G$ 是阶与 $\ell$ 互素的有限群。则任一表示 $\rho:G\to\GL_2(\F_\ell)$ 都可提升为表示
\[
\tilde\rho:G\longrightarrow\GL_2(\OO_{K,(\lambda)}).
\]

\noindent\textbf{证明}。\textbf{第一步：提升到 $\GL_2(\Z_\ell)$}。

\begin{itemize}
\item 只需证明：任一 $G\to\GL_2(\Z/\ell^n\Z)$ 都可提升为 $G\to\GL_2(\Z/\ell^{n+1}\Z)$。
\item 为此考虑短正合列
\[
1\longrightarrow 1+\ell^nM_2(\F_\ell)\longrightarrow\GL_2(\Z/\ell^{n+1}\Z)\longrightarrow\GL_2(\Z/\ell^n\Z)\longrightarrow 1.
\]
\item 第一个群通过 $1+\ell^nx\mapsto x$ 同构于 $M_2(\F_\ell)$。
\item 现在用非交换群上同调：考虑左作用，$G$ 通过
\[
(g\cdot\phi)(x):=\phi(g)^{-1}\phi(gx)
\]
作用在 $\mathrm{Map}\bigl(G,\GL_2(\Z/\ell^n\Z)\bigr)$ 上。\znoteplain{原文此处在提取文本中作 $g\cdot\phi:=\phi(g)\phi(g)^{-1}$，明显残缺。按非交换群上同调的标准扭作用应为 $(g\cdot\phi)(x)=\phi(g)^{-1}\phi(gx)$：其不动点恰为同态 $\phi:G\to\GL_2(\Z/\ell^n\Z)$，故有下一个等式。译文据此重建。}
\item 于是
\[
\Hom\bigl(G,\GL_2(\Z/\ell^n\Z)\bigr)=\mathrm{Map}\bigl(G,\GL_2(\Z/\ell^n\Z)\bigr)^G.
\]
\item 对上述短正合列施加 $\mathrm{Map}$，再由 $G$ 得到非交换群上同调的长正合列
\[
\cdots\to\Hom(G,\GL_2(\Z/\ell^{n+1}))\to\Hom(G,\GL_2(\Z/\ell^n))\to H^1\bigl(G,\mathrm{Map}(G,M_2(\F_\ell))\bigr).
\]
\item 群上同调的事实：若 $|G|$ 与 $|A|$ 有限且互素，则 $H^1(G,A)=0$。
\end{itemize}
% ===== PDF p.086 =====
\noindent\textbf{第二步：已在数域上定义。}

\begin{itemize}
\item 由于 $G$ 有限，$\rho$ 半单。
\item 迹必为阶 $\le|G|\le A$ 的单位根之和，从而落在 $K$ 中。
\item 由 Brauer--Nesbitt 可知 $\rho$ 已在 $K$ 上定义。
\item 最后，$K\cap\overline{\Z}_\ell=\OO_{K,(\lambda)}$。\hfill$\square$
\end{itemize}

\subsection*{回到定理 10.2 的证明（权 1 形式的 Galois 表示）}

\begin{itemize}
\item 把它应用于 $\rho_{f,\ell}:G_\ell\hookrightarrow\GL_2(\F_\ell)$，得到 $G_\ell\to\GL_2(\OO_{K,(\lambda)})$。\znoteplain{原文这一段的提取文本把下标时而写作 $g$、时而写作 $f$（例如 $\rho_{g,\ell}$、$\tilde\rho_{g,\ell}$ 与 $\rho_{f,\ell}$ 并存）。按上下文，本讲自始至终讨论的是同一个权 $1$ 本征形式 $f$，故译文统一作下标 $f$。}
\item 再与 $G_\Q\to G_\ell$ 复合，得到 Galois 表示
\[
\tilde\rho_{f,\ell}:G_\Q\longrightarrow\GL_2(\OO_{K,(\lambda)}).
\]
\end{itemize}

\noindent\textbf{命题 10.16}\phantomsection\label{prop:10.16}\quad 提升 $\tilde\rho_{f,\ell}:G_\Q\to\GL_2(\OO_{K,(\lambda)})$ 满足：对一切 $p\nmid\ell N$，
\[
\operatorname{charpoly}(\Frob_p\mid\tilde\rho_{f,\ell})=X^2-a_pX+\varepsilon(p).
\]

关键在于变动 $\ell$：回忆我们对所有 $\ell\in L$ 都得到 $\tilde\rho_{f,\ell}$。

\noindent\textbf{定义 10.17}\phantomsection\label{def:10.17}\quad 令 $Y$ 为形如 $(X-\alpha)(X-\beta)$ 的多项式之集，其中 $\alpha,\beta\in\OO_{K,(\lambda)}$ 是阶 $\le A$ 的单位根。这是有限集。

\begin{itemize}
\item \textbf{注}：给定 $n\le A$、$\ell\in L$ 以及 $n$ 次单位根 $x\in\F_\ell^\times$，存在唯一的提升为 $K$ 中的单位根（存在性：$n\le A$；唯一性：$\ell\ge A$）。
\item \textbf{推论}：若 $F,G\in Y$ 且 $F\equiv G\pmod\lambda$，则 $F=G$。
\item \textbf{想法}：把它应用于 $F=X^2-a_pX+\varepsilon(p)$ 与 $G=\operatorname{charpoly}(\Frob_p\mid\tilde\rho_{f,\ell})$。
\end{itemize}

\noindent\textbf{引理 10.18}\phantomsection\label{lem:10.18}\quad $X^2-a_pX+\varepsilon(p)\in Y$。

\noindent\textbf{证明}\quad 对每个位于 $L$ 中素数之上的 $\lambda$，由于 $\rho_{f,\ell}$ 的像有限，存在 $F\in Y$ 使得
\[
\operatorname{charpoly}(\Frob_p\mid\rho_{f,\ell})\equiv F\pmod\lambda.
\]
于是也有
\[
X^2-a_pX+\varepsilon(p)\equiv F\pmod\lambda.
\]
因 $Y$ 有限，存在某个 $F$ 使上式对无穷多个 $\lambda$ 成立。故 $X^2-a_pX+\varepsilon(p)=F\in Y$。\hfill$\square$

\noindent\textbf{命题的证明}。
\begin{itemize}
\item 由于 $G_\ell$ 有限且阶 $\le A$，有 $\operatorname{charpoly}(\Frob_p\mid\tilde\rho_{f,\ell})\in Y$。
\item 又
\[
\operatorname{charpoly}(\Frob_p\mid\tilde\rho_{f,\ell})\equiv\operatorname{charpoly}(\Frob_p\mid\rho_{f,\ell})\equiv X^2-a_pX+\varepsilon(p)\pmod\lambda.
\]
\item 两边都在 $Y$ 中，故 $\operatorname{charpoly}(\Frob_p\mid\tilde\rho_{f,\ell})=X^2-a_pX+\varepsilon(p)$。\hfill$\square$
\end{itemize}

尚需证明：

\begin{itemize}
\item $\rho$ 为奇：这仅因为 $\varepsilon$ 为奇且 $c$ 的阶为 $2$。
\item $\rho$ 不可约：用 Rankin 的一个复分析估计。\hfill$\square$
\end{itemize}

至此完成了「本征形式 $\rightsquigarrow$ Galois 表示」的构造！
% ===== PDF p.087 =====
\subsection*{权 1 构造的总结}

\begin{center}
\begin{tikzcd}[row sep=3.4em, column sep=2.6em]
\text{权 }1\text{ 本征形式}
  \arrow[r, "{\text{同余}}"]
  \arrow[d, "{\text{模 }\ell\text{ 约化}}"']
& \text{权 }\ge 2\text{ 本征形式}
  \arrow[r, "{\text{Deligne 构造}}"]
  \arrow[d, "{\text{模 }\ell\text{ 约化}}"']
& \ell\text{-进 Galois 表示}
  \arrow[d, "{\text{模 }\ell\text{ 约化}}"]
&
\\
\text{模 }\ell\text{ 的权 }1\text{ 本征形式}
  \arrow[r, "{\cdot\,\mathrm{Ha}}"]
& \text{模 }\ell\text{ 的权 }\ge 2\text{ 本征形式}
  \arrow[u, "{\text{DS 提升}}"]
  \arrow[r, "{\text{模 }\ell\text{ 构造（定理 10.8）}}"]
& \text{模 }\ell\text{ 的 Galois 表示}
  \arrow[r, "{\text{提升}}"]
& \text{复 Galois 表示}
\end{tikzcd}
\end{center}

\znoteplain{原第 87 页是一张示意图（方框加箭头）。此处按原图的七个方框重排为两行四列的交换图：「$\cdot\,\mathrm{Ha}$」表示乘以 Hasse 不变量，「DS 提升」即引理 10.10，「模 $\ell$ 构造」用定理 10.8，最后的「提升」即本讲末尾把 $\rho_{f,\ell}$ 提升到 $\C$ 的步骤。原图为扫描排版，箭头走向按图中文字标注复原，宜与原书对照。}

\chapter{一般群的 Langlands 纲领 I}

% ===== PDF p.088 =====
\noindent\textbf{接下来的目标}
\begin{itemize}
\item 对 $\GL_n$、乃至 $\Q$ 上一般的约化群 $G$，给出（部分）Langlands 纲领的或多或少精确的陈述。
\end{itemize}

\medskip
\noindent\textbf{$\GL_{n,F}$ 的 Langlands（–Clozel–Fontaine–Mazur）互反}
\begin{itemize}
\item 设 $F$ 为任意数域。
\item 固定某个素数 $\ell$ 与同构 $\C\cong\Ql$。
\item Langlands（–Clozel–Fontaine–Mazur）互反猜想断言：对任意 $n\ge1$，存在一个（唯一的）双射，介于
\begin{itemize}
\item[(i)] $\GL_n(\A_F)$ 的 $L$-代数尖点自守表示之集，
\item[(ii)] 不可约连续表示 $\Gal(\overline F/F)\to\GL_n(\Ql)$ 的同构类之集，这些表示几乎处处非分歧，且在整除 $\ell$ 的位处为 de Rham 表示，
\end{itemize}
之间，使得该双射把 Satake 参数与 Frobenius 元素的特征值相匹配。\znote{原文此处作 “Langland(-Clozel-Fontaine-Mazur)”，应为 “Langlands”，漏了一个 s；译文按 Langlands 处理。}
\end{itemize}

\medskip
\noindent\textbf{评论}
\begin{itemize}
\item 这只是 $G=\Res_{F/\Q}\GL_{n,F}$ 的 Langlands 纲领的一部分。
\item 它既没有描述 $\GL_{n,F}$ 的全部自守表示，也没有描述全部的 Galois 表示。
\item 对 $F=\Q$、$n\ge2$ 未知。例如：$\GL_2(\A)$ 存在 $L$-代数尖点自守表示（与 Maass 形式相伴），它们应当对应于 $\Gal(\overline\Q/\Q)$ 的偶的、不可约的二维 $\ell$-进表示（像有限）。
\item Langlands–Tunnell：若像可解，则该 Galois 表示的自守性已知，参见 \texttt{[GH19, Theorem 13.4.5.]}。这并不覆盖像为 $\SL_2(\F_5)$ 的、几乎处处非分歧的二维不可约 Galois 表示。
\item 唯一性可由双方的（强）重数一定理得到。
\item 今日：解释「$L$-代数」是什么意思。
\end{itemize}

\medskip
\noindent\textbf{一般群的自守形式}
\begin{itemize}
\item $G$ 为 $\Q$ 上任意的约化群。
\item 我们将为 $L^2([G])$ 引入一个方便的（= 更代数的）替代物。
\item 回忆第 3 节中，对 $G=\GL_2=\GL_{2,\Q}$、$k\in\Z$ 的嵌入
\[
\Phi\colon M_k\longrightarrow C^\infty(\GL_2(\A)),\qquad f\longmapsto\phi_f.
\]
\item 回忆 $\Phi(M_k)$ 与 $\Phi(S_k)$ 的像的描述，即：
\begin{itemize}
\item 对 $k\in\Z$，元素 $\phi\in C^\infty(\GL_2(\A))$ 落在 $\Phi\colon M_k\to C^\infty(\GL_2(\A))$ 的像中，当且仅当
\begin{itemize}
\item 对 $(g,g_\infty)\in\GL_2(\A)$、$z\in\C^\times\subseteq\GL_2(\R)$，有 $\phi(g,g_\infty z)=z^{-k}\phi(g,g_\infty)$；
\item $\phi$ 是适度增长的（这蕴涵负 Fourier 系数为零）；
\item 对每个 $g\in\GL_2(\Af)$ 有 $Y*\phi(g,-)=0$，
% ===== PDF p.089 =====
其中 $Y\in\g_\C=\mathfrak{gl}_2(\R)_\C$ 是适当构造的自然元素（这蕴涵全纯性）；
\item 对一切 $\gamma\in\GL_2(\Q)$ 有 $\phi(\gamma g)=\phi(g)$。
\end{itemize}
\item $f$ 还是尖点的\idx{尖点的}，当且仅当对一切含于真抛物子群 $P\subseteq G$（定义在 $\Q$ 上）的幺幂根基 $N$，有
\[
\int_{N(\Q)\backslash N(\A)}\phi_f(ng)\,\mathrm{d}n=0.
\]
\end{itemize}
\item 固定 $G(\R)$ 的一个极大紧子群 $K_\infty$（通常不连通），例如 $\mathrm{O}_2(\R)=\SO_2(\R)\cup s\SO_2(\R)\subseteq\GL_2(\R)$，其中
\[
s=\begin{pmatrix}-1&0\\0&1\end{pmatrix}.
\]
\end{itemize}

\medskip
\noindent\textbf{定义 11.1}（参见 \texttt{[GH19, Definition 6.5]}）\phantomsection\label{def:11.1}\quad 设 $G$ 为 $\Q$ 上的约化群。$G$ 的一个（阿代尔）\textbf{自守形式}是光滑函数 $\phi\colon G(\A)\to\C$，使得
\begin{itemize}
\item[1)] $\phi$ 是适度增长的（第 4 节已定义；本讲不需要）；
\item[2)] $\phi$ 是右 $K_\infty$-有限的，即函数 $g\mapsto\phi(gk)$（$k\in K_\infty\subseteq G(\A)$）张成有限维向量空间；
\item[3)] $\phi$ 被 $\cZ(\g_\C)$ 中一个有限余维的理想杀死，其中 $\g_\C$ 记 $G(\R)$ 的（复化）Lie 代数，$\cZ(\g_\C)$ 记 $\g_\C$ 的泛包络代数 $U(\g_\C)$ 的中心；
\item[4)] $\phi$ 是左 $G(\Q)$-不变的，即对一切 $\gamma\in G(\Q)$ 有 $\phi(\gamma g)=\phi(g)$。
\end{itemize}

\medskip
\noindent\textbf{$\GL_2$ 的 Casimir 元}
\begin{itemize}
\item 我们来验证：对 $f\in M_k$，函数 $\phi_f\in C^\infty(\GL_2(\A))$ 确实满足 3)。
\item 也就是说，元素 $Y$ 并不在 $\cZ(\mathfrak{gl}_{2,\C})$ 中。
\item 置
\[
H:=\begin{pmatrix}0&-i\\ i&0\end{pmatrix},\qquad
X:=\frac12\begin{pmatrix}-1&-i\\ -i&1\end{pmatrix},\qquad
Y:=\frac12\begin{pmatrix}-1&i\\ i&1\end{pmatrix},\qquad
Z:=\begin{pmatrix}1&0\\ 0&1\end{pmatrix}.
\]
\item 于是
\begin{itemize}
\item[1)] $H*\phi_f=-k\phi_f$，因为对 $z\in\C^\times$ 有 $z\cdot\phi_f=z^{-k}\phi_f$；
\item[2)] $Y*\phi_f=0$。
\end{itemize}
\item $[X,Y]=XY-YX=H$（因为 $H,X,Y$ 构成 $\mathfrak{sl}_2$-三元组）。
\item 定义 Casimir 算子
\begin{equation}\tag{1}\label{eq:11.1}
\Delta:=\tfrac14\bigl(H^2+2XY+2YX\bigr).
\end{equation}
则 $\Delta\in\cZ(\mathfrak{gl}_{2,\C})$。
\item 事实上，$\cZ(\mathfrak{gl}_{2,\C})$ 作为 $\C$-代数由 $\Delta$ 与 $Z\in\mathfrak{gl}_{2,\C}$ 生成（这可由 Harish-Chandra 同构推出，参见 \texttt{[GH19, Theorem 4.6.1]}）。
\item 计算（利用 $Y*\phi_f=0$）得
\begin{equation}\tag{2}\label{eq:11.2}
\Delta*\phi_f=\tfrac14\bigl(H^2+4YX+2H\bigr)*\phi_f=\tfrac14\bigl(H^2+2H\bigr)*\phi_f=\tfrac14(k^2-2k)\phi_f,
\end{equation}
于是 $\phi_f$ 确实被 $\cZ(\mathfrak{gl}_{2,\C})$ 中有限余维的理想 $\bigl\langle\Delta-\tfrac14(k^2-2k),\,Z+k\bigr\rangle$ 杀死。\znote{原文此处作 $\Delta*\phi_f=\frac14(H^2-2XY+2YX)*\phi_f=\frac14(H^2-2H)*\phi_f=\frac14(k^2-2k)\phi_f$：中间式与 $\Delta$ 的定义不符，且由 $H*\phi_f=-k\phi_f$，该式会给出 $\frac14(k^2+2k)\phi_f$，与末式矛盾。按 $[X,Y]=H$（即 $XY=YX+H$）有 $2XY+2YX=4YX+2H$；再由 $Y*\phi_f=0$ 得 $YX*\phi_f=0$，故 $\Delta*\phi_f=\frac14(H^2+2H)*\phi_f=\frac14(k^2-2k)\phi_f$。译文按更正后的计算写出。[来源：底本 p.~89；见疑点清单 p.89 条]}
\end{itemize}

% ===== PDF p.090 =====
\medskip
\noindent\textbf{关于自守形式的补充}
\begin{itemize}
\item $G$ 为 $\Q$ 上任意的约化群。
\item 记 $\cA(G)$ 为自守形式构成的空间，$\cA([G])$ 为其中在 $A_G$（经 $G(\A)$ 上的左平移或右平移作用）下不变的自守形式构成的子空间。\znoteplain{记号 $[G]$ 沿用第 2 讲：$[G]=G(\Q)A_G\backslash G(\A)$。}
\item 遗憾的是，空间 $\cA([G])$ 与 $\cA(G)$ 在 $G(\R)$ 下并不稳定，因为 $K_\infty$-有限性未必被保持。
\end{itemize}

\medskip
\noindent\textbf{$K_\infty$-有限性不被保持的例子（据 Andreas Mihatsch）}
\begin{itemize}
\item 取 $G=\GL_2(\R)$、$K:=K_\infty^\circ=\SO_2(\R)\cong S^1$，并取 $V$ 为 $G$ 上函数空间中的任一表示，只要它包含 $C_c^\infty(G)$。我们给出一个例子，说明 $K$-有限性在右平移作用下不被保持。
\item 更精确地说：存在 $\phi\in C_c^\infty(G)$ 与 $h\in G$，使得 $\phi$ 是 $K$-有限的，但不是 $K_h:=h^{-1}Kh$-有限的。
\item 为此：取 $h\in G$ 使 $K_h\cap K=\{\pm1\}$。
\item 再取非零的 $\phi\in C_c^\infty(G)$，它在 $K$ 下不变，且在 $K_h$ 的某个开子集上 $\phi\equiv0$（例如：把 $G/K$ 上的函数拉回以构造 $\phi$；注意 $K_h$ 在 $G/K$ 中的像维数 $>0$）。
\item $K_h$ 上任何非零的 $K_h$-有限函数都不能在开子集上为零。事实上，$K_h\cong S^1$ 上的 $K_h$-有限函数恰为 Laurent 多项式的限制。
\end{itemize}

\medskip
\noindent\textbf{关于自守形式的补充（续）}
\begin{itemize}
\item 显然，$G(\Af)$ 经 $\phi\mapsto\bigl(g\mapsto\phi(gh)\bigr)$（$h\in G(\Af)$）作用在 $\cA(G)$ 上。
\item 类似地，$K_\infty$、乃至 $A_GK_\infty$，作用在 $\cA(G)$ 上。
\item 不那么显然、但为真的是（参见 \texttt{[GH19, Proposition 4.4.2]}）：$\g_\C$ 作用在 $\cA(G)$ 上。$\Rightarrow$ $U(\g_\C)$、特别是 $\cZ(\g_\C)\subseteq U(\g_\C)$，经微分算子作用在 $\cA(G)$ 上。
\item $\g_\C$-作用与 $K_\infty$-作用使 $\cA(G)$ 成为下述意义下的 $(\g_\C,K_\infty)$-模。
\end{itemize}

\medskip
\noindent\textbf{定义 11.2}（参见 \texttt{[GH19, Definition 4.5]}）\phantomsection\label{def:11.2}\quad 一个 $(\g_\C,K_\infty)$-模是一个 $\C$-向量空间 $V$，它同时是 $\g_\C$-模与 $K_\infty$-模，且满足
\begin{itemize}
\item[1)] $V$ 是有限维 $K_\infty$-稳定子空间的并；
\item[2)] 对 $X\in\mathfrak k:=\Lie(K_\infty)_\C$ 与 $\phi\in V$，有
\[
X*\phi=\frac{\partial}{\partial t}\bigl(\exp(tX)\phi\bigr)\Big|_{t=0};
\]
\item[3)] 对 $h\in K_\infty$、$X\in\g_\C$ 与 $\phi\in V$，有
\[
h\bigl(X*(h^{-1}\phi)\bigr)=\bigl(\mathrm{Ad}(h)(X)\bigr)*\phi.
\]
\end{itemize}

\medskip
\noindent\textbf{评论}
\begin{itemize}
\item 我们对 $V$ 不要求任何拓扑，因此 $(\g_\C,K_\infty)$-模比 $G(\R)$-表示代数得多，也就更容易处理。
\item 2) 中的极限是对 $V$ 中某个包含 $\phi$ 的有限维 $K_\infty$-稳定子空间的自然拓扑取的。
\item 对 $G(\R)$ 在某个 Hilbert 空间（视为 $\C$-向量空间）上的每个表示，光滑且 $K_\infty$-有限的向量构成的空间自然地是 $(\g_\C,K_\infty)$-模，参见 \texttt{[GH19, Proposition 4.4.2]}；而对酉表示不会丢失任何信息，参见 \texttt{[GH19, Theorem 4.4.4]}，即：$G(\R)$ 的不可约酉表示的等类单射到不可约 $(\g_\C,K_\infty)$-模的等类中。
\item 此外，如此得到的 $(\g_\C,K_\infty)$-模是容许的，即：对 $K_\infty$ 的每个（有限维）不可约表示 $\sigma$，其 $\sigma$-等型分量是有限维的，参见 \texttt{[GH19, Theorem 4.4.1]}。
\item Schur 引理的一个变体：在每个不可约容许 $(\g_\C,K_\infty)$-模 $\pi$ 上，$U(\g_\C)$ 的中心 $\cZ(\g_\C)$ 经某个 $\C$-代数同态 $\omega_\pi\colon\cZ(\g_\C)\to\C$ 作用；$\omega_\pi$ 称为 $\pi$ 的\textbf{无穷小特征}（infinitesimal character）。
\item Harish-Chandra 同构（参见 \texttt{[GH19, Theorem 4.6.1]}）给出：对每个 Cartan 子代数 $\mathfrak t\subseteq\g$ 与相应的 Weyl 群 $W$，有 $\cZ(\g_\C)\cong U(\mathfrak t_\C)^W$。注意 $U(\mathfrak t_\C)\cong\Sym^\bullet\mathfrak t_\C$，因为 $\mathfrak t_\C$ 是 abelian Lie 代数。
\item 由于 $W$ 有限，$\C$-代数同态 $\cZ(\g_\C)\cong U(\mathfrak t_\C)^W\to\C$ 与 $\C$-代数同态 $U(\mathfrak t_\C)\cong\Sym^\bullet\mathfrak t_\C\to\C$ 的 $W$-轨道一一对应。此外
\[
\Hom_{\C\text{-}\mathrm{alg}}\bigl(U(\mathfrak t_\C),\C\bigr)\cong\Hom_{\C\text{-}\mathrm{lin}}\bigl(\mathfrak t_\C,\C\bigr)=:\mathfrak t_\C^\vee,
\]
且 $U(\mathfrak t_\C)^W\hookrightarrow U(\mathfrak t_\C)$ 是有限扩张。
\item 于是无穷小特征定义了「权空间」$\mathfrak t_\C^\vee$ 中一个典范的 $W$-轨道。
\item 现在设对某个极大环面 $T\subseteq G_\R$ 有 $\mathfrak t=\Lie(T)$。则 $W$-等变地有 $X^*(T_\C)\subseteq\mathfrak t_\C^\vee$。
% ===== PDF p.091 =====
\item 不可约容许 $(\g_\C,K_\infty)$-模 $\pi$ 称为\textbf{$L$-代数的}，若与无穷小特征 $\omega_\pi$ 相应的 $\mathfrak t_\C^\vee$ 中的 $W$-轨道落在有限自由 $\Z$-模 $X^*(T_\C)\subseteq\mathfrak t_\C^\vee$ 中，参见 \texttt{[BG10, Definition 2.3.1]}、\texttt{[Tho, Definition 93]}。
\end{itemize}

\medskip
\noindent\textbf{由模形式生成的 $(\mathfrak{gl}_{2,\C},\mathrm{O}_2(\R))$-模}
\begin{itemize}
\item 设 $f\in M_k$。我们来分析由 $\phi:=\phi_f$ 生成的 $(\mathfrak{gl}_{2,\C},\mathrm{O}_2(\R))$-模 $V\subseteq\cA(\GL_2)$。
\item 回忆 $Y*\phi=0$，且对 $z\in\C^\times$ 有 $z\cdot\phi=z^{-k}\phi$。
\item 不难看出：存在唯一的（在同构意义下）由这样的元素 $\phi$ 生成的 $(\mathfrak{gl}_{2,\C},\mathrm{O}_2(\R))$-模 $D'_{k-1}$（即 \texttt{[Del73, Section 2.1]} 中 $D_{k-1}$ 的对偶），并且 $D'_{k-1}\cong V$ 不可约。
\item 特别地，它只依赖于 $k$，不依赖于 $f$。
\item 正如之前某处提到的，我们得到
\[
M_k\cong\Hom_{(\mathfrak{gl}_{2,\C},\mathrm{O}_2(\R))}\bigl(D'_{k-1},\cA(\GL_2)\bigr)
\]
（注：\texttt{[Del73, Scholie 2.1.3]} 中的空间 $H(G_\A)$ 与我们的 $\cA(\GL_2)$ 只在 $\GL_2(\A)$ 上取逆的意义下一致）。
% ===== PDF p.092 =====
\item 我们算得 $\Delta*\phi=\tfrac14(k^2-2k)\phi$、$Z*\phi=-k\phi$，这确定了无穷小特征 $\omega_{D'_{k-1}}$。
\item 在 Harish-Chandra 同构下，与 $\omega_{D'_{k-1}}$ 相应的 $\C$-代数同态 $U(\mathfrak t_\C)\cong\C[Z,H]\to\C$ 的 $W$-轨道由同态
\[
Z\longmapsto-k,\qquad H\longmapsto\pm(k-1)
\]
给出。
\item 格 $X^*(T)\subseteq\langle Z,H\rangle^\vee_\C$ 由满足 $a,b\in\Z$ 且 $a\equiv b\pmod2$ 的 $\C$-线性映射 $Z\mapsto a$、$H\mapsto b$ 给出。
\item 因此 $D'_{k-1}$ 不是 $L$-代数的（它在 \texttt{[BG10]} 的意义下是 $\C$-代数的）。\znoteplain{这里 $L$-代数性要求该 $W$-轨道落在 $X^*(T)$ 中，即要求 $Z\mapsto-k$ 与 $H\mapsto\pm(k-1)$ 这两个坐标同奇偶；但 $-k$ 与 $k-1$ 的奇偶性总是相反，故 $D'_{k-1}$ 从不 $L$-代数。扭上 $|\det|^{1/2+\alpha}$ 之后，$Z$ 的坐标变为 $1-k+2\alpha$，它与 $k-1$ 同奇偶，故扭后确是 $L$-代数的。另外注意这里用的是特征格 $X^*(T)$（$W$-等变地含于权空间 $\mathfrak t^\vee_\C$），而非余特征格 $X_*(T)$。}
\item 然而，对任意 $a\in\Z$，$D'_{k-1}\otimes_\C|\det|^{1/2+a}$ 是 $L$-代数的，因为（与 $|\det|^{1/2}$ 相伴的 $(\mathfrak{gl}_{2,\C},\mathrm{O}_2(\R))$-模）的无穷小特征对应于 $Z\mapsto1$、$H\mapsto0$。
\end{itemize}

\medskip
\noindent\textbf{要点}
\begin{itemize}
\item 对任意 $G$ 都存在（阿代尔）自守形式的空间 $\cA(G)$，当 $G=\GL_2$ 时它自然地包含诸 $\phi_f$（$f\in M_k$）。
\item 空间 $\cA(G)$ 带有 $(\g_\C,K_\infty)\times G(\Af)$-作用，是带 $G(\A)$-作用的 $L^2([G])$ 的一个更代数的替代物。
\end{itemize}

\medskip
\noindent\textbf{自守表示的一个新概念}
\begin{itemize}
\item 我们用 $\cA(G)$ 代替 $L^2([G])$ 来重新定义自守表示。
\item 即：$G$ 的自守表示是 $\cA(G)$ 的不可约 $G(\Af)\times(\g_\C,K_\infty)$-子商，参见 \texttt{[GH19, Definition 6.8]}。
\item 此前的自守表示今后称为「$L^2$-自守表示」。
\item 两种概念之间的关系参见 \texttt{[GH19, Section 6.5]}。
\end{itemize}

\medskip
\noindent\textbf{Flath 定理（参见 \texttt{[GH19, Theorem 5.7.1]}）}
\begin{itemize}
\item 对几乎所有素数 $p$，约化群 $G_{\Q_p}:=G\times_{\Spec(\Q)}\Spec(\Q_p)$ 是「非分歧的」，即它可以延拓为约化群概形 $G_p\to\Spec(\Z_p)$。
\item 等价地，$G_{\Q_p}$ 在一个非分歧扩张上是准分裂的（即含一个定义在 $\Q_p$ 上的 Borel 子群）并且是分裂的（即含一个极大分裂环面）。
\item 于是
\begin{equation}\tag{3}\label{eq:11.3}
G(\A)=\sideset{}{'}\prod_{p}\bigl(G(\Q_p),G_p(\Z_p)\bigr)\times G(\R),
\end{equation}
其中 $G_p(\Z_p)\subseteq G(\Q_p)$ 是紧开拓子群（参见 \texttt{[GH19, Proposition 2.3.1]}），称为\textbf{超特殊的}。
% ===== PDF p.093 =====
\item Flath 定理（参见 \texttt{[GH19, Section 5.7]} 与 \texttt{[Fla79]}）断言：不可约容许 $G(\Af)\times(\g_\C,K_\infty)$-模 $\pi$（唯一地）分解为「限制张量积」
\begin{equation}\tag{4}\label{eq:11.4}
\pi\cong\sideset{}{'}\bigotimes_{p\ \text{素数}}\pi_p\otimes\pi_\infty,
\end{equation}
其中 $\pi_p$ 是不可约光滑（甚至是容许的）$G_{\Q_p}$-表示，$\pi_\infty$ 是不可约容许 $(\g_\C,K_\infty)$-模。
\item 此外，对几乎所有素数 $p$，表示 $\pi_p$ 非分歧，即群 $G_{\Q_p}$ 非分歧且 $\pi_p^{G_p(\Z_p)}\ne0$，其中 $G_p$ 是 $G_{\Q_p}$ 的约化模型。
\item 重要（参见 \texttt{[GH19, Theorem 7.5.1]}）：$\cH(G(\Q_p),G_p(\Z_p))$ 仍然是交换的！（$\Rightarrow$ 若 $\pi_p^{G_p(\Z_p)}\ne0$，则其维数为 $1$。）
\end{itemize}

\medskip
\noindent\textbf{要点}
\begin{itemize}
\item 于是不可约容许 $G(\Af)\times(\g_\C,K_\infty)$-模是一族「局部数据」。
\item 而「是自守的」在这些局部数据之间加上了很强的关系。
\item 对几乎所有素数 $p$ 得到同态
\[
\chi_p\colon\cH(G(\Q_p),G_p(\Z_p))\longrightarrow\C,
\]
与 $\GL_2$ 的情形类似；$\Rightarrow$ 这将给出「Hecke 特征系统」的类比。
\item 对 $\pi$ 置 $\pi_f:=\sideset{}{'}\bigotimes_p\pi_p$，于是 $\pi\cong\pi_f\otimes\pi_\infty$。
\end{itemize}

\medskip
\noindent\textbf{定义 11.3}（$L$-代数自守表示）\phantomsection\label{def:11.3}\quad 自守表示 $\pi\cong\pi_f\otimes\pi_\infty$ 称为\textbf{$L$-代数的}，若不可约容许 $(\g_\C,K_\infty)$-模 $\pi_\infty$ 是 $L$-代数的，参见 \texttt{[BG10, Definition 3.1.1]}。

\medskip
\noindent\textbf{例}
\begin{itemize}
\item 由模形式生成的自守表示 $\pi\subseteq\cA(\GL_2)$ 不是 $L$-代数的，但对任意 $a\in\Z$，其扭 $\pi\otimes_\C|\det|^{1/2+a}_{\text{adèlic}}$ 是 $L$-代数的。
\item 连续特征标 $\psi\colon\Q^\times\backslash\A^\times\to\C$ 是 $L$-代数的，当且仅当对某个 $k\in\Z$ 有 $\psi=\chi|\cdot|^k_{\text{adèlic}}$，其中 $|\cdot|_{\text{adèlic}}$ 是阿代尔范数，$\chi\colon\Q^\times\backslash\A^\times\to\C$ 是有限阶特征标。
\end{itemize}

\chapter{一般群的 Langlands 纲领 II}

% ===== PDF p.094 =====
\noindent\textbf{上次}
\begin{itemize}
\item $F$ 为任意数域。
\item $\GL_{n,F}$ 的 Langlands（–Clozel–Fontaine–Mazur）互反猜想断言：
\begin{itemize}
\item $\ell$ 为某个素数。固定一个同构 $\C\cong\Ql$。
\item 则对任意 $n\ge1$，存在一个（唯一的）双射，介于
\begin{itemize}
\item[(i)] $\GL_n(\A_F)$ 的 $L$-代数尖点自守表示之集，
\item[(ii)] 不可约连续表示 $\Gal(\overline F/F)\to\GL_n(\Ql)$ 的同构类之集，这些表示几乎处处非分歧，且在整除 $\ell$ 的位处为 de Rham 表示，
\end{itemize}
之间，使得该双射把 Satake 参数\idx{Satake 参数}与 Frobenius 元素的特征值相匹配。
\end{itemize}
\item 对任意 $G$ 引入了（阿代尔）自守形式的空间 $\cA(G)$。
\item 由 Flath 定理 $\Rightarrow$ 自守表示 $\pi$ 分解：$\pi\cong\sideset{}{'}\bigotimes_p\pi_p\otimes\pi_\infty$。
\item 解释了「$L$-代数」（= $\pi_\infty$ 的无穷小特征是「整的」）。
\end{itemize}

\medskip
\noindent\textbf{今日}
\begin{itemize}
\item 引入「尖点」自守表示。
\item 解释 Satake 参数。
\item 讨论一般 $G$ 的 Langlands 纲领的一些预期。
\end{itemize}

\medskip
\noindent\textbf{定义 12.1}（\texttt{[GH19, Definition 9.2]}）\phantomsection\label{def:12.1}\quad 设 $G/\Q$ 约化。称 $\phi\in L^2([G])$（相应地 $\phi\in\cA(G)$）为\textbf{尖点的}，若对一切真抛物子群 $P\subseteq G$（定义在 $\Q$ 上）的幺幂根基 $N$ 以及几乎所有 $g\in G(\A)$，有
\begin{equation}\tag{1}\label{eq:12.1}
\int_{N(\Q)\backslash N(\A)}\phi(ng)\,\mathrm{d}n=0.
\end{equation}

\medskip
\noindent\textbf{尖点自守表示}
\begin{itemize}
\item 记 $L^2_{\mathrm{cusp}}([G])$、相应地 $\cA_{\mathrm{cusp}}(G)$、$\cA_{\mathrm{cusp}}([G])$，为尖点元素构成的子空间。
\item 存在像稠密的嵌入
\[
\cA_{\mathrm{cusp}}([G])\subseteq L^2_{\mathrm{cusp}}([G]),
\]
即：尖点自守形式满足足够强的增长条件，使其属于 $L^2$（它们是「速降」的），参见 \texttt{[GH19, Section 6.5.]}。
\item $G(\A)$ 的 $L^2$-自守表示称为\textbf{尖点的}，若它同构于 $L^2_{\mathrm{cusp}}([G])$ 的某个子商（实际上是子表示）。
\item 类似地，自守表示称为尖点的，若它作为 $\cA_{\mathrm{cusp}}(G)$ 的子商出现。
\item Gelfand、Piatetski-Shapiro：作为酉 $G(\A)$-表示，
% ===== PDF p.095 =====
\begin{equation}\tag{2}\label{eq:12.2}
L^2_{\mathrm{cusp}}([G])\cong\widehat{\bigoplus}_{\pi\in\widehat{G(\A)}}m_\pi\,\pi
\end{equation}
其中每个 $m_\pi$ 有限，即 $m_\pi\in\N\cup\{0\}$，参见 \texttt{[GH19, Corollary 9.1.2.]}。
\item 特别地，$L^2_{\mathrm{cusp}}([G])\subset L^2_{\mathrm{disc}}([G])$（注意：平凡表示出现在 $L^2_{\mathrm{disc}}([G])\setminus L^2_{\mathrm{cusp}}([G])$ 中）。
\item 由
\[
L^2_{\mathrm{cusp}}([G])\cong\widehat{\bigoplus}_{\pi\in\widehat{G(\A)}}m_\pi\,\pi
\]
推出分解（分解为不可约容许 $G(\Af)\times(\g_\C,K_\infty)$-模）
\begin{equation}\tag{3}\label{eq:12.3}
\cA_{\mathrm{cusp}}([G])\cong\bigoplus_{\pi\in\widehat{G(\A)}}m_\pi\,\pi^{K_\infty\text{-}\mathrm{finite}},
\end{equation}
其中 $\pi^{K_\infty\text{-}\mathrm{finite}}\subseteq\pi$ 记 $K_\infty$-有限（特别是光滑）向量构成的（稠密）子空间，参见 \texttt{[GH19, Theorem 4.4.4.]}。
\item 于是，$L^2_{\mathrm{cusp}}([G])$ 的分解与 $\cA_{\mathrm{cusp}}([G])$ 的分解「相同」。
\item Piatetski-Shapiro、Shalika，参见 \texttt{[GH19, Theorem 11.3.4]}：设 $F$ 为数域、$G=\Res_{F/\Q}\GL_{n,F}$。则 $L^2_{\mathrm{cusp}}([G])$ 上的 $G(\A)$-表示是重数自由的（「重数一」）。
\item 「强重数一」：设 $\pi,\pi'$ 是 $\GL_{n,F}$ 的尖点自守表示。若对几乎所有素数 $p$ 有 $\pi_p\cong\pi'_p$，则 $\pi\cong\pi'$，参见 \texttt{[GH19, 11.7.2]}。
\item 去掉尖点性后不成立，参见 \texttt{[BG10, page 38]}。
\item Mœglin、Waldspurger，参见 \texttt{[GH19, Theorem 10.7.1.]}、\texttt{[MW89]}：设 $G=\Res_{F/\Q}\GL_{n,F}$。则在把 $d\mid n$ 的 $\GL_{d,F}$ 的尖点自守表示参数化的意义下，$L^2_{\mathrm{disc}}([G])$ 是已知的，并且它被证明是重数自由的。
\end{itemize}

我们现在可以补上一个缺口。即：我们从未给新形式配上自守表示。但用现在陈述的材料，我们至少可以给出（两种）形式定义。设 $f\in S_k$ 为新形式。则我们可以给 $f$ 配上 $\cA(\GL_2)$ 中由自守形式 $\phi_f$ 生成的（不可约）$\GL_2(\Af)\times(\mathfrak{gl}_{2,\C},\mathrm{O}_2(\R))$-子模（参见第 3 节），也可以配上 $L^2([\GL_2])$ 中由 $\tilde\phi_f$ 生成的不可约酉子表示（参见第 3 节）。具体用哪种构造，要在相应的具体情形中检验。

\medskip
\noindent\textbf{Satake 参数}
\begin{itemize}
\item 描述 $\pi$ 在 $p$ 处非分歧时的 $\pi_p$，即 $\pi_p^{G_p(\Z_p)}\ne0$，其中 $G_p\to\Spec(\Z_p)$ 是 $G_{\Q_p}$ 的约化模型（对几乎所有素数 $p$ 存在）。
\item 也就是描述 $\C$-代数同态
\[
\cH(G(\Q_p),K)\longrightarrow\C,\qquad K:=G_p(\Z_p).
\]
\item Langlands/Satake，参见 \texttt{[GH19, Theorem 7.5.1., Corollary 7.5.2.]}、\texttt{[BG10, Section 2.1.]}：存在
\begin{itemize}
\item 一个「自然的」$\Q$ 上代数群 $\widehat G$，
\item 一个自同构 $\Frob_p\colon\widehat G(\Q)\cong\widehat G(\Q)$，
% ===== PDF p.096 =====
\item 以及一个双射
\begin{equation}\tag{4}\label{eq:12.4}
\Hom_{\C\text{-}\mathrm{alg}}\bigl(\cH(G(\Q_p),K),\C\bigr)\;\overset{1:1}{\longleftrightarrow}\;\bigl\{\widehat G(\C)\rtimes\Frob_p^{\Z}\ \text{中的 Frobenius 半单}\ \widehat G(\C)\text{-共轭类}\bigr\}.
\end{equation}
\end{itemize}
\item 群 $\widehat G_\C$ 与 $G$ 具有对偶的根数据。因此例如（参见 \texttt{[Bor79]}）：
\begin{itemize}
\item 若 $G=\GL_n$，则 $\widehat G=\GL_{n,\C}$；
\item 若 $G=\SL_n$，则 $\widehat G=\mathrm{PGL}_{n,\C}$；
\item 若 $G=\SO_{2n}$，则 $\widehat G=\SO_{2n,\C}$；
\item 若 $G=\SO_{2n+1}$，则 $\widehat G=\mathrm{Sp}_{2n,\C}$；
\item 若 $G=\mathrm{GSp}_{2n}$，则 $\widehat G=\mathrm{GSpin}_{2n+1,\C}$。
\end{itemize}
\item 若 $G_{\Q_p}$ 分裂，则自同构 $\Frob_p$ 平凡，而 Frobenius 半单共轭类就是 $\widehat G(\C)$ 中半单元素的共轭类。
\item 若 $G=\GL_n$，则半单共轭类由其元素的特征多项式唯一确定。
\item 设 $\pi=\sideset{}{'}\bigotimes_p\pi_p\otimes\pi_\infty$ 是 $G$ 的自守表示。则我们得到下述「Hecke 特征系统」的类比：
\begin{itemize}
\item 素数的一个有限集 $S$，使得 $p\notin S$ 时 $\pi_p$ 非分歧，
\item 对每个 $p\notin S$，$\widehat G(\C)\rtimes\Frob_p^{\Z}$ 中一个 Frobenius 半单共轭类 $c_p(\pi)$。
\end{itemize}
\item 对 $\GL_n$，$c_p(\pi)$ 中（每个元素的）特征值就是 $\pi$ 的 Satake 参数。
\item 具体地，若双陪集
\[
\GL_n(\Z_p)\diag\bigl(\underbrace{p,\dots,p}_{i\ \text{个}},\underbrace{1,\dots,1}_{(n-i)\ \text{个}}\bigr)\GL_n(\Z_p)
\]
在 $\pi_p^{\GL_n(\Z_p)}$ 上的特征值为 $e_{a_{p,i}}$，则 $c_p(\pi)$ 中元素的特征多项式为
\begin{equation}\tag{5}\label{eq:12.5}
X^n-p^{\frac{1-n}{2}}e_{a_{p,1}}X^{n-1}+\cdots+(-1)^ip^{\frac{i(i-n)}{2}}e_{a_{p,i}}X^i+\cdots+(-1)^ne_{a_{p,n}},
\end{equation}
参见 \texttt{[GH19, Section 7.2.]}。
\item 更具体地，对由新形式 $f=\sum_{n\ge1}a_nq^n\in S_k(\Gamma_0(N),\chi)$ 生成的 $\pi$，$p\nmid N$ 时的 $c_p(\pi)$ 具有特征多项式
\begin{equation}\tag{6}\label{eq:12.6}
X^2-p^{-1/2}p\,a_pX+\chi(p)p^k.
\end{equation}
\end{itemize}

回忆 $\pi\otimes_\C|\det|^{1/2}_{\text{adèlic}}$ 是 $L$-代数的。对 $\pi\otimes_\C|\det|^{1/2}_{\text{adèlic}}$ 得到多项式
\[
X^2-a_pX+\chi(p)p^{k-1},
\]
（在选定同构 $\Ql\cong\C$ 之后）它曾是 $p$ 处算术 Frobenius 的特征多项式。\znote{原文第 (6) 式确作 $X^2-p^{-1/2}p\,a_pX+\chi(p)p^k$，即一次项系数为 $p^{1/2}a_p$，而非 $p^{-1/2}a_p$。按此处的 Satake 参数归一化，扭 $|\det|^{1/2}$ 会把每个 Satake 参数乘上 $p^{-1/2}$：于是 $p^{1/2}a_p\cdot p^{-1/2}=a_p$，$\chi(p)p^k\cdot p^{-1}=\chi(p)p^{k-1}$，正得下一式。两式自洽，故此处并非笔误。[来源：底本 p.~96 原式]}

\medskip
\noindent\textbf{L-群}
\begin{itemize}
\item $G$ 为 $\Q$ 上的约化群。
\item $\ell$ 为素数。
\item 稍作工作（参见 \texttt{[GH19, Section 7.3.]}）可知，群 $\Gal(\overline\Q/\Q)$ 作用在 $\Q$ 上的约化群 $\widehat G$ 上（若 $G$ 分裂，则此作用平凡）。
% ===== PDF p.097 =====
\item 置
\begin{equation}\tag{7}\label{eq:12.7}
{}^LG:=\widehat G(\Ql)\rtimes\Gal(\overline\Q/\Q).
\end{equation}
\item 所谓\textbf{L-参数}，按定义是连续同态
\[
\Gal(\overline\Q/\Q)\longrightarrow{}^LG
\]
的 $\widehat G(\Ql)$-共轭类，这些同态到 $\Gal(\overline\Q/\Q)$ 的投影为恒等映射。
\item 若 $F/\Q$ 有限且 $G=\Res_{F/\Q}\GL_{n,F}$，则一个 L-参数等同于 $\Gal(\overline F/F)$ 的一个 $n$ 维 $\ell$-进 Galois 表示的同构类。
\end{itemize}

\medskip
\noindent\textbf{关于 $L$-代数自守表示的 Buzzard–Gee 猜想}
\begin{itemize}
\item 设 $\ell$ 为素数。
\item 固定一个同构 $\iota\colon\C\cong\Ql$。
\item 设 $\pi$ 是 $G$ 的 $L$-代数自守表示。
\item 则 Buzzard–Gee（参见 \texttt{[BG10, Conjecture 3.2.2.]}）猜想存在 L-参数
\[
\rho_\pi\colon\Gal(\overline\Q/\Q)\longrightarrow{}^LG,
\]
使得（特别地）
\begin{itemize}
\item 若 $p\ne\ell$ 且 $\pi$ 在 $p$ 处非分歧，则每个算术 Frobenius $\Frob_p\in\Gal(\overline{\Q}_p/\Q_p)$ 在
\[
\Gal(\overline{\Q}_p/\Q_p)\longrightarrow\widehat G(\Ql)\rtimes\Gal(\overline{\Q}_p/\Q_p)\twoheadrightarrow\widehat G(\Ql)\rtimes\Frob_p^{\widehat\Z}
\]
下的像落在 $\iota(c_p(\pi))\in\widehat G(\Ql)\rtimes\Frob_p^\Z$ 的共轭类中；
\item 对每个连续表示 ${}^LG\to\GL_N(\Ql)$（它在 $\widehat G(\Ql)$ 上是代数的），复合
\[
\Gal(\overline{\Q}_\ell/\Q_\ell)\xrightarrow{\ \rho_\pi\ }{}^LG\longrightarrow\GL_N(\Ql)
\]
是 de Rham 的。
\end{itemize}
\item 这推广了「把 Galois 表示附着到某个新形式所对应的（某个扭的）自守表示上」这一构造。
\item 与 $\GL_n$ 的情形相反，这里并不猜想 L-参数 $\rho_\pi$ 唯一，参见 \texttt{[BG10, Remark 3.2.4.]}。
\item 此外，对不同的 $L$-代数自守表示 $\pi$，我们可能得到同一个 L-参数 $\rho$。
\item 对 $\GL_n$，Buzzard–Gee 猜想在此前的 Langlands–Clozel–Fontaine–Mazur 猜想之外，还预言存在附着于（可能非尖点的）$L$-代数自守表示的 Galois 表示。此时这些 Galois 表示不再被猜想为不可约。
\end{itemize}

\medskip
\noindent\textbf{Langlands 群\idx{Langlands 群} $\cL$}
\begin{itemize}
\item 猜想存在一个（非常大的）局部紧群 $\cL$，即（整体）\textbf{Langlands 群}，具有（至少）下述性质，参见 \texttt{[GH19, Section 12.6.]}、\texttt{[Art02]}、\texttt{[LR87]}：\znoteplain{原文把这个群记作花体的 $L$，与 $L$-函数、L-群等术语中的字母 $L$ 同形，容易混淆；译文一律记作 $\cL$。}
\begin{itemize}
% ===== PDF p.098 =====
\item 存在典范满射 $\cL\twoheadrightarrow\Gal(\overline\Q/\Q)$（事实上 $\cL\twoheadrightarrow W_\Q$，其中 $W_\Q$ 是整体 Weil 群）。
\item 对每个 $n\ge0$，$\cL$ 的连续不可约 $n$ 维 $\C$-表示的等类 $\mathrm{Irr}_n(\cL)$ 自然地与 $\GL_n(\A)$ 的尖点自守表示之集 $\mathrm{Csp}_n$ 一一对应。
\item 固定一个同构 $\Ql\cong\C$。则存在单射，从不可约、几乎处处非分歧且在 $\ell$ 处为 de Rham 的 $\ell$-进 Galois 表示 $\Gal(\overline\Q/\Q)$ 的同构类之集 $\mathcal{G}_{\mathrm{geom}}$，映入 $\cL$ 的不可约表示的等类之集 $\mathrm{Irr}$，其像对应于 $L$-代数尖点自守表示。
\item 对 $\Q$ 的每个位 $v$，存在局部 Langlands 群的单射 $\cL_v\to\cL$，它在 $\cL$ 中共轭意义下良定义。这里：
\begin{itemize}
\item 若 $v=\infty$，则 $\cL_v\cong W_\R$，即 $\C^\times$ 被 $\Gal(\C/\R)$ 的非分裂扩张；
\item 若 $v=p$ 为素数，则 $\cL_v\cong W_{\Q_p}\times\mathrm{SU}_2(\R)$，其中
\[
W_{\Q_p}\cong I_{\Q_p}\rtimes\Z\subseteq\Gal(\overline{\Q}_p/\Q_p)\cong I_{\Q_p}\rtimes\widehat\Z
\]
是映入 Frobenius 的整数幂的元素构成的子群。
\end{itemize}
\item 若 $\pi\cong\sideset{}{'}\bigotimes_v\pi_v\in\mathrm{Csp}_n$ 对应于连续表示 $\rho_\pi\colon\cL\to\GL_n(\C)$，则在 $\GL_n$ 的局部 Langlands 对应下，$\pi_v$ 应当对应于 $\rho_\pi|_{\cL_v}$，参见 \texttt{[GH19, Section 12.5.]}。
\end{itemize}
\item 若这些条件成立，则必然有 $\cL^{\mathrm{ab}}\cong\Q^\times\backslash\A^\times$。注意由整体类域论存在满射 $\Q^\times\backslash\A^\times\twoheadrightarrow\Gal(\overline\Q/\Q)^{\mathrm{ab}}$。
\item $\cL$ 的存在性目前完全无法企及；更进一步地，猜想 $\cL$ 应满射到 $\C$ 上的「动机 Galois 群」$G_{\mathrm{mot}}$，参见 \texttt{[Art02]}、\texttt{[LR87]}。
\end{itemize}

\medskip
\noindent\textbf{一般 $G$ 的整体 Langlands 对应（非常、非常粗略）}
\begin{itemize}
\item 这里不追求任何精确性，所有陈述都近乎空泛！
\item 精确的陈述可查阅 \texttt{[GH19, Chapter 12.]} 或 \texttt{[Art94]}。
\item 为了参数化非 $L$-代数的自守表示（至少是那些与 $L^2_{\mathrm{disc}}([G])$ 的分解有关的），应当把此前 Galois 版本的 L-参数换成（某些）L-参数
\[
\cL\longrightarrow{}^LG
\]
（实际上是 L-群的 Weil 形式 $W_\Q\ltimes\widehat G(\C)$）。
\item 于是人们希望构造满射（参见 \texttt{[GH19, Conjecture 12.6.2.]}）
\[
\{(\text{某些})\ \text{自守表示}\}\ \xrightarrow{\ \mathcal{LL}\ }\ \{(\text{某些})\ \text{L-参数}\}.
\]
\item $\mathcal{LL}$ 的纤维称为\textbf{L-包}（L-packet），它们可能无限。
\item 局部 Langlands 对应的精确构造（参见 \texttt{[GH19, Conjecture 12.5.1.]}）应当给出 L-包的参数化，参见 \texttt{[GH19, (12.23)]}。
% ===== PDF p.099 =====
\item 对每个 $\pi\subseteq L^2_{\mathrm{disc}}([G])$，其重数应当可以通过 L-参数 $\mathcal{LL}(\pi)$ 与 L-包的精确参数化算出，参见 \texttt{[GH19, Conjecture 12.6.3.]}。
\item 然而，对非缓增（non-tempered）表示，上述做法应当不合理；此时应考虑 Arthur 参数
\[
\cL\times\SL_2(\C)\longrightarrow{}^LG
\]
代替 L-参数，参见 \texttt{[Art94]}。
\end{itemize}

\medskip
\noindent\textbf{本讲未提及的重要内容}
\begin{itemize}
\item L-函数，\texttt{[Bor79]}、\texttt{[GH19, Chapter 11, Section 12.7.]}；
\item 函子性，\texttt{[GH19, Section 12.6.]}；
\item 迹公式，\texttt{[Art05]}、\texttt{[GH19, Section 18]}；
\item 函数域的 Langlands 纲领；
\item ……
\end{itemize}

\medskip
\noindent\textbf{Shimura 簇\idx{Shimura 簇}一瞥}
\begin{itemize}
\item 设 $K_\infty\subseteq G(\R)$ 为极大的（连通）紧子群。
\item 设 $K\subseteq G(\Af)$ 为（充分小的）紧开拓子群。
\item 回忆：
\begin{itemize}
\item $[G]\to[G]/K=G(\Q)A_G\backslash G(\A)/K$ 是实流形上的射有限覆盖；
\item $[G]/K\to X_K:=[G]/KK_\infty=G(\Q)A_G\backslash G(\A)/KK_\infty$ 是 $X_K$ 上的 $K_\infty$-丛，而 $X_K$ 是算术流形的无交并（= 对称空间对算术子群的商）。
\end{itemize}
\item 注意 $X_K$ 隐含地依赖于 $K_\infty$。
\item 为简化后面的一些记号，置
\[
\widehat X:=\varprojlim_{K}X_K.
\]
\end{itemize}

\medskip
\noindent\textbf{对 $G=\GL_2$ 已出现过的东西}
\begin{itemize}
\item[(1)] 上半/下半平面 $\HH^\pm\cong G(\R)/K_\infty$，它自然地是复流形；
\item[(2)] 全纯嵌入 $\HH^\pm\to\Pone_\C$；
\item[(3)] $\HH^\pm$ 上 $G(\R)$-等变的向量丛 $\omega_k$，或经拉回成为 $K\subseteq G(\Af)$ 紧开拓时 $X_K$ 上的向量丛；
\item[(4)] 对 $k\in\Z$，模形式空间 $M_k\subseteq H^0(\widehat X,\omega_k)$，其中右端指全纯截面；
\item[(5)] $X_K$ 的典范紧化 $X_K^*$，以及 $\omega_k$ 在其上的延拓；
\item[(6)] 一个概形 $\widehat X\to\Spec(\Q)$，其 $\C$-值点自然地同构于 $\widehat X$；
% ===== PDF p.100 =====
\item[(7)] 由 $\widetilde X$ 上某些 $G(\Af)$-等变的 $\Ql$-局部系统 $\mathbb L$（例如 $\mathbb L\cong\Ql$）得到的 $\Gal(\overline\Q/\Q)\times G(\Af)$-表示\znote{原文在第 (4)、(6) 条用 $\widehat X$，在第 (7)–(9) 条改用 $\widetilde X$，所指应当相同（即前述 $\varprojlim_KX_K$）；译文保留原文的两种写法。}
\[
\widetilde H^1_{\mathrm{\acute{e}t}}\bigl(\widetilde X_{\overline\Q},\mathbb L\bigr);
\]
\item[(8)] Eichler–Shimura 同构：把 $k\ge2$ 的 $S_k$ 与某些 $\widetilde H^1_{\mathrm{\acute{e}t}}(\widetilde X_{\overline\Q},\mathbb L)$ 联系起来；
\item[(9)] Eichler–Shimura 关系：表达 $\widetilde H^1_{\mathrm{\acute{e}t}}(\widetilde X_{\overline\Q},\mathbb L)$ 上 $\Gal(\overline\Q/\Q)$-作用与 $G(\Af)$-作用之间的关系，它是某条模曲线的 mod $p$ 几何的推论；
\item[(10)] 一个映射
\[
\mathcal{LL}\colon A_{\mathrm{mod}}\longrightarrow G_{\mathrm{mod}},
\]
其像（猜想）可以描述。
\end{itemize}

对一般的 $G$，可惜并没有类比——但是等等：存在能给出 Shimura 簇的 $G$！

\begin{itemize}
\item 甚至第 (1) 条对一般的 $G$ 就失效：空间 $G(\R)/K_\infty$ 不必是复流形！（例如 $G=\GL_3$ 时，实维数
\[
\dim\bigl(\GL_3(\R)/A_{\GL_3}(\R)\SO_3(\R)\bigr)=\frac{3(3+1)}{2}-1=5
\]
是奇数。）
\item 特别地，一切与全纯性有关的东西对这样的 $G$ 都没有类比，例如第 (2)、(3)、(4)、(6)、(7) 条等。
\item 这是坏消息！
\item 好消息：对能给出 Shimura 数据\idx{Shimura 数据}的群 $G$（例如 $F/\Q$ 全实时的 $\Res_{F/\Q}\GL_{2,F}$，或 $\mathrm{GSp}_{2n}$ 等，参见 \texttt{[Lan17]}、\texttt{[Mil05]}、\texttt{[Del71a]}、\texttt{[Del79]}），实流形
\[
G(\R)/K_\infty A_G
\]
是复流形。
\item 一个 \textbf{Shimura 数据}是一对 $(G,X)$，其中 $G$ 是 $\Q$ 上的约化群，$X$ 是态射
\[
h\colon\mathbb S:=\Res_{\C/\R}\Gm\longrightarrow G_\R
\]
的 $G(\R)$-共轭类，使得
\begin{itemize}
\item[1)] 对每个 $h\in X$，$\mathbb S(\R)=\C^\times$ 经 $\mathrm{Ad}\circ h$ 在 $\Lie(G_\R)_\C$ 上作用的特征标或为 $z\mapsto z/\bar z$，或为 $z\mapsto1$，或为 $z\mapsto\bar z/z$；
\item[2)] 对每个 $h\in X$，群 $\{g\in G(\C)\mid h(i)gh(i)^{-1}=g\}$ 模其中心是紧的；
\item[3)] 伴随群 $G^{\mathrm{ad}}$ 没有定义在 $\Q$ 上、且 $h\in X$ 的投影在其上平凡的单因子。
\end{itemize}
\item 例如取 $G=\GL_2$，$X$ 为通常嵌入 $\C^\times\to\GL_2(\R)$ 的共轭类。注意 $X\cong\HH^\pm$。
\item 一般地，对 $h\in X$，$h$ 在 $G(\R)$ 中的稳定化子 $K_h$ 模 $G(\R)$ 的中心是紧的（这由条件 2) 推出），且 $X\cong G(\R)/K_h$（$\cong G(\R)/A_GK_\infty$）。
\end{itemize}

% ===== PDF p.101 =====
\begin{itemize}
\item 对 $K\subseteq G(\Af)$，回忆
\[
X_K:=G(\Q)\backslash\bigl(G(\Af)/K\times X\bigr).
\]
另一种记号：$Sh_K(G,X)$，即附着于 $(G,X)$ 的级 $K$ 的 \textbf{Shimura 簇}。
\item 我们得到：
\begin{itemize}
\item[(1)] $X$ 是 Hermite 对称域的无交并，参见 \texttt{[Mil05]}、\texttt{[Del79]}。
\item[(2)] 对每个 $h\in X$ 自然地附着抛物子群 $P_h\subseteq G(\C)$（参见 \texttt{[CS17, Section 2.1]}），且 $h\mapsto P_h$ 定义了开嵌入
\[
\pi\colon X\cong G(\R)/K_h\longrightarrow\mathcal{Fl}\cong G(\C)/P_h,
\]
即把 $X$ 嵌入一个旗簇（这推广了嵌入 $\HH^\pm\to\Pone(\C)$）。\znoteplain{原文脚注（16）：把 $X$ 视为「（极化）Hodge 结构的模空间」可以解释这个嵌入，参见 \texttt{[Del79]}。}
特别地，$X$ 带有自然的复结构。更强的是（Baily–Borel）：$X_K$ 自然地是拟射影簇，参见 \texttt{[Lan17, Theorem 2.4.1.]}。
\item[(3)] 旗簇\idx{旗簇} $\mathcal{Fl}$ 上 $G(\C)$-等变向量丛的拉回给出 $G(\Af)\times X$ 上 $G(\Q)$-等变向量丛。经下降，对任意紧开拓 $K\subseteq G(\Af)$ 得到 $X_K$ 上的\textbf{自守向量丛\idx{自守向量丛}} $\mathcal E$，即 $\omega_k$（$k\in\Z$）的类比，参见 \texttt{[Har88]}。注意 $\mathcal{Fl}$ 上 $G(\C)$-等变向量丛等价于 $P_h$ 的（代数）表示。
\item[(4)、(5)] 若 $\dim_\C(X_K)>1$，则确实存在紧化 $X_K$，但它们不再是典范的。不过，可以在 $H^*(X_K,\mathcal E)$ 的上同调中定义子空间，以推广 $k\in\Z$ 的 $M_k$，参见 \texttt{[Har88]}。
\item[(6)] 空间 $\widehat X$（= 复流形的逆向极限）再次作为概形 $\widehat X\to\Spec(E)$ 的 $\C$-点出现，该概形即\textbf{典范模型\idx{典范模型}}，定义在某个数域（「Shimura 数据的反射域」）上，参见 \texttt{[Lan17, Theorem 2.4.3.]}。
\item[(8)、(9)、(10)] 有了典范模型，就可以尝试（有时会成功）得到第 (8)、(9)、(10) 条的类比，并最终把（某些）自守表示与（各种变体的）Galois 表示联系起来，参见 \texttt{[Har88]}、\texttt{[BRb]}。不用说，一切都（紧化、Eichler–Shimura 同构/关系、整模型、mod $p$ 几何、内上同调与相交上同调……）要复杂得多！\znote{原文此处作 “obtain analogs of points (18),(19),(20)”，与本条开头的 “(8),(9),(10)” 不一致，按上下文此处应为 (8)、(9)、(10)，译文据此处理。又，原文 “By descent one one obtains” 衍一 “one” 字，译文按 “one obtains” 处理。}
\end{itemize}
\end{itemize}

\backmatter
\chapter{参考文献}

% ===== PDF p.102 =====
\begin{references}
\refitem{Art94}{Art94} James Arthur. \emph{Unipotent automorphic representations: conjectures}. Department of Mathematics, University of Toronto, 1994.
\refitem{Art02}{Art02} James Arthur. A note on the automorphic Langlands group. \emph{Canadian Mathematical Bulletin}, 45(4):466--482, 2002.
\refitem{Art05}{Art05} James Arthur. An introduction to the trace formula. \emph{Harmonic analysis, the trace formula, and Shimura varieties}, 4:1--263, 2005.
\refitem{BC09}{BC09} Olivier Brinon and Brian Conrad. \emph{CMI summer school notes on $p$-adic Hodge theory}, 2009.
\refitem{BG10}{BG10} Kevin Buzzard and Toby Gee. The conjectural connections between automorphic representations and Galois representations. \emph{Automorphic forms and Galois representations}, 1:135--187, 2010.
\refitem{Bor79}{Bor79} Armand Borel. Automorphic $L$-functions. In \emph{Automorphic forms, representations and $L$-functions} (Proc. Sympos. Pure Math., Oregon State Univ., Corvallis, Ore., 1977), Part 2, pages 27--61, 1979.
\refitem{Bor80}{Bor80} A. Borel. Stable and $L^2$-cohomology of arithmetic groups. \emph{Bull. Amer. Math. Soc. (N.S.)}, 3(3):1025--1027, 1980.
\refitem{BRa}{BRa} J. N. Bernstein and K. Rumelhart. Lectures on $p$-adic groups. 未出版, 见 \url{http://www.math.tau.ac.il/~bernstei/Publication_list}.
\refitem{BRb}{BRb} Don Blasius and J. Rogawski. Zeta functions of Shimura varieties. \emph{Motives (Seattle, WA, 1991)}, 525--571. In \emph{Proc. Sympos. Pure Math.}, volume 55.
\refitem{BT13}{BT13} Raoul Bott and Loring W. Tu. \emph{Differential forms in algebraic topology}, volume 82. Springer Science \& Business Media, 2013.
\refitem{Cas73}{Cas73} William Casselman. On some results of Atkin and Lehner. \emph{Mathematische Annalen}, 201(4):301--314, 1973.
\refitem{CGP15}{CGP15} Brian Conrad, Ofer Gabber, and Gopal Prasad. \emph{Pseudo-reductive groups}, volume 26 of New Mathematical Monographs. Cambridge University Press, Cambridge, second edition, 2015.
\refitem{Col11}{Col11} P. Colmez. \emph{Éléments d'analyse et d'algèbre (et de théorie des nombres)}, second ed. Éditions de l'École Polytechnique, Palaiseau, 2011.
\refitem{CS17}{CS17} Ana Caraiani and Peter Scholze. On the generic part of the cohomology of compact unitary Shimura varieties. \emph{Annals of Mathematics}, pages 649--766, 2017.
\refitem{DE14}{DE14} Anton Deitmar and Siegfried Echterhoff. \emph{Principles of harmonic analysis}. Springer, 2014.
\refitem{Del71a}{Del71a} P. Deligne. Travaux de Shimura, Sém. Bourbaki 389. \emph{Lecture Notes in Math}, 244, 1971.
\refitem{Del71b}{Del71b} Pierre Deligne. Formes modulaires et représentations $\ell$-adiques. In \emph{Séminaire Bourbaki vol. 1968/69, Exposés 347--363}, pages 139--172. Springer, 1971.
\refitem{Del73}{Del73} Pierre Deligne. Formes modulaires et représentations de $GL(2)$. In \emph{Modular functions of one variable II}, pages 55--105. Springer, 1973.
\refitem{Del75}{Del75} Pierre Deligne. Courbes elliptiques: formulaire d'après J. Tate. In \emph{Modular functions of one variable, IV} (Proc. Internat. Summer School, Univ. Antwerp, Antwerp, 1972), volume 476, pages 53--73, 1975.
\refitem{Del79}{Del79} Pierre Deligne. Variétés de Shimura: interprétation modulaire, et techniques de construction de modèles canoniques. In \emph{Automorphic forms, representations and $L$-functions} (Proc. Sympos. Pure Math., Oregon State Univ., Corvallis, Ore., 1977), Part 2, pages 247--289, 1979.
\refitem{DR73}{DR73} Pierre Deligne and Michael Rapoport. Les schémas de modules de courbes elliptiques. In \emph{Modular functions of one variable II}, pages 143--316. Springer, 1973.
\refitem{DS74}{DS74} Pierre Deligne and Jean-Pierre Serre. Formes modulaires de poids 1. In \emph{Annales scientifiques de l'École Normale Supérieure}, volume 7, pages 507--530, 1974.
\refitem{DS05}{DS05} Fred Diamond and Jerry Michael Shurman. \emph{A first course in modular forms}, volume 228. Springer, 2005.
\refitem{Fla79}{Fla79} Dan Flath. Decomposition of representations into tensor products. In \emph{Automorphic forms, representations and $L$-functions} (Proc. Sympos. Pure Math., Oregon State Univ., Corvallis, Ore., 1977), Part 1, pages 179--183, 1979.

% ===== PDF p.103 =====
\refitem{FM95}{FM95} Jean-Marc Fontaine and Barry Mazur. Geometric Galois representations. \emph{Elliptic curves, modular forms, \& Fermat's last theorem (Hong Kong, 1993), Ser. Number Theory, I}, pages 41--78, 1995.
\refitem{Fon94}{Fon94} Jean-Marc Fontaine. Le corps des périodes $p$-adiques. \emph{Astérisque}, 223:59--111, 1994.
\refitem{Gel75}{Gel75} Stephen S. Gelbart. \emph{Automorphic forms on adele groups}, number 83. Princeton University Press, 1975.
\refitem{GH19}{GH19} Jayce R. Getz and Heekyoung Hahn. \emph{An introduction to automorphic representations}. 见网页 \url{https://www.math.duke.edu/~hahn/GTM.pdf}, 2019.
\refitem{GS17}{GS17} Philippe Gille and Tamás Szamuely. \emph{Central simple algebras and Galois cohomology}, volume 165. Cambridge University Press, 2017.
\refitem{Har88}{Har88} Michael Harris. Automorphic forms and the cohomology of vector bundles on Shimura varieties. In \emph{Proceedings of the Conference on Automorphic Forms, Shimura Varieties, and $L$-functions}, Ann Arbor, pages 41--91, 1988.
\refitem{Lan17}{Lan17} Kai-Wen Lan. \emph{An example-based introduction to Shimura varieties}. Preprint, 2017.
\refitem{LR87}{LR87} R. P. Langlands and M. Rapoport. Shimuravarietäten und Gerben. \emph{J. Reine Angew. Math.}, 378:113--220, 1987.
\refitem{Lur}{Lur} Jacob Lurie. \emph{Tamagawa numbers via nonabelian Poincaré duality}. 讲义见 \url{http://www.math.harvard.edu/~lurie/282y.html}.
\refitem{Mil05}{Mil05} James S. Milne. Introduction to Shimura varieties. \emph{Harmonic analysis, the trace formula, and Shimura varieties}, 4:265--378, 2005.
\refitem{MW89}{MW89} Colette Mœglin and J.-L. Waldspurger. Le spectre résiduel de $GL(n)$. In \emph{Annales scientifiques de l'École normale supérieure}, volume 22, pages 605--674, 1989.
\refitem{Pan19}{Pan19} Lue Pan. The Fontaine--Mazur conjecture in the residually reducible case. \emph{arXiv preprint} arXiv:1901.07166, 2019.
\refitem{PS16}{PS16} Vincent Pilloni and Benoît Stroh. Arithmétique $p$-adique des formes de Hilbert: surconvergence, ramification et modularité. \emph{Astérisque}, 382:195--266, 2016.
\refitem{Rib77}{Rib77} Kenneth A. Ribet. Galois representations attached to eigenforms with nebentypus. In \emph{Modular functions of one variable V}, pages 18--52. Springer, 1977.
\refitem{Rib90}{Rib90} Kenneth A. Ribet. From the Taniyama--Shimura conjecture to Fermat's last theorem. \emph{Annales de la Faculté des sciences de Toulouse: Mathématiques}, Ser. 5, 11(1):116--139, 1990.
\refitem{Ser97}{Ser97} Jean-Pierre Serre. \emph{Abelian $\ell$-adic representations and elliptic curves}. CRC Press, 1997.
\refitem{Tho}{Tho} Jack Thorne. \emph{Topics in automorphic forms} (notes taken by Chao Li). 课程讲义见 \url{http://www.math.columbia.edu/~chaoli/docs/AutomorphicForm.html}.
\refitem{Thu82}{Thu82} William P. Thurston. Three dimensional manifolds, Kleinian groups and hyperbolic geometry. \emph{Bulletin of the American Mathematical Society}, 6(3):357--381, 1982.
\refitem{Wil95}{Wil95} Andrew Wiles. Modular elliptic curves and Fermat's last theorem. \emph{Annals of Mathematics}, 141(3):443--551, 1995.
\end{references}

\vspace{2em}
\noindent\textbf{作者地址}：Mathematisches Institut, Universität Bonn, Endenicher Allee 60, 53115 Bonn, Deutschland。电子邮件：\texttt{ja@math.uni-bonn.de}。

\printindex
\end{document}
